% !TEX root = ../print.tex
% Appendix A: the discrete stratum on the line. The draft's sections 2.2-2.5 and Proposition 2.24
% (draft/spatial-hemigroup-pyramids.md, lines 162-259 and 317-324), moved here by adr/0004. The
% draft marks these statements as ports of the causal companion note (Paper II), which is
% unreleased and never cited; every statement is proved here in full (plan decision 3). Paper I's
% Remark 13.1 is not cited either: the canonical discrete gauge is lem:family-dilations(3).
%
% The families of this appendix are read through their kernels (lem:line-axioms-kernels, chapter
% 2), so that the lemmas of chapter 3 apply at d = 1 (the draft's Remark 3.16). Three lemmas are new
% at the transcription: lem:line-exponent-data (the data of an exponent), which carries what the
% draft left to "matching the representations"; lem:knot-exponent-ladder, the three facts the
% draft's prose definition of the modulation orbit asserts; and lem:laplace-atoms-admissible, the
% Frullani-type identity and the admissibility of sums of Laplace atoms, used by three nodes.
% Displays are tagged (\ref{eq:line-e1})-(\ref{eq:line-e9}), apart from the section numbers A.1-A.5.
%
% The draft's review flags on this material are answered at the nodes: Theorem 2.9's "modulated
% by" (thm:annulus-parametrization-phase), Remark 2.18(1) (prop:classical-pyramids-approximate and
% the paragraph after it), section 2.4 (section A.3, with the extension of
% notes/phase-entropy-without-covariant-member.md), the "published form" of the limit kernel (the
% paragraph after thm:matern-pyramid-line, where the comparison is cut).
%
% Stays prose, and why, one line each:
% - the paragraph after lem:gauge-group (what the gauge does and does not fix): explanation;
% - the paragraph after thm:annulus-parametrization (Maejima-Sato as the prior form): positioning,
%   and its wording is the author's (HANDOFF, item G's follow-ups);
% - the paragraph after prop:strata-cylinder-coordinates ("two resolutions"): explanation;
% - the paragraph after lem:step-field-determines-family (ladder, modulation orbit): it leads into
%   def:modulation-orbit, which carries the content;
% - the paragraphs before def:phase-entropy and after thm:phase-entropy-modulated-selection:
%   motivation and the open question (marked open);
% - the paragraph after prop:classical-kernels-no-exact-cascade (what the classical constraints
%   capture): positioning in the literature;
% - the paragraph after thm:matern-pyramid-line (the Matern pyramid as canonical): reading;
% - the four readings of non-creation before prop:ladder-non-creation-line: positioning;
% - Remark 2.18(4), multiresolution analysis: positioning; its check (the nonnegative scaling
%   function) is cut, not answered;
% - Remark 2.19 (sampling): explanation; its check is answered by prop:sampled-laplace-exact;
% - the paragraph after prop:extreme-rays-random-walks: reading.

\chapter{The discrete stratum on the line}
\label{ch:line-appendix}

This appendix is the theory of the families covariant under one ratio $q > 1$ on the line, the
discrete stratum of Theorem~\ref{thm:lattice-line}, carried as far as the line allows. The plane
uses it at $d = 1$ only through chapter~3. It proves five things. A family of the stratum is the
same thing as a path of exponents over one octave of scales, glued at its ends
(\S\ref{sec:line-annulus}). The \Levy{} data of the path are one positive measure on the quadrant of
scales and displacements, invariant under a diagonal dilation, and the extreme members of that cone
are the knot-atom cascades, whose circle averages are the Cin rays of the continuous stratum
(\S\ref{sec:line-surface}). On a ladder with no Gaussian part, the phase entropy selects the
fully covariant member when there is one. It selects a modulated covariant member on a strictly
larger class of ladders, and nothing on a ladder whose step measure has a singular part, which
includes every grid pyramid (\S\ref{sec:line-phase-entropy}). The finite masks of the classical
pyramids are steps of no family. The \Matern{} pyramid is a fully covariant family whose octave
steps are single Laplace kernels, and its ladder creates no extrema (\S\ref{sec:line-examples}).
The $q$-Cin rays, the extreme rays of the cone of semi-selfdecomposable exponents, are limits of
cascades of random walks, and for an integer ratio each truncation of the cascade is carried by
one lattice (\S\ref{sec:line-random-walks}).

\section{The setting, the gauge and the annulus parametrization}
\label{sec:line-annulus}

Throughout, $q > 1$ and $\ell := \log q$. For $c > 0$ and a measure $\mu$ on $\RR$ or on
$(0,\infty)$, $T_c\mu$ is its image under $x \mapsto cx$, so that
$\widehat{T_c\mu}(\omega) = \widehat\mu(c\omega)$. On $\RR$ it is the dilation $D_c$ of chapter~3
at $d = 1$. The $q^{\ZZ}$-covariant families, their canonical discrete gauge and the fully
covariant families are those of Definition~\ref{def:q-covariant-line}; fully covariant means (A8)
for every $\lambda > 0$. By Lemma~\ref{lem:line-axioms-kernels} such a family is, read through its
kernels, a symmetric $q$-covariant family in dimension~$1$ (Definition~\ref{def:family-dimension-d}),
with the same bijections, so the lemmas of chapter~3 apply at $d = 1$. Write $\mathrm{ND}_s$ for the
set of functions
\[
  \psi(\omega) = a\,\omega^2 + \int_{(0,\infty)}\bigl(1 - \cos\omega v\bigr)\,\varpi(dv), \qquad
  a \ge 0, \quad \int(1 \wedge v^2)\,\varpi(dv) < \infty ,
\]
and call $(a, \varpi)$ the \emph{data} of $\psi$: a Gaussian coefficient and a folded \Levy{}
measure. The first lemma says what this appendix uses of the \Levy{}--Khintchine formula.

% draft: none. New at the transcription: the uniqueness of the data and two rules for computing
% with them, which the draft uses under "matching the representations" and "unique by [V-line
% v0.1, Prop. 2.3(3)]" (Lemma 2.3's proof, line 140; Theorems 2.9-2.10).
\begin{lemma}[the data of an exponent]
\label{lem:line-exponent-data}
\lean{Pyramids.line_exponent_data_1, Pyramids.line_exponent_data_1_converse,
Pyramids.line_exponent_data_1_unique, Pyramids.line_exponent_data_2, Pyramids.line_exponent_data_3}\leanok
\uses{thm:cited-A39}
\begin{enumerate}
\item Every $\psi \in \mathrm{ND}_s$ is $-\log\widehat\mu$ for exactly one symmetric infinitely
divisible law $\mu$ on $\RR$, and every such $-\log\widehat\mu$ lies in $\mathrm{ND}_s$. The data of
$\psi$ are unique.
\item If $\psi$ has the data $(a, \varpi)$ and $\lambda > 0$, then $\psi(\lambda\,\cdot)$ has the
data $(\lambda^2a, T_\lambda\varpi)$.
\item If $\psi, \psi' \in \mathrm{ND}_s$ have the data $(a, \varpi)$ and $(a', \varpi')$, then
$\psi' - \psi \in \mathrm{ND}_s$ iff $a' \ge a$ and $\varpi' \ge \varpi$, and then $\psi' - \psi$
has the data $(a' - a, \varpi' - \varpi)$.
\end{enumerate}
% CHANGED (fidelity review, 2026-11-15, R26/R50/R69): the status line lists the declarations' hypotheses beside the node's (LINKAGE rule 6).
\statusA\quad\ledger{A39}\quad\emph{(\textbf{Assignment.} \ledger{A39} carries (1): the
generating triplet of an infinitely divisible law exists, is unique, and every triplet occurs, read
for symmetric laws, whose triplet is $(2a, \nu, 0)$ with $\nu$ the symmetric measure whose image
under $x \mapsto |x|$ is $\varpi$. Clause (2) is proved here from (1) by a change of variables,
and clause (3) from (1) by linearity. \textbf{Hypotheses} (the node's, each with what its
declarations require): (1) $\psi \in \mathrm{ND}_s$: same (\texttt{IsSymLevyExponent}); $\mu$ a
symmetric infinitely divisible law: same (\texttt{IsProbabilityMeasure}, \texttt{IsSymmetric},
\texttt{IsInfinitelyDivisible}); $\psi$ with the data $(a, \varpi)$: same (\texttt{HasExpData}, the
measure $\varpi$ a measure on $\RR$ null on $(-\infty, 0]$); (2) $\lambda > 0$: same; (3) $\psi,
\psi'$ with the data $(a, \varpi)$, $(a', \varpi')$: same.)}
\end{lemma}

% CHANGED (wave 2b, 2026-11-13): (3) says why the difference of two Levy measures is a measure, as the
% machine-checked proof constructs it; the closing sentence says which steps are checked and which
% are cited interfaces.
% CHANGED (the theory paper, 2026-11-16): the closing sentence on checking and cited facts shortened; no step changed.
\begin{proof}
\leanok
\uses{thm:cited-A39}
(1) Let $\nu$ be a symmetric \Levy{} measure on $\RR$. The imaginary part of the integrand,
$\sin\omega x - \omega x\mathbf 1_{|x| \le 1}$, is odd and $\nu$-integrable, so its integral
vanishes, and
\[
  \int\bigl(e^{i\omega x} - 1 - i\omega x\mathbf 1_{|x| \le 1}\bigr)\,\nu(dx)
  = -\int(1 - \cos\omega x)\,\nu(dx).
\]
So the law with the triplet $(2a, \nu, 0)$, $\nu$ the symmetrization $\tfrac12(\varpi + \varpi(-\,\cdot))$
of $\varpi$, has the exponent $\psi$. It exists by [\sscite{sato1999levy}, Thm.~8.1(iii),
p.~38]. It is symmetric because its transform $e^{-\psi}$ is real and even. A symmetric law
has a symmetric triplet, by the uniqueness of the triplet [\sscite{sato1999levy}, Thm.~8.1(ii),
p.~38] applied to the law and its
reflection, so its exponent is of this form. The law with exponent $\psi$ is unique because the
transform determines the law. The data are unique because the triplet is.

(2) $\int(1 - \cos\lambda\omega v)\,\varpi(dv) = \int(1 - \cos\omega u)\,(T_\lambda\varpi)(du)$ by
the substitution $u = \lambda v$, and $a(\lambda\omega)^2 = \lambda^2a\,\omega^2$. Then (1) gives
the uniqueness.

(3) Let $a' \ge a$ and $\varpi' \ge \varpi$. The difference $\varpi' - \varpi$ is a measure: $\varpi$
is finite on each of the disjoint sets $\{v \le 0\} \cup \{v > 1\}$ and
$\{1/(n+1) < v \le 1/n\}$, $n \ge 1$, which cover $\RR$, and $\varpi' - \varpi$ is the sum of the
differences of the restrictions to these sets. It is at most $\varpi'$, so it is a \Levy{} measure
on $(0,\infty)$, and $\psi' - \psi$ has the data $(a' - a, \varpi' - \varpi)$ by linearity of the
formula. Conversely, if $\psi' - \psi$ has data $(a'', \varpi'')$, then $\psi' = \psi + (\psi' - \psi)$
has the data $(a + a'', \varpi + \varpi'')$, which are $(a', \varpi')$ by (1).

Every step is machine-checked, and the cited fact is \ledger{A39}, read on $\RR$ inside $\RR^1$
through $x \mapsto xe_1$.
\end{proof}

For a
$q^{\ZZ}$-covariant family write $\mu_t := \mu_{0,t}$, $\Psi_t := -\log\widehat\mu_t$ and
$(a(t), \varpi_t)$ for the data of $\Psi_t$. By Lemma~\ref{lem:family-structure} these exist,
$a(t)$ and $\varpi_t$ are nondecreasing in $t$, and $\mu_{s,t}$ has the data
$(a(t) - a(s), \varpi_t - \varpi_s)$. By Lemma~\ref{lem:family-dilations}(3) the scale may be
relabelled so that $S_qt = qt$, the \emph{canonical discrete gauge}. Then
Lemma~\ref{lem:discrete-covariance-levy} gives
\[
  \Psi_{qt}(\omega) = \Psi_t(q\omega), \qquad a(qt) = q^2a(t), \qquad \varpi_{qt} = T_q\varpi_t
  \qquad (t > 0),
\]
which is Lemma~\ref{lem:discrete-covariance-line}. The scales $q^n$ are the knots, the laws
$\mu_{q^n}$ the knot kernels and the octave steps $\mu_{q^n,q^{n+1}}$ the ladder
(Definition~\ref{def:knots-and-ladder}). The first
question is how much the canonical discrete gauge fixes.

% draft: Lemma 2.8 (section 2.2), lines 164-166 of draft/spatial-hemigroup-pyramids.md.
% CHANGED: clause (1) restated as lem:family-dilations(3) states it (relabellings chi of [0,inf)),
% and the "infinite-dimensional" sentence made precise as the bijection of clause (3).
% CHANGED: the clause on the continuous stratum is stated and proved for two gauges that both
% conjugate every S_lambda to a dilation; its reading as "the case beta = 1 of the homogeneous
% gauges of [V-line v0.1, Remark 6.7]" moves to the prose after the lemma.
\begin{lemma}[the gauge group]
\label{lem:gauge-group}
\lean{Pyramids.gauge_group_1, Pyramids.gauge_group_2, Pyramids.gauge_group_3,
Pyramids.gauge_group_4}
\leanok
\uses{lem:family-dilations}
Let $\mathcal{G}_q$ be the group of increasing homeomorphisms $g$ of $(0,\infty)$ with
$g(qt) = q\,g(t)$ for all $t > 0$.
\begin{enumerate}
% CHANGED (2026-11-14, a strengthening): (1) and (4) are stated for increasing bijections S_q and
% S_lambda of [0,infinity), such as those of a family; the checked proof reads nothing else of the
% family (records/pyramids/WAVE2-report-2e.md).
\item Let $S_q$ be an increasing bijection of $[0,\infty)$, such as the bijection of a
$q^{\ZZ}$-covariant family. Two increasing homeomorphisms
$\chi, \chi'$ of $[0,\infty)$ with $\chi S_q = q\chi$ and $\chi'S_q = q\chi'$ satisfy
$\chi' = g\chi$ with $g \in \mathcal{G}_q$, and $g\chi$ is such a homeomorphism for every
$g \in \mathcal{G}_q$. So the canonical discrete gauge is determined up to $\mathcal{G}_q$.
\item In the coordinate $u = \log t$, the map $g \mapsto \tilde g$,
$\tilde g(u) := \log g(e^u)$, is an isomorphism of $\mathcal{G}_q$ onto the group of increasing
homeomorphisms $\tilde g$ of $\RR$ with $\tilde g(u + \ell) = \tilde g(u) + \ell$, the lifts of the
orientation-preserving homeomorphisms of the circle $\RR/\ell\ZZ$.
\item Restriction to $[1,q]$ is a bijection of $\mathcal{G}_q$ onto the set of continuous, strictly
increasing maps $\phi : [1,q] \to (0,\infty)$ with $\phi(q) = q\,\phi(1)$.
\item Let $S_\lambda$, $\lambda > 0$, be increasing bijections of $[0,\infty)$, such as the
bijections of a family that is $\lambda$-covariant for every $\lambda > 0$. If $\chi$ and $\chi'$ are
increasing homeomorphisms of $[0,\infty)$ with $\chi S_\lambda = \lambda\chi$ and
$\chi'S_\lambda = \lambda\chi'$ for every $\lambda > 0$, then $\chi' = c\chi$ for a constant
$c > 0$.
\end{enumerate}
% CHANGED (fidelity review, 2026-11-15, R64): the status line says which declarations keep a phantom
% family binder and which clause is stated in decomposed form.
% CHANGED (fidelity review, 2026-11-15, R26/R50/R69): the status line lists the declarations' hypotheses beside the node's (LINKAGE rule 6).
\statusT\quad\emph{(\textbf{Hypotheses} (the node's, each with what its declarations require):
$q > 1$: same for \texttt{gauge\_group\_1}, \texttt{\_2}, \texttt{\_3}, absent
from \texttt{gauge\_group\_4}, which has no $q$; (1) $S_q$ an increasing bijection of
$[0,\infty)$: same, read as \texttt{IsCovariantAt}~$\mu$~$q$~$S$ for a family $\mu$ that is a free
binder of \texttt{gauge\_group\_1}, which the laws $\mu \equiv 0$ satisfy for every increasing
bijection $S$, so the binder is phantom; $\chi$ an increasing homeomorphism: weaker,
\texttt{IsGauge}~$\chi$ (an increasing bijection of $[0,\infty)$, continuity not required, though it
follows); $\chi'$ an increasing homeomorphism: same (\texttt{IsGauge} and \texttt{ContinuousOn});
$\chi S_q = q\chi$, $\chi'S_q = q\chi'$: same, on $[0,\infty)$; (2), (3): $g \in \mathcal{G}_q$:
same (\texttt{IsInGaugeGroup}); (4) the $S_\lambda$ increasing bijections: same, through
\texttt{IsCovariantAt} with the same phantom binder $\mu$ in \texttt{gauge\_group\_4}; $\chi, \chi'$
increasing homeomorphisms: weaker, \texttt{IsGauge} only; the two equations for every
$\lambda > 0$: same. The isomorphism of (2) is stated in decomposed form, the reader composing it:
\texttt{gauge\_group\_2} gives that $g \mapsto \tilde g$ lands in the lifts, is onto, is injective up
to equality on $(0,\infty)$ and carries composition to composition.)}
\end{lemma}

% CHANGED (wave 2d, 2026-11-13): in (1) the converse reads continuity off monotonicity, as the
% checked proof does, and a closing paragraph says which steps are checked; otherwise the route is
% the printed one.
% CHANGED (the theory paper, 2026-11-16): the closing sentence on checking and cited facts shortened; no step changed.
\begin{proof}
\leanok
\uses{lem:family-dilations}
(1) is Lemma~\ref{lem:family-dilations}(3). In detail, $g := \chi'\chi^{-1}$ is an increasing
homeomorphism of $[0,\infty)$ fixing $0$, and $g(qt) = \chi'S_q\chi^{-1}(t) = q\chi'\chi^{-1}(t) =
q\,g(t)$. So its restriction to $(0,\infty)$ lies in $\mathcal{G}_q$. Conversely, for
$g \in \mathcal{G}_q$ extended by $g(0) := 0$, the map $g\chi$ is strictly increasing on
$[0,\infty)$ and maps it onto itself, so it has no jump and is an increasing homeomorphism of
$[0,\infty)$; and $(g\chi)S_q = g(q\chi) = q\,g\chi$.

(2) $\tilde g$ is an increasing homeomorphism of $\RR$, and $g(qt) = q\,g(t)$ reads
$\tilde g(u + \ell) = \tilde g(u) + \ell$ at $t = e^u$. Composition is preserved, and every such
$\tilde g$ comes from $g(t) := \exp\tilde g(\log t)$.

(3) For $g \in \mathcal{G}_q$ the restriction $\phi := g|_{[1,q]}$ has the stated properties, and
$g(q^ns) = q^n\phi(s)$ for $s \in [1,q)$ and $n \in \ZZ$, so the restriction determines $g$.
Conversely, given $\phi$, put $g(q^ns) := q^n\phi(s)$ for $s \in [1,q)$, $n \in \ZZ$. The intervals
$q^n[\phi(1), q\phi(1))$, $n \in \ZZ$, partition $(0,\infty)$ in increasing order, and $g$ maps
$q^n[1,q)$ increasingly onto the $n$-th of them, so $g$ is an increasing bijection of $(0,\infty)$.
It is continuous at the points $q^n$ because $\phi(q) = q\phi(1)$, and $g(qt) = q\,g(t)$ by
construction.

(4) $g := \chi'\chi^{-1}$ has $g(\lambda t) = \chi'S_\lambda\chi^{-1}(t) = \lambda\,g(t)$ for all
$\lambda, t > 0$. With $t = 1$, $g(\lambda) = \lambda\,g(1)$. So $\chi' = c\chi$ with $c = g(1)$.

% CHANGED (2026-11-14, a strengthening): the closing paragraph follows the restated (1) and (4).
Every step is machine-checked, and the proof cites no fact of its own.
\end{proof}

The canonical discrete gauge fixes the knots only up to $\mathcal{G}_q$, and which scale is called
$1$ is a choice in that group. A relabelling in $\mathcal{G}_q$ that fixes the knots changes the
laws between the knots and leaves the knot kernels and the ladder alone. The families with a given
ladder are more than these relabellings (Proposition~\ref{prop:ladder-knot-kernel}(2)). In the full
stratum the freedom shrinks to the scale unit, clause~(4), and the gauge freedom of the full
stratum is discussed in \prog{V-line v0.1}{Remark~6.7, pp.~32--33}. Statements below are made
in one canonical discrete gauge, fixed once.

% draft: Theorem 2.9 (section 2.2), lines 170-174 of draft/spatial-hemigroup-pyramids.md; the
% bijection only, the log-periodic form being the next node.
% CHANGED: split from Theorem 2.9: the bijection here, the form (2.2) in
% thm:annulus-parametrization-phase (separate proof, separate grade).
% CHANGED: condition (2) narrowed to differences that are not identically zero. Without it the
% bijection is false: a path constant on a subinterval of [1,q] satisfies the draft's (1)-(4) and
% gives laws mu_{s,t} = delta_0 with s < t, against (ND). This answers the draft's check on (ND).
\begin{theorem}[annulus parametrization]
\label{thm:annulus-parametrization}
\lean{Pyramids.annulus_parametrization_map, Pyramids.annulus_parametrization_recover,
Pyramids.annulus_parametrization_injective, Pyramids.annulus_parametrization_surjective}
\leanok
\uses{lem:line-axioms-kernels, lem:line-exponent-data}
In the canonical discrete gauge, the $q^{\ZZ}$-covariant families are in bijection with the paths
$(\Psi_t)_{t \in [1,q]}$ of functions on $\RR$ such that
\begin{enumerate}
\item $\Psi_t \in \mathrm{ND}_s$ and $\Psi_t \not\equiv 0$ for each $t \in [1,q]$;
\item $\Psi_{t'} - \Psi_t \in \mathrm{ND}_s$ and $\Psi_{t'} - \Psi_t \not\equiv 0$ for
$1 \le t < t' \le q$;
\item $\Psi_q = \Psi_1(q\,\cdot)$ (gluing);
\item $t \mapsto \Psi_t(\omega)$ is continuous on $[1,q]$ for each $\omega$.
\end{enumerate}
The family is recovered by $\Psi_{q^nt} := \Psi_t(q^n\,\cdot)$ for $t \in [1,q)$ and $n \in \ZZ$,
$\Psi_0 := 0$, and $\widehat\mu_{s,t} = e^{-(\Psi_t - \Psi_s)}$.
% CHANGED (fidelity review, 2026-11-15, R26/R50/R69): the status line lists the declarations' hypotheses beside the node's (LINKAGE rule 6).
\statusT\quad\emph{(\textbf{Hypotheses} (the node's, each with what its declarations require): $q > 1$:
same for \texttt{annulus\_parametrization\_recover}, \texttt{\_injective} and \texttt{\_surjective},
absent from \texttt{annulus\_parametrization\_map}; a $q^{\ZZ}$-covariant family in the canonical
discrete gauge: same, \texttt{IsQFamily}, the symmetric family of laws in dimension~$1$ with
$S_qt = qt$ that Lemma~\ref{lem:line-axioms-kernels} reads off the operators; the path conditions
(1)--(4): same, \texttt{IsAnnulusPath}, with the strict pairs and $\not\equiv 0$ of (2).)}
\end{theorem}

% CHANGED (wave 2e, 2026-11-14): the proof follows the machine-checked route. The gluing and the
% recovery read covariance directly, in place of Lemma discrete-covariance-levy; surjectivity says
% how the folded data become data in R^1, uses the uniqueness of folded data for the gluing, and
% reads the laws back on the line; a closing paragraph says which steps are checked.
% CHANGED (the theory paper, 2026-11-16): the closing sentence on checking and cited facts shortened; no step changed.
\begin{proof}
\leanok
\uses{lem:line-axioms-kernels, lem:line-exponent-data, lem:family-structure, lem:synthesis}
By Lemma~\ref{lem:line-axioms-kernels} it suffices to treat symmetric $q$-covariant families
$(\mu_{s,t})$ in dimension~$1$ with $S_qt = qt$.

\emph{The map.} Send such a family to $(\Psi_t)_{t \in [1,q]}$. By
Lemma~\ref{lem:family-structure} at $d = 1$, read through the embedding $x \mapsto xe_1$ of $\RR$
in $\RR^1$, every $\mu_{s,t}$ is infinitely divisible, and $\widehat\mu_t > 0$. So $\Psi_t$ is the
exponent of the symmetric infinitely divisible law $\mu_t$, hence $\Psi_t \in \mathrm{ND}_s$
(Lemma~\ref{lem:line-exponent-data}(1)), and $\Psi_t \not\equiv 0$ because $\mu_t \ne \delta_0$
by (ND). By (I1), $\widehat\mu_{t,t'} = e^{-(\Psi_{t'} - \Psi_t)}$, so likewise
$\Psi_{t'} - \Psi_t \in \mathrm{ND}_s$, not identically zero by (ND): this is (1) and (2).
Covariance at the pair $(0, 1)$, $T_q\mu_1 = \mu_q$, gives (3), and (I2) with
$\widehat\mu_t > 0$ gives (4). The map is injective: covariance at the pairs $(0, t)$ gives
$\Psi_{qt} = \Psi_t(q\,\cdot)$, and by iteration $\Psi_{q^nt} = \Psi_t(q^n\,\cdot)$ for
$n \in \ZZ$, so the path, extended as in the statement, is $\Psi_t$ at every $t \ge 0$; and
$\widehat\mu_{s,t} = e^{-(\Psi_t - \Psi_s)}$ by (I1). This is the recovery, and two families
with the same path have the same transforms, hence the same laws.

\emph{Surjectivity.} Let a path satisfy (1)--(4), and let $(a(t), \varpi_t)$ be the folded data of
$\Psi_t$. In $\RR^1$ they are the \Levy{} data $(a(t), \nu_t)$, with
$\nu_t := \tfrac12(e_*\varpi_t + (-e)_*\varpi_t)$ the symmetrization of $\varpi_t$, $e = e_1$,
whose exponent at $z$ is $\Psi_t(z_1)$. By (2) and Lemma~\ref{lem:line-exponent-data}(3), for
$t < t'$ the difference $(a(t') - a(t), \varpi_{t'} - \varpi_t)$ is nonnegative, so
$a(t') - a(t) \ge 0$ and $\nu_t \le \nu_{t'}$; it is not zero, since
$\Psi_{t'} - \Psi_t \not\equiv 0$. By (3) and Lemma~\ref{lem:line-exponent-data}(2), the
function $\Psi_q = \Psi_1(q\,\cdot)$ has the folded data $(q^2a(1), T_q\varpi_1)$, so
$(a(q), \varpi_q) = (q^2a(1), T_q\varpi_1)$ by the uniqueness of folded data
(Lemma~\ref{lem:line-exponent-data}(1)), and $\nu_q = D_q\nu_1$, the symmetrization commuting
with the dilation. With (4), these are the hypotheses of Lemma~\ref{lem:synthesis} at $d = 1$.
It gives laws on $\RR^1$ with the transforms $e^{-(\Psi_t - \Psi_s)}$ of the extended path, read
at $z_1$, and their images under $y \mapsto y_1$ are laws on $\RR$ with the transforms
$e^{-(\Psi_t - \Psi_s)}$. Conversely, for any laws on $\RR$ with these transforms, their images
in $\RR^1$ have the transforms of Lemma~\ref{lem:synthesis}, so they form a symmetric $q$-covariant
family in the canonical gauge, and so do the laws themselves, read back on $\RR$. Its exponent at
$t \in [1, q)$ is $\Psi_t$, and at $t = q$ it is $\Psi_1(q\,\cdot) = \Psi_q$ by (3): its image under
the map is the path.

Every step is machine-checked, and the cited facts enter only through
Lemmas~\ref{lem:family-structure}, \ref{lem:line-exponent-data} and \ref{lem:synthesis}.
\end{proof}

The bijection is, for laws, a known construction: Maejima and Sato build the semi-selfsimilar
additive processes from paths of laws over one period with these conditions
[\sscite{maejima1999semi}, Thm.~7, p.~361]. What this appendix adds is the reading in the
scale-space axioms and what follows from it.
% The anchor, as the supplement prints it after the citation audit (pass 7), verified in
% records/pyramids/CITATIONS-acquired.md, item 7.

% draft: Theorem 2.9 (section 2.2), the display (2.2) and the sentence after it, lines 170-172 of
% draft/spatial-hemigroup-pyramids.md.
% CHANGED: split from Theorem 2.9 (see thm:annulus-parametrization).
% CHANGED: renamed the periodic function chi to Xi: chi is the gauge of Theorem 2.0.
% CHANGED: "every pyramid is a similarity family modulated by a log-periodic function of the
% scale" removed (review flag, "Thm. 2.9's 'modulated by' overstates (2.2)"): (2.2) is a change of
% variables, not a factorization. The statement now says what (2.2) is, and when the phase drops
% out, with the gauge made explicit (a fully covariant family in the canonical discrete gauge need
% not have the form F(t omega) until it is relabelled).
% CHANGED (supplement review, 2026-11-16): (2) says what "F admissible" means for a function on
% [0,infinity), the even extension's being an admissible exponent, which the proof used to say; the
% first variable of Xi renamed z (w is the mask). No claim changed.
\begin{theorem}[the octave phase]
\label{thm:annulus-parametrization-phase}
\lean{Pyramids.annulus_phase_1, Pyramids.annulus_phase_1_converse, Pyramids.annulus_phase_2,
Pyramids.annulus_phase_3}
\leanok
\uses{thm:annulus-parametrization, thm:cited-A4, lem:gauge-group}
Let $(\Psi_t)_{t > 0}$ be the exponents of a $q^{\ZZ}$-covariant family in the canonical discrete
gauge, and put $\psi(u,v) := \Psi_{e^u}(e^v)$ for $u, v \in \RR$.
\begin{enumerate}
\item $\psi(u + \ell, v) = \psi(u, v + \ell)$. The function $\Xi(z;\xi) := \psi(\xi, z - \xi)$ is
$\ell$-periodic in $\xi$, and
\[
  \Psi_t(\omega) = \Xi\bigl(\log(t|\omega|);\ \log t\bigr) \qquad (t > 0,\ \omega \ne 0).
\]
Conversely, for every function $\Xi$ on $\RR^2$, $\ell$-periodic in its second argument,
$\psi(u,v) := \Xi(u + v; u)$ satisfies $\psi(u + \ell, v) = \psi(u, v + \ell)$.
\item $\Xi$ does not depend on its second argument iff $\Psi_t(\omega) = F(t|\omega|)$ for a
function $F$ on $[0,\infty)$. Then $F$ is admissible, that is, its even extension $F(|\cdot|)$ is an
admissible exponent (\ref{eq:admissible-exponent}). Moreover $F \not\equiv 0$, and the family
satisfies (A8) with $S_\lambda t = \lambda t$ for every $\lambda > 0$.
\item Every $q^{\ZZ}$-covariant family in the canonical discrete gauge that satisfies (A8) for
every $\lambda > 0$ takes the form of (2) after the relabelling $t \mapsto g(t)$ by some
$g \in \mathcal{G}_q$.
\end{enumerate}
% CHANGED (2026-11-14, a strengthening): an unused hypothesis dropped from the Lean of the converse
% in (1); the print keeps the setting of the node.
% CHANGED (wave 3d, 2026-11-14): every clause is machine-checked, so statement and proof are marked
% proved; the status line says what is checked.
% CHANGED (fidelity review, 2026-11-15, R26/R50/R69): the status line lists the declarations' hypotheses beside the node's (LINKAGE rule 6).
\statusA\quad\ledger{A4}\quad\emph{(\textbf{Assignment.} \ledger{A4}, the classification of the
fully covariant families on the line (Theorem~\ref{thm:cited-A4}), carries the admissibility of $F$ in (2) and the
form $F(\chi(t)\,\cdot)$ used in (3). Clause (1), the covariance in (2), the dilation from
$F'(\chi(1)\,\cdot)$ to $F$ in (2), which the unnormalized pair requires, and the step from
$\chi(qt) = q\chi(t)$ to $\chi \in \mathcal{G}_q$ in (3) are proved here. Every clause is
machine-checked, Theorem~\ref{thm:cited-A4} being itself a theorem from \ledger{A39}; clause (1)
uses no cited fact. \textbf{Hypotheses} (the node's, each with what its declarations require):
$q > 1$: same, absent from \texttt{annulus\_phase\_1\_converse}; the exponents of a
$q^{\ZZ}$-covariant family in the canonical discrete gauge: same (\texttt{IsQFamily}); (1) the
converse, $\Xi$ $\ell$-periodic in its second argument: same, and no family is required; (2) none
beyond the family; (3) (A8) for every $\lambda > 0$: same, some bijection $S$ with
\texttt{IsCovariantAt}~$\mu$~$\lambda$~$S$ for each $\lambda > 0$.)}
\end{theorem}

% CHANGED (wave 3d, 2026-11-14): the proof follows the machine-checked route. Two steps the print
% left implicit are said: Theorem cited-A4 is applied to the operator family of the kernels, which
% Lemma line-axioms-kernels (2)-(3) provides; and chi in G_q is Lemma gauge-group(1) at the identity
% gauge. The rest is the printed argument.
% CHANGED (fourth blind pass, 2026-11-15): the identity Psi_{qt} = Psi_t(q .) is cited from
% lem:discrete-covariance-line (now in the uses); (2) says that "F admissible" means that F(|.|) is
% of the form (2.0), and uses the evenness of F' for F not identically 0; (3) names the relabelling
% g := chi and says that the relabelled kernels are again a family in the canonical discrete gauge.
\begin{proof}
\leanok
\uses{thm:annulus-parametrization, thm:cited-A4, lem:line-axioms-kernels, lem:family-structure, lem:gauge-group,
lem:discrete-covariance-line}
(1) By Lemma~\ref{lem:discrete-covariance-line}, the family being in the canonical discrete gauge,
$\Psi_{qt}(\omega) = \Psi_t(q\omega)$. This reads $\psi(u + \ell, v) = \psi(u, v + \ell)$ at
$t = e^u$, $\omega = e^v$, and the exponents are even in $\omega$. Then
$\Xi(z;\xi + \ell) = \psi(\xi + \ell, z - \xi - \ell) = \psi(\xi, z - \xi) = \Xi(z;\xi)$, and
$\Xi(\log t + \log|\omega|;\log t) = \psi(\log t, \log|\omega|) = \Psi_t(\omega)$. Conversely
$\psi(u + \ell, v) = \Xi(u + v + \ell; u + \ell) = \Xi(u + v + \ell; u) = \psi(u, v + \ell)$.

(2) If $\Xi(z;\xi) = \Xi_0(z)$, put $F(r) := \Xi_0(\log r)$ for $r > 0$ and $F(0) := 0$. Then
$\Psi_t(\omega) = F(t|\omega|)$. Conversely $\Psi_t(\omega) = F(t|\omega|)$ gives
$\Xi(z;\xi) = F(e^z)$. In that case
\[
  \widehat\mu_{\lambda s,\lambda t}(\omega) = e^{-(F(\lambda t|\omega|) - F(\lambda s|\omega|))}
  = \widehat\mu_{s,t}(\lambda\omega),
\]
so $T_\lambda\mu_{s,t} = \mu_{\lambda s,\lambda t}$, and (A8) holds with $S_\lambda t = \lambda t$
by Lemma~\ref{lem:line-axioms-kernels}(3), for the family with (A1)--(A7) and (ND) whose kernels
are the $\mu_{s,t}$ (Lemma~\ref{lem:line-axioms-kernels}(2)). By Theorem~\ref{thm:cited-A4}
(\prog{V-line v0.1}{Thm.~7.3, p.~37}) there are a gauge $\chi$ and an admissible $F' \not\equiv 0$
with $\Psi_t = F'(\chi(t)\,|\cdot|)$; the kernels of the theorem are the $\mu_{s,t}$, since the
kernel of a convolution operator is determined by it. Put $c := \chi(1)$, which is positive because $\chi$ is an
increasing bijection of $[0,\infty)$. At $t = 1$, $F(|\omega|) = \Psi_1(\omega) = F'(c|\omega|)$.
The function $F'(c\,\cdot)$ is admissible: if $F'$ has the data $(a, k)$, the substitution
$y = cx$ gives $F'(c\omega) = c^2a\,\omega^2 + \int_0^\infty(1 - \cos\omega y)\,k(y/c)\,dy/y$, with
$k(\cdot/c)$ nonincreasing and
$\int(1 \wedge y^2)\,k(y/c)\,dy/y = \int(1 \wedge c^2x^2)\,k(x)\,dx/x \le
\max(1, c^2)\int(1 \wedge x^2)\,k(x)\,dx/x$. So $F(|\cdot|) = F'(c\,|\cdot|) = F'(c\,\cdot)$, the
last because $F'$ is even, is of the form (\ref{eq:admissible-exponent}): $F$ is admissible. And $F \not\equiv 0$:
$F'$ is even and not identically $0$, so $F'(r) \ne 0$ for some $r \ge 0$, and then
$F(r/c) = F'(r) \ne 0$.

(3) Theorem~\ref{thm:cited-A4}, applied as in (2) to the family of the kernels, gives
$\Psi_t = F(\chi(t)\,|\cdot|)$ with $F$ admissible, $F \not\equiv 0$ and
$\chi$ an increasing bijection of $[0,\infty)$, hence a homeomorphism. Covariance at $q$ with
$S_qt = qt$ gives $F(\chi(qt)\,\cdot) = \Psi_{qt} = \Psi_t(q\,\cdot) = F(q\chi(t)\,\cdot)$. The map
$r \mapsto F(r\,\cdot)$ is injective on $(0,\infty)$. Indeed, if $F(r\,\cdot) = F(r'\,\cdot)$ with
$r < r'$, then $F(\omega) = F((r/r')^n\omega) \to F(0) = 0$ for every $\omega$, against
$F \not\equiv 0$. The limit holds because an exponent of the form (\ref{eq:admissible-exponent}) is
continuous with $F(0) = 0$, by dominated convergence. So
$\chi(qt) = q\chi(t)$, and $g := \chi$ restricted to $(0,\infty)$ lies in $\mathcal{G}_q$: this is
Lemma~\ref{lem:gauge-group}(1) for the identity gauge, which the canonical discrete gauge
$S_qt = qt$ satisfies, and the gauge $\chi$. The relabelling $t \mapsto g(t)$ gives the scale $t$
the new label $g(t)$, so the relabelled kernels are $\mu'_{s,t} := \mu_{\chi^{-1}(s),\chi^{-1}(t)}$,
with exponents $\Psi'_t := \Psi_{\chi^{-1}(t)} = F(t\,|\cdot|)$. They are again a symmetric family
in dimension $1$: $\chi^{-1}$ is an increasing homeomorphism of $[0,\infty)$ fixing $0$, so it
keeps (I1), (I2) and the inequality $s < t$ of (ND), and the laws are unchanged. And
$\Psi'_{qt} = F(qt\,|\cdot|) = \Psi'_t(q\,\cdot)$, so by Lemma~\ref{lem:discrete-covariance-line}
the relabelled family is in the canonical discrete gauge, $S'_qt = qt$, and has the form of (2).
\end{proof}

So $\Psi_t(\omega)$ is a function of two coordinates: the log scale-frequency $\log(t|\omega|)$, on
which the similarity families depend alone, and the octave phase $\log t \operatorname{mod} \ell$. A pyramid
is a similarity family exactly when the phase drops out, by (2), and a fully covariant pyramid
becomes one after a relabelling in $\mathcal{G}_q$, by (3).

\section{The surface measure and the knot-atom cascades}
\label{sec:line-surface}

The conditions (2) of Theorem~\ref{thm:annulus-parametrization} form a continuum, one for each pair
of scales. Disintegrated along the quadrant of scales and displacements, the \Levy{} data turn them
into the positivity of one measure, and the extreme members of the cone of such measures are single
orbits of a dilation. The first lemma bounds the \Levy{} integral of the data of an exponent by the
values of the exponent.

% draft: none. New at the supplement review (2026-11-16): the averaging bound, proved inside the
% proofs of thm:surface-measure and prop:strata-cylinder-coordinates-q-cin, stated once and cited
% there (records/pyramids/REVIEW-draft-theory-supplement-dispositions.md, the PROSE flag on the
% repeated arguments).
\begin{lemma}[the averaging bound]
\label{lem:averaging-bound}
\lean{Pyramids.min_one_sq_le_avg, Pyramids.lintegral_min_one_sq_le,
Pyramids.SymLevyPair_a_le_exponent_one, Pyramids.levy_integrable_of_cos_le}
\leanok
\uses{lem:line-exponent-data}
\begin{enumerate}
\item For real $v$, $1 \wedge v^2 \le 3(1 - \cos v) + \int_0^2(1 - \cos\omega v)\,d\omega$.
\item If $\psi \in \mathrm{ND}_s$ has the data $(\alpha, \rho)$, then $\alpha \le \psi(1)$ and
\[
  \int(1 \wedge v^2)\,\rho(dv) \ \le\ 3\psi(1) + \int_0^2\psi(\omega)\,d\omega .
\]
\item If $\psi \in \mathrm{ND}_s$, and a measure $\rho$ on $\RR$ that is a countable sum of finite
measures has $\int(1 - \cos\omega v)\,\rho(dv) \le \psi(\omega)$ for every real $\omega$, then
$\int(1 \wedge v^2)\,\rho(dv) < \infty$.
\end{enumerate}
\statusT\quad\emph{(Cites nothing. \textbf{Hypotheses} (the node's, each with what its declarations
require): (1) $v$ real: same; (2) $\psi \in \mathrm{ND}_s$ with the data $(\alpha, \rho)$: same, $\psi$
read as the exponent of a pair (\texttt{SymLevyPair}), in $[0,\infty]$ in
\texttt{lintegral\_min\_one\_sq\_le}; (3) $\rho$ a countable sum of finite measures: same
(\texttt{SFinite}); $\psi \in \mathrm{ND}_s$: same, the exponent of a pair; the bound for every
$\omega$: same, in $[0,\infty]$.)}
\end{lemma}

\begin{proof}
\leanok
(1) For $|v| \le 1$, $v^2 \le 3(1 - \cos v)$, by the Taylor bound of $\cos$, and the integral is
nonnegative. For $|v| \ge 1$ the integral is $2 - \sin(2v)/v \ge 1 \ge 1 \wedge v^2$.

(2) $\alpha = \alpha\cdot1^2 \le \psi(1)$. Integrate (1) against $\rho$. By Tonelli's theorem,
\[
  \int(1 \wedge v^2)\,\rho(dv) \le 3\int(1 - \cos v)\,\rho(dv)
  + \int_0^2\int(1 - \cos\omega v)\,\rho(dv)\,d\omega ,
\]
and $\int(1 - \cos\omega v)\,\rho(dv) \le \psi(\omega)$ for every $\omega$, the Gaussian term being
nonnegative.

(3) The same computation applies, Tonelli's theorem holding for a countable sum of finite measures,
and gives $\int(1 \wedge v^2)\,\rho(dv) \le 3\psi(1) + \int_0^2\psi(\omega)\,d\omega$. The right side
is finite. If $\psi$ has the data $(a, \varpi)$, then for $0 < \omega \le 2$ one has
$1 - \cos\omega v \le 2(1 \wedge v^2)$, so $\psi(\omega) \le 4a + 2\int(1 \wedge v^2)\,\varpi(dv)$.

Every step is machine-checked, with no cited fact.
\end{proof}

% draft: Theorem 2.10 (section 2.3), lines 180-185 of draft/spatial-hemigroup-pyramids.md.
% CHANGED: added conditions (iii) (no mass on a vertical line, no atom of B) and (iv)
% (nondegeneracy). Without them the stated correspondence is not onto the families: a Lambda with
% mass on a line {y} x (0,inf) gives a path that jumps at y, against (I2), and the knot-atom
% cascades of the next theorem are of that kind.
% CHANGED: the diagonal dilation renamed J (D is the dilation of measures and KL the relative
% entropy; review flag on the three uses of D).
% CHANGED (proof): the passage from the continuity of the exponents to the continuity of t -> varpi_t
% is proved here (the averaging bound), in place of [V-line v0.1, Lemma 4.2] "with Prop. 2.8",
% which states the continuity of the exponents and Levy's continuity theorem and not this passage
% (librarian, records/pyramids/VERBATIM-V-line.md). Ledger A11 is not used.
% CHANGED: the conditions labelled (i)-(iv), as every node cites them (blind review 2, flag 17), and
% the displays tagged (\ref{eq:line-e1})-(\ref{eq:line-e2}), apart from the section numbers A.1-A.5 (flag 18).
\begin{theorem}[the surface measure]
\label{thm:surface-measure}
\lean{Pyramids.surface_measure_family, Pyramids.surface_measure_unique,
Pyramids.surface_measure_pair, Pyramids.surface_measure_injective,
Pyramids.surface_measure_exponent, Pyramids.surface_measure_integrable_iff}
\leanok
\uses{thm:annulus-parametrization, lem:line-exponent-data}
Let $J(y,v) := (qy, qv)$ on $(0,\infty)^2_{(y,v)}$. In the canonical discrete gauge, the
$q^{\ZZ}$-covariant families are in bijection with the pairs $(\Lambda, \mathcal{B})$ of a positive
measure $\Lambda$ on $(0,\infty)^2$ and a positive measure $\mathcal{B}$ on $(0,\infty)_y$ such that
\begin{enumerate}
\item[(i)] $J_*\Lambda = \Lambda$, and $\int(1 \wedge v^2)\,\Lambda((0,t] \times dv) < \infty$ for
some, equivalently every, $t > 0$;
\item[(ii)] $\mathcal{B}((0,t]) < \infty$ for $t > 0$, and $\mathcal{B}(qA) = q^2\mathcal{B}(A)$ for
Borel $A \subseteq (0,\infty)$;
\item[(iii)] $\Lambda(\{y\} \times (0,\infty)) = 0$ and $\mathcal{B}(\{y\}) = 0$ for every $y > 0$;
\item[(iv)] $\Lambda((s,t] \times (0,\infty)) + \mathcal{B}((s,t]) > 0$ for $0 < s < t$;
\end{enumerate}
the family being given by
\begin{equation}\label{eq:line-e1}
  \varpi_t = \Lambda\bigl((0,t] \times \cdot\,\bigr), \qquad a(t) = \mathcal{B}\bigl((0,t]\bigr),
\end{equation}
equivalently
\begin{equation}\label{eq:line-e2}
  \Psi_t(\omega) = a(t)\,\omega^2 + \int_{(0,t] \times (0,\infty)}\bigl(1 - \cos\omega v\bigr)\,
  \Lambda(dy, dv).
\end{equation}
% CHANGED (2026-11-14, a strengthening): unused hypotheses dropped from the Lean; the print keeps
% the setting of the node, which its other clauses read.
% CHANGED (fidelity review, 2026-11-15, R26/R50/R69): the status line lists the declarations' hypotheses beside the node's (LINKAGE rule 6).
\statusT\quad\emph{(Cites nothing of its own. Three declarations require less than the node, their
unused hypotheses dropped on 2026-11-14. \textbf{Hypotheses} (the node's, each with what its
declarations require): $q > 1$: same for \texttt{surface\_measure\_family}, \texttt{\_pair},
\texttt{\_injective} and \texttt{\_integrable\_iff}, absent from \texttt{surface\_measure\_unique}
and \texttt{\_exponent}; a $q^{\ZZ}$-covariant family in the canonical discrete gauge: same
(\texttt{IsQFamily}) for \texttt{surface\_measure\_family}, \texttt{\_injective} and
\texttt{\_exponent}, weaker in \texttt{surface\_measure\_unique}, which requires only that the kernels
$\mu_{0,t}$ have the data of the pair $(\Lambda, \mathcal{B})$ and of the competitor, the pair being
fixed by the transforms of the kernels alone, and of the competitor only that it be carried by
$(0,\infty)^2$ and $(0,\infty)$; $\Lambda$ a positive measure on $(0,\infty)^2$: same, a measure on
$\RR^2$ null off the open quadrant, not required by \texttt{surface\_measure\_integrable\_iff};
$\mathcal{B}$ a positive measure on $(0,\infty)$: same, a measure on $\RR$ null on $(-\infty,0]$;
(i)--(iv): same (\texttt{IsSurfacePair}), the integrability of (i) required at every $t > 0$, and
\texttt{surface\_measure\_integrable\_iff} requires only $J_*\Lambda = \Lambda$ for ``some, equivalently
every''.)}
\end{theorem}

% CHANGED (wave 3a, 2026-11-14): the proof follows the machine-checked route. The averaging bound
% averages over omega in (0,2] with the term 3(1 - cos v); the measure Lambda is built as a measure
% W(dv) kappa_v(dy) on each strip (0,T], from the Radon-Nikodym densities of the varpi_t against
% varpi_T, in place of the Lebesgue-Stieltjes theorem in two variables, and the strips are glued by
% the uniqueness; the gluing of the converse is a change of variables under J.
% CHANGED (the theory paper, 2026-11-16): the closing sentence on checking and cited facts shortened; no step changed.
% CHANGED (supplement review, 2026-11-16): the averaging bound is cited from lem:averaging-bound.
\begin{proof}
\leanok
\uses{thm:annulus-parametrization, lem:line-exponent-data, lem:family-structure,
lem:discrete-covariance-levy, lem:averaging-bound}
\emph{From a family to $(\Lambda, \mathcal{B})$: the data.} Let the kernels have the data
$(a(t), \varpi_t)$, $t \ge 0$, so that $(a(0), \varpi_0) = (0, 0)$, as $\mu_0 = \delta_0$. For
$s < t$ the step $\mu_{s,t}$ is symmetric and infinitely divisible
(Lemma~\ref{lem:family-structure}) with the exponent $\Psi_t - \Psi_s$, by (I1); so
$\Psi_t - \Psi_s \in \mathrm{ND}_s$, and by Lemma~\ref{lem:line-exponent-data}(3)
$a(s) \le a(t)$, $\varpi_s \le \varpi_t$, and $\Psi_t - \Psi_s$ has the data
$(a(t) - a(s), \varpi_t - \varpi_s)$. Covariance at the pairs $(0, t)$ gives
$\Psi_{qt} = \Psi_t(q\,\cdot)$, so $(a(qt), \varpi_{qt}) = (q^2a(t), T_q\varpi_t)$ by
Lemma~\ref{lem:line-exponent-data}(1),(2). Put
$G(t) := 3\Psi_t(1) + \int_0^2\Psi_t(\omega)\,d\omega$. The averaging bound,
Lemma~\ref{lem:averaging-bound}(2), for $\Psi_t - \Psi_s$ gives
\[
  a(t) - a(s) \le \Psi_t(1) - \Psi_s(1), \qquad
  (1 \wedge \varepsilon^2)\,(\varpi_t(A) - \varpi_s(A)) \le G(t) - G(s)
\]
for every Borel $A \subseteq (\varepsilon, \infty)$, $\varepsilon > 0$. By (I2) and
$\widehat\mu_t > 0$, $t \mapsto \Psi_t(\omega)$ is continuous on $[0,\infty)$; and
$0 \le \Psi_t \le \Psi_T$ for $t \le T$ with $\Psi_T$ continuous, so $G$ is continuous by
dominated convergence. Hence $a$ is continuous and nondecreasing, and so is $t \mapsto \varpi_t(A)$
for every Borel $A$ bounded away from $0$.

\emph{The measure $\Lambda$.} Fix $T > 0$ and put $W := \varpi_T$, a measure on $(0,\infty)$
finite on every $(\varepsilon,\infty)$, hence $\sigma$-finite. For rational $r$ let
$h_r := 1 \wedge d\varpi_r/dW$ when $0 < r < T$, $h_r := 0$ when $r \le 0$ and $h_r := 1$ when
$r \ge T$. Then $\int_A h_r\,dW = \varpi_{r \wedge T}(A)$ for Borel $A$, since
$d\varpi_r/dW \le 1$ $W$-a.e., and for rationals $r' \le r$ we have $h_{r'} \le h_r$ $W$-a.e., as
$\varpi_{r'} \le \varpi_r$. Put $f_r(v) := \sup_{r' \le r}h_{r'}(v)$ and
$F_r(v) := \inf_{r' > r}f_{r'}(v)$, over rationals. Then $f_r = h_r$ $W$-a.e., and for every $v$
the function $r \mapsto F_r(v)$ is nondecreasing and right-continuous on the rationals, vanishes for
$r < 0$ and equals $1$ for $r \ge T$. So it extends to a distribution function on $\RR$, whose
law $\kappa_v$ depends measurably on $v$, with $\kappa_v((-\infty, t]) = \inf_{r > t}f_r(v)$. For
Borel $A \subseteq (\varepsilon,\infty)$ and $t \ge 0$, monotone convergence (the $f_r$ are at most
$1$ and $W(A) < \infty$) and the right-continuity of $t \mapsto \varpi_t(A)$ give
\[
  \int_A\kappa_v\bigl((-\infty, t]\bigr)\,W(dv) = \inf_{r > t}\int_A f_r\,dW
  = \inf_{r > t}\varpi_{r \wedge T}(A) = \varpi_{t \wedge T}(A).
\]
Let $\Lambda_T$ be the image of the measure $W(dv)\,\kappa_v(dy)$ under $(v, y) \mapsto (y, v)$. At
$t = 0$ the display gives $\Lambda_T((-\infty,0] \times (\varepsilon,\infty)) = 0$, and
$\Lambda_T(\RR \times (-\infty, 0]) = W((-\infty,0]) = 0$, so $\Lambda_T$ is carried by
$(0,\infty)^2$; and $\Lambda_T((0,t] \times A) = \varpi_{t \wedge T}(A)$ for $t > 0$. Two measures
on $(0,\infty)$ that agree on the Borel sets bounded away from $0$ are equal, so
$\Lambda_T((0,t] \times \cdot\,) = \varpi_{t \wedge T}$.

\emph{Uniqueness on the quadrant.} Two measures $\Lambda, \Lambda'$ on $(0,\infty)^2$ with the
same slabs $\Lambda((0,t] \times \cdot\,)$, finite on every $(\varepsilon,\infty)$, are equal.
Indeed, for Borel $A \subseteq (\varepsilon,\infty)$ the measures $S \mapsto \Lambda(S \times A)$
and $S \mapsto \Lambda'(S \times A)$ on $(0,\infty)$ agree, finitely, on the intervals $(0,t]$,
hence on all Borel sets. The products $S \times A$ with
$A \subseteq (-\infty,0] \cup (\varepsilon,\infty)$ for some $\varepsilon > 0$ form a
$\pi$-system that generates the Borel sets of $\RR^2$; both measures agree on them, and countably
many of them, the sets $(-n,n] \times ((-\infty,0] \cup (1/(n+1),\infty))$, have finite measure
and cover $\RR^2$. So $\Lambda = \Lambda'$. In particular, for $m \le n$ the restriction of
$\Lambda_{n+1}$ to $(0, m+1] \times \RR$ is $\Lambda_{m+1}$, and
$\Lambda := \sum_{n \ge 0}\Lambda_{n+1}|_{(n, n+1] \times \RR}$ has
$\Lambda((0,t] \times \cdot\,) = \varpi_t$ for every $t > 0$.

Let $\mathcal{B}$ be the Stieltjes measure of the continuous nondecreasing function
$t \mapsto a(t \vee 0)$. Then $a(t) = \mathcal{B}((0,t])$ and $\mathcal{B}((-\infty,0]) = 0$.
This is (\ref{eq:line-e1}), and (\ref{eq:line-e2}) follows by substituting. Any $\Lambda'$ on $(0,\infty)^2$ and
$\mathcal{B}'$ on $(0,\infty)$ with (\ref{eq:line-e1}) have the slabs $\Lambda'((0,t] \times \cdot\,)$ and the
masses $\mathcal{B}'((0,t])$ of the data of $\mu_t$, which are unique
(Lemma~\ref{lem:line-exponent-data}(1)). So $\Lambda' = \Lambda$ by the uniqueness on the
quadrant, and $\mathcal{B}' = \mathcal{B}$, two measures on $(0,\infty)$ that agree, finitely, on
every $(0,t]$ being equal.

(i) $J^{-1}((0,t] \times A) = (0,t/q] \times q^{-1}A$, so the slabs of $J_*\Lambda$ are
$T_q\varpi_{t/q} = \varpi_t$, and $J_*\Lambda = \Lambda$ by the uniqueness on the quadrant.
Integrability at every $t$ holds because $\varpi_t$ is a \Levy{} measure. For a general invariant
$\Lambda$,
\[
  \int(1 \wedge v^2)\,\Lambda((0,qt] \times dv) = \int(1 \wedge q^2v^2)\,\Lambda((0,t] \times dv)
  \le q^2\int(1 \wedge v^2)\,\Lambda((0,t] \times dv),
\]
and the integral is nondecreasing in $t$. So integrability at one $t$ gives it at all.
(ii) The measures $A \mapsto \mathcal{B}(qA)$ and $q^2\mathcal{B}$ on $(0,\infty)$ agree, finitely,
on the intervals $(0,t]$, since $a(qt) = q^2a(t)$, hence on all Borel sets.
(iii) For $0 < s < y$ and $\varepsilon > 0$,
$\Lambda(\{y\} \times (\varepsilon,\infty)) \le \Lambda((s,y] \times (\varepsilon,\infty))
= \varpi_y((\varepsilon,\infty)) - \varpi_s((\varepsilon,\infty))
\le (G(y) - G(s))/(1 \wedge \varepsilon^2)$, which tends to $0$ as $s \uparrow y$. So
$\Lambda(\{y\} \times (\varepsilon,\infty)) = 0$ for every $\varepsilon > 0$, and
$\Lambda(\{y\} \times (0,\infty)) = 0$. $\mathcal{B}$ has no atoms because $a$ is continuous.
(iv) If $\Lambda((s,t] \times (0,\infty)) = 0 = \mathcal{B}((s,t])$, then $a(t) = a(s)$ and
$\varpi_t = \varpi_s$, the data of $\mu_{s,t}$ vanish, and $\mu_{s,t} = \delta_0$, against (ND).

\emph{From $(\Lambda, \mathcal{B})$ to a family.} Given a pair with (i)--(iv), define $a(t)$ and
$\varpi_t$ by (\ref{eq:line-e1}). By (i), $\varpi_t$ is a \Levy{} measure on $(0,\infty)$, and $\Psi_t$ given
by (\ref{eq:line-e2}) lies in $\mathrm{ND}_s$ with the data $(a(t), \varpi_t)$. For $0 \le s < t$,
$\Psi_t \ne \Psi_s$: otherwise their data agree (Lemma~\ref{lem:line-exponent-data}(1)), so
$\mathcal{B}((s,t]) = 0$ and $\Lambda((s,t] \times (\varepsilon,\infty)) = 0$ for every
$\varepsilon > 0$, the masses being finite, and (iv) fails on $(s \vee t/2, t]$. With $s = 0$, where
$\Psi_0 = 0$, this gives $\Psi_t \not\equiv 0$. For $t < t'$ in $[1,q]$ the data are ordered,
$a(t) \le a(t')$ and $\varpi_t \le \varpi_{t'}$, so condition (2) of
Theorem~\ref{thm:annulus-parametrization} holds by Lemma~\ref{lem:line-exponent-data}(3).
Gluing: by $J_*\Lambda = \Lambda$ and the change of variables $(y, v) = J(y', v')$,
$\int_{(0,qt] \times (0,\infty)}(1 - \cos\omega v)\,\Lambda(dy, dv)
= \int_{(0,t] \times (0,\infty)}(1 - \cos q\omega v')\,\Lambda(dy', dv')$, and
$a(qt) = \mathcal{B}(q(0,t]) = q^2a(t)$ by (ii), so $\Psi_{qt} = \Psi_t(q\,\cdot)$ for every
$t > 0$, in particular $\Psi_q = \Psi_1(q\,\cdot)$. Continuity: on $(0,T] \times (0,\infty)$ the
integrand of (\ref{eq:line-e2}) is integrable, its integral being at most the exponent of the data at $T$; by
(iii) no vertical line carries mass of $\Lambda$, and $\mathcal{B}$ has no atoms. So $a(t)$ and the
integral in (\ref{eq:line-e2}) are continuous in $t$ by dominated convergence, the indicator of $(0,t]$
converging off the line $\{t_0\} \times (0,\infty)$ as $t \to t_0$.
Theorem~\ref{thm:annulus-parametrization} gives a family whose exponents at every $t \ge 0$ are the
path extended by $\Psi_{q^nt} := \Psi_t(q^n\,\cdot)$, which by $\Psi_{qt} = \Psi_t(q\,\cdot)$ is
(\ref{eq:line-e2}) itself. So its kernels have the data (\ref{eq:line-e1}), and its pair, constructed in the first part, is
$(\Lambda, \mathcal{B})$ by the uniqueness. The map from families to pairs is injective: two
families with the data (\ref{eq:line-e1}) have the same exponents at every $t > 0$, so their kernels agree on
$[1,q)$ and the families coincide (Lemma~\ref{lem:discrete-covariance-levy}). So the two maps are
inverse bijections. (\ref{eq:line-e1}) and (\ref{eq:line-e2}) are equivalent for a pair with (i)--(iv), since the
exponent of the data in (\ref{eq:line-e1}) is (\ref{eq:line-e2}) and the data of an exponent are unique.

Every step is machine-checked, and the cited facts enter only through the statements the proof
cites.
\end{proof}

The extreme members of a cone of measures are found by splitting a measure along a Borel set. A
measure that no such split decomposes is a point mass.

% draft: none. New at the supplement review (2026-11-16): the argument "every Borel set null or
% conull gives a point mass", proved inside the proofs of thm:knot-atom-cascades(3) and
% prop:strata-cylinder-coordinates-q-cin(1), stated once and cited there.
\begin{lemma}[a measure with no proper part is a point mass]
\label{lem:null-conull-point-mass}
\lean{Pyramids.eq_smul_dirac_of_zero_one, Pyramids.eq_smul_dirac_prod_of_zero_one}
\leanok
Let $\kappa$ be a measure on $\RR$ or on $\RR^2$ such that every Borel set $E$ has $\kappa(E) = 0$
or $\kappa(E^c) = 0$. Then $\kappa = \kappa(\{p\})\,\delta_p$ for some point $p$.
\statusT\quad\emph{(Cites nothing. \textbf{Hypotheses} (the node's, each with what its declarations
require): $\kappa$ a measure on $\RR$, every Borel set null or conull: same
(\texttt{eq\_smul\_dirac\_of\_zero\_one}); the same on $\RR^2$: same
(\texttt{eq\_smul\_dirac\_prod\_of\_zero\_one}).)}
\end{lemma}

\begin{proof}
\leanok
On $\RR$, let $S := \{x : \kappa((x,\infty)) = 0\}$. If $S$ is empty, then
$\kappa((-\infty, n]) = 0$ for every integer $n$, and if $S$ is unbounded below, then
$\kappa((-n,\infty)) = 0$ for every $n$; in both cases $\kappa = 0$ by continuity from below, and
any $p$ will do. Otherwise $\tau := \inf S$ is finite. Every $x > \tau$ lies in $S$, so
$\kappa((\tau,\infty)) = 0$, and every $x < \tau$ lies outside $S$, so $\kappa((-\infty,x]) = 0$ and
$\kappa((-\infty,\tau)) = 0$. So $\kappa$ is carried by $\{\tau\}$, and $\kappa = \kappa(\{\tau\})\,
\delta_\tau$.

On $\RR^2$, the images of $\kappa$ under the two coordinate maps have the same property, so they are
carried by points $\theta$ and $r$. The complement of $\{(\theta, r)\}$ lies in the union of
$(\RR \setminus \{\theta\}) \times \RR$ and $\RR \times (\RR \setminus \{r\})$, both $\kappa$-null,
so $\kappa$ is carried by $(\theta, r)$.

Every step is machine-checked, with no cited fact.
\end{proof}

% draft: Theorem 2.11 (section 2.3), lines 187-191 of draft/spatial-hemigroup-pyramids.md.
% CHANGED: "the extreme rays of the cone of q^Z-covariant families are the knot-atom cascades" is
% false as written: a knot-atom cascade jumps at its knots, so it violates (I2) and is not a family.
% The statement now places the extreme rays in the cone of invariant surface measures, says that
% the knot-atom cascades are its members and not families, and says which mixtures are families.
% CHANGED: "free and proper" stated as free and properly discontinuous, with the orbit space and a
% fundamental domain.
% CHANGED: the generator reading ([V-cone v0.1, section 6, Prop. 6.4(5)]) moves to the prose; the
% statement says what happens to the kernel at a knot.
% CHANGED: the integrability is proved here from Theorem 2.10(i), with explicit two-sided bounds;
% this answers the draft's check "against the integrability of [V-cone v0.1, Prop. 3.3]" (its
% comparison at Haar sigma is prop:strata-cylinder-coordinates(1)).
% CHANGED: clause (4) says that the knot Gaussians are not families either (blind review 2, flag 15).
\begin{theorem}[the knot-atom cascades]
\label{thm:knot-atom-cascades}
\lean{Pyramids.knot_atom_cascades_1_action, Pyramids.knot_atom_cascades_1_orbits,
Pyramids.knot_atom_cascades_2_bij, Pyramids.knot_atom_cascades_2_integrable,
Pyramids.knot_atom_cascades_2_gaussian, Pyramids.knot_atom_cascades_3,
Pyramids.knot_atom_cascades_3_gaussian, Pyramids.knot_atom_cascades_4,
Pyramids.knot_atom_cascades_4_gaussian, Pyramids.knot_atom_cascades_4_families}
\leanok
\uses{thm:surface-measure}
For $y > 0$ let $[y] \in [1,q)$ be its representative modulo $q^{\ZZ}$, and let
$\mathcal{C} := [1,q)_\theta \times (0,\infty)_r$, the \emph{cylinder}.
\begin{enumerate}
\item The group $\{J^n : n \in \ZZ\}$ acts freely and properly discontinuously on $(0,\infty)^2$.
The invariant map $(y,v) \mapsto (\log y \operatorname{mod} \ell,\ v/y)$ induces a homeomorphism of the orbit
space onto $(\RR/\ell\ZZ) \times (0,\infty)$, and $[1,q) \times (0,\infty)$ meets every orbit
exactly once. The ratio $r := v/y$ is the invariant \emph{displacement-to-scale ratio}.
\item Let $\sigma$ be the image of the restriction of $\Lambda$ to $[1,q) \times (0,\infty)$ under
$(y,v) \mapsto (y, v/y)$. Then $\Lambda \mapsto \sigma$ is a bijection of the $J$-invariant positive
measures on $(0,\infty)^2$ onto the positive measures on $\mathcal{C}$, with inverse
\[
  \sigma \mapsto \Lambda_\sigma := \int_{\mathcal{C}}\Lambda^{(\theta,r)}\,\sigma(d\theta,dr),
  \qquad \Lambda^{(\theta,r)} := \sum_{n \in \ZZ}\delta_{(\theta q^n,\, r\theta q^n)} .
\]
$\Lambda_\sigma$ satisfies the integrability of Theorem~\ref{thm:surface-measure}(i) iff
\begin{equation}\label{eq:line-e3}
  \int_{\mathcal{C}}\Bigl[(r\theta)^2 \wedge \bigl(1 + \log_+(r\theta)\bigr)\Bigr]\,
  \sigma(d\theta,dr) < \infty ,
\end{equation}
and it has no mass on a line $\{y\} \times (0,\infty)$ iff
$\sigma(\{\theta\} \times (0,\infty)) = 0$ for every $\theta$. Likewise $\mathcal{B}$ with
Theorem~\ref{thm:surface-measure}(ii) corresponds to its restriction $\beta$ to $[1,q)$, by
$\mathcal{B}(A) = \sum_nq^{2n}\beta(q^{-n}A \cap [1,q))$.
\item The extreme rays of the cone of $J$-invariant $\Lambda$ with the integrability (\ref{eq:line-e3}) are the
rays of the $\Lambda^{(\theta,r)}$, $(\theta,r) \in \mathcal{C}$, and every such $\Lambda$ is the
mixture $\Lambda_\sigma$ with $\sigma$ unique. With $\mathcal{B} = 0$, the exponents (\ref{eq:line-e2}) of
$\Lambda^{(\theta,r)}$ are the \emph{knot-atom cascade}
\begin{equation}\label{eq:line-e4}
  \Psi^{(\theta,r)}_t(\omega) = \sum_{n :\ \theta q^n \le t}\bigl(1 - \cos(\omega\,r\theta q^n)\bigr).
\end{equation}
The extreme rays of the cone of the $\mathcal{B}$ are those of
$\mathcal{B}^{(\theta)} := \sum_n(\theta q^n)^2\delta_{\theta q^n}$, the \emph{knot Gaussians}
$a(t) = \sum_{n :\ \theta q^n \le t}(\theta q^n)^2$.
\item A knot-atom cascade is constant in $t$ between consecutive knots $\theta q^n$ of phase
$\theta$, and at each of them its kernel is convolved with the law of transform
$e^{-(1 - \cos(r\theta q^n\omega))}$, the compound Poisson law of unit rate with jumps
$\pm r\theta q^n$. It violates (I2), so it is a member of the cone and not a family. A knot
Gaussian jumps by $(\theta q^n)^2$ at each knot $\theta q^n$, so its $\mathcal{B}^{(\theta)}$ has
atoms, against Theorem~\ref{thm:surface-measure}(iii), and it is not a family either. The families
are the pairs $(\Lambda_\sigma, \mathcal{B})$ with (\ref{eq:line-e3}), with $\sigma(\{\theta\} \times (0,\infty)) = 0$
for every $\theta$, $\beta$ without atoms, and Theorem~\ref{thm:surface-measure}(iv).
\end{enumerate}
% CHANGED (fidelity review, 2026-11-15, R64): the status line says that the orbit-space
% homeomorphism of (1) is stated in decomposed form.
% CHANGED (fidelity review, 2026-11-15, R26/R50/R69): the status line lists the declarations' hypotheses beside the node's (LINKAGE rule 6).
\statusT\quad\emph{(\textbf{Hypotheses} (the node's, each with what its declarations require):
$q > 1$: same, in every declaration; (2) $\Lambda$ a $J$-invariant positive measure on
$(0,\infty)^2$: same, null off the quadrant with $J_*\Lambda = \Lambda$; $\sigma$ a positive measure
on $\mathcal{C}$: same, null off the cylinder; $\mathcal{B}$ with (ii): same
(\texttt{IsSurfaceGaussian}); $\beta$ the restriction of such a $\mathcal{B}$: same in effect,
\texttt{knot\_atom\_cascades\_2\_gaussian} requiring $\beta$ finite and carried by $[1,q)$, which the
restriction of a $\mathcal{B}$ with (ii) is; (3) the cones: same, \texttt{surfaceCone} ($J$-invariant,
carried by the quadrant, (\ref{eq:line-e3})) and \texttt{gaussianCone} ((ii)); (3) and (4) $(\theta, r) \in
\mathcal{C}$: same, and $t > 0$ for the cascade; (4) the families: same, $\sigma$ carried by the
cylinder and $\mathcal{B}$ with (ii) in \texttt{knot\_atom\_cascades\_4\_families}. The homeomorphism
of the orbit space in (1) is stated in decomposed form, the reader composing it:
\texttt{knot\_atom\_cascades\_1\_orbits} gives that the invariant map is continuous and open on the
quadrant, onto $(\RR/\ell\ZZ) \times (0,\infty)$, with the $J^{\ZZ}$-orbits as its fibres.)}
\end{theorem}

% CHANGED (wave 3a, 2026-11-14): the proof follows the machine-checked route. In (1) the induced map
% is a homeomorphism because p is open, the composite of the homeomorphism (y,v) -> (log y, v/y) onto
% R x (0,inf) with the open quotient map. In (2) the integrability is compared with the explicit
% constants of the checked bounds and holds at every t by Theorem surface-measure(i). In (3) the
% extreme members of the cone of sigma are found by splitting along a Borel set, which makes every
% Borel set null or conull, and such a measure is carried by one point through its two marginals,
% in place of the support argument. In (4) the transform of the compound Poisson law is computed
% and the failure of (I2) is shown at one frequency.
% CHANGED (the theory paper, 2026-11-16): the closing sentence on checking and cited facts shortened; no step changed.
% CHANGED (supplement review, 2026-11-16): in (3) the point-mass argument is cited from
% lem:null-conull-point-mass; the weight w(x) renamed Theta(x) (w is the mask).
\begin{proof}
\leanok
\uses{thm:surface-measure, lem:null-conull-point-mass}
(1) $J^n(y,v) = (y,v)$ forces $q^ny = y$, so $n = 0$: the action is free. For a compact
$K \subset (0,\infty)^2$ with $y$-coordinates in $[y_1, y_2]$, $J^nK \cap K \ne \varnothing$ needs
$q^n \in [y_1/y_2, y_2/y_1]$, which holds for finitely many $n$. The map
$p(y,v) := (\log y \operatorname{mod} \ell, v/y)$ is continuous and invariant. It is the map
$\phi(y,v) := (\log y, v/y)$ followed by the quotient map $\RR \to \RR/\ell\ZZ$ on the first
coordinate. The map $\phi$ is a homeomorphism of $(0,\infty)^2$ onto $\RR \times (0,\infty)$, with
inverse $(u, r) \mapsto (e^u, re^u)$, and the quotient map of a topological group is open. So $p$
is open, and it is onto $(\RR/\ell\ZZ) \times (0,\infty)$. Two points with the same image have
$\log y' - \log y \in \ell\ZZ$, that is $y' = q^ny$, and $v'/y' = v/y$, so they lie in one orbit.
A continuous open surjection whose fibres are the orbits induces a homeomorphism of the orbit space.
Each orbit $\{(q^ny, q^nv)\}$ has exactly one $n$ with $q^ny \in [1,q)$, namely
$n = -\lfloor\log_qy\rfloor$.

(2) Write $F := [1,q) \times (0,\infty)$. The sets $J^nF$, $n \in \ZZ$, partition $(0,\infty)^2$, so
an invariant $\Lambda$ satisfies $\Lambda(E) = \sum_n\Lambda(J^{-n}E \cap F)$ and is determined by
its restriction to $F$. Conversely, for any measure on $F$ that sum defines an invariant measure.
Transporting $F$ to $\mathcal{C}$ by $(y,v) \mapsto (y, v/y)$ gives the bijection. Restricted to $F$,
$\Lambda^{(\theta,r)}$ is the single atom at $(\theta, r\theta)$, since only $n = 0$ has
$\theta q^n \in [1,q)$. So the restriction of $\Lambda_\sigma$ to $F$ is the transport of $\sigma$,
and $\Lambda_\sigma$ is the inverse image.

For the integrability, by Tonelli and Theorem~\ref{thm:surface-measure}(i) at $t = 1$,
\[
  \int(1 \wedge v^2)\,\Lambda_\sigma((0,1] \times dv)
  = \int_{\mathcal{C}}\Bigl[f(r\theta) + \mathbf 1_{\{\theta = 1\}}\bigl(1 \wedge (r\theta)^2\bigr)\Bigr]
  \sigma(d\theta,dr), \qquad f(x) := \sum_{m \ge 1}\bigl(1 \wedge x^2q^{-2m}\bigr),
\]
because $\theta q^n \le 1$ means $n \le -1$, or $n \le 0$ when $\theta = 1$. Write
$\Theta(x) := x^2 \wedge (1 + \log_+x)$, and note $\log_+x \le \Theta(x)$ and $1 \wedge x^2 \le \Theta(x)$ for
$x \ge 0$ (as $\log x \le x - 1 \le x^2$). On the one side, the terms of $f$ with $q^m \le x$ number
at most $\log_+x/\ell$, and the others are at most $(1 \wedge x^2)q^{-2j}$, $j \ge 0$, so
$f(x) \le \log_+x/\ell + \tfrac{q^2}{q^2 - 1}(1 \wedge x^2) \le C_q\Theta(x)$ with
$C_q := 1/\ell + q^2/(q^2 - 1)$. On the other side, for $x \le q$ the term $m = 1$ gives
$f(x) \ge x^2/q^2$, so $\Theta(x) \le x^2 \le q^2f(x)$. For $x > q$ the term $m = 1$ is $1$, so
$f(x) \ge 1$, and the terms equal to $1$ number $\lfloor\log_qx\rfloor$, so
$\log x \le \ell(f(x) + 1) \le 2\ell f(x)$ and $\Theta(x) \le 1 + \log x \le (1 + 2\ell)f(x)$. So the
integral is finite iff (\ref{eq:line-e3}) holds, and then it is finite at every $t > 0$ by
Theorem~\ref{thm:surface-measure}(i), $\Lambda_\sigma$ being invariant. A line
$\{y\} \times (0,\infty)$ meets the orbit of $(\theta, r)$ only if $\theta = q^{-n}y$ for some $n$,
so its mass under $\Lambda_\sigma$ is $\sum_n\sigma(\{q^{-n}y\} \times (0,\infty))$, and conversely
$\sigma(\{\theta\} \times (0,\infty))$ is at most the mass of the line $\{\theta\} \times (0,\infty)$,
the term $n = 0$. The statement on $\mathcal{B}$ is the same argument for the action $y \mapsto qy$
with the weight $q^2$: the condition $\mathcal{B}(qA) = q^2\mathcal{B}(A)$ gives
$q^{2n}\mathcal{B}(q^{-n}A) = \mathcal{B}(A)$ for every $n \in \ZZ$, the intervals
$q^n[1,q)$ partition $(0,\infty)$, and for $\beta$ finite on $[1,q)$ the sum
$\sum_nq^{2n}\beta(q^{-n}(0,t] \cap [1,q))$ has no terms with $q^n > t$ and a geometric tail.

(3) By (2) the map $\sigma \mapsto \Lambda_\sigma$ is a bijection, additive and positively
homogeneous, with inverse $\Lambda \mapsto \sigma$, also additive and positively homogeneous. So
$\Lambda = \Lambda_\sigma$ with $\sigma$ unique, and a decomposition $\Lambda = \Lambda' + \Lambda''$
in the cone is the decomposition $\sigma = \sigma' + \sigma''$ of the cylinder measures, each with
(\ref{eq:line-e3}). Let $\Lambda_\sigma$ span an extreme ray, and let $E \subset \mathcal{C}$ be Borel. Then
$\sigma = \sigma|_E + \sigma|_{E^c}$, both terms satisfy (\ref{eq:line-e3}), and extremality gives
$\sigma|_E = c\sigma$ for some $c \ge 0$. Evaluated on $E^c$ this is $0 = c\,\sigma(E^c)$, so either
$\sigma(E^c) = 0$, or $c = 0$ and $\sigma(E) = 0$: every Borel set is $\sigma$-null or
$\sigma$-conull. By Lemma~\ref{lem:null-conull-point-mass}, $\sigma = c\,\delta_{(\theta,r)}$ with
$c = \sigma(\{(\theta,r)\})$ for a point $(\theta, r)$, and $c > 0$ as $\sigma \ne 0$. The point
lies in $\mathcal{C}$, and $c < \infty$ by (\ref{eq:line-e3}), the weight $\Theta(r\theta)$ being positive. So
$\Lambda = c\,\Lambda^{(\theta,r)}$. Conversely, $c\,\delta_{(\theta,r)}$ with $0 < c < \infty$
satisfies (\ref{eq:line-e3}), and in a decomposition $c\,\delta_{(\theta,r)} = \sigma' + \sigma''$ the measure
$\sigma'$ vanishes off $(\theta, r)$, so it is a multiple $c'\delta_{(\theta,r)}$, $c' \le c$, and
$\Lambda_{\sigma'} = (c'/c)\,c\Lambda^{(\theta,r)}$. Inserting $\Lambda^{(\theta,r)}$ in (\ref{eq:line-e2}), the
integral over $(0,t] \times (0,\infty)$ is the sum over the atoms $(\theta q^n, r\theta q^n)$ with
$\theta q^n \le t$, which is (\ref{eq:line-e4}). The sum converges, since its terms vanish for $\theta q^n > t$
and are at most $\tfrac12(\omega r\theta)^2q^{2n}$, by $1 - \cos u \le u^2/2$. For $\mathcal{B}$ the
same argument applies to $\beta$, on $\RR$, where $\beta$ is finite. The point mass at $\theta$
gives $\sum_nq^{2n}\delta_{\theta q^n}$, which is $\theta^{-2}\mathcal{B}^{(\theta)}$, and
$\mathcal{B}^{(\theta)}((0,t])$ is the sum of the atoms $(\theta q^n)^2$ with $\theta q^n \le t$.

(4) For $t \in [\theta q^n, \theta q^{n+1})$, $\theta q^m \le t$ holds iff $m \le n$. So the sum
(\ref{eq:line-e4}) is constant on $[\theta q^n, \theta q^{n+1})$, and from $t \in [\theta q^{n-1}, \theta q^n)$
to $\theta q^n$ it gains exactly the term $m = n$, $1 - \cos(\omega r\theta q^n)$. That function is
the exponent of the compound Poisson law
$\exp[\delta^2_v] = e^{-1}\sum_{k \ge 0}(\tfrac12(\delta_v + \delta_{-v}))^{*k}/k!$, $v = r\theta q^n$:
the step $\tfrac12(\delta_v + \delta_{-v})$ has the transform $\cos v\omega$, so the law has the
transform $e^{-1}\sum_k(\cos v\omega)^k/k! = e^{-(1 - \cos v\omega)}$. The transforms multiply, so
the kernel at $\theta q^n$ is the kernel at $t$ convolved with $\exp[\delta^2_v]$, the transform
determining the law. The cascade violates (I2): if a family of laws $\mu_{s,t}$ with (I1), (I2) had
$\widehat{\mu_{0,t}} = e^{-\Psi^{(\theta,r)}_t}$, then at $\omega = \pi/(r\theta)$ the function
$t \mapsto \widehat{\mu_{0,t}}(\omega)$ would be continuous at $t = \theta$; it is constant on
$[\theta/q, \theta)$, and at $\theta$ the exponent gains $1 - \cos\pi = 2$, so the value there differs
from the limit from the left. The knot Gaussian $a(t) = \sum_{\theta q^n \le t}(\theta q^n)^2$
jumps by $(\theta q^n)^2$ at $\theta q^n$, since the points $\theta q^m$ are distinct, so
$\mathcal{B}^{(\theta)}(\{\theta q^n\}) > 0$, against Theorem~\ref{thm:surface-measure}(iii): no pair
with (i)--(iv) has this $\mathcal{B}$. The description of the families is (i)--(iv) read through
(2): $\Lambda_\sigma$ is invariant and carried by $(0,\infty)^2$, its integrability is (\ref{eq:line-e3}), and it
has no mass on a vertical line iff $\sigma(\{\theta\} \times (0,\infty)) = 0$ for every $\theta$;
$\mathcal{B}$ has no atoms iff $\beta$ has none, since
$\mathcal{B}(\{y\}) = \sum_nq^{2n}\beta(\{q^{-n}y\})$ and an atom of $\beta$ at $y \in [1,q)$ is an
atom of $\mathcal{B}$.

Every step is machine-checked, with no cited fact.
\end{proof}

In the generator form (\ref{eq:generator-form}), with the symmetric second difference
$\delta^2_h$ defined there, the knot-atom cascade applies one step $\delta^2_{r\theta q^n}$ at each
knot and does nothing between knots. The Cin ray's generator is a single symmetric second
difference at every scale \prog{V-cone v0.1}{Prop.~6.4(5), p.~58}.

% draft: Proposition 2.12 (section 2.3), lines 193-197 of draft/spatial-hemigroup-pyramids.md;
% clauses (1)-(3) here, the cone of span-q exponents in the next node.
% CHANGED: split from Proposition 2.12: the marginal cone and the q-Cin rays are
% prop:strata-cylinder-coordinates-q-cin.
% CHANGED: "fully covariant" made precise as condition (a), Psi_{lambda t} = Psi_t(lambda .), the
% Gaussian part added (the draft's statement omits B), and the equivalence with the profile k and
% its integrability added in (1); "modulation content" renamed "modulation spectrum" (review flag,
% one term for one thing) and stated as clause (3).
% CHANGED: clause (1) split into the jump part and the Gaussian part, each with its own three
% equivalent conditions, so that the nodes that use the jump part alone cite a statement (blind
% review 2, flag 9). Condition (a) of the draft, Psi_{lambda t} = Psi_t(lambda .), is the
% conjunction of the two new (a)'s, the last sentence of the clause.
\begin{proposition}[the strata in cylinder coordinates; circle averaging]
\label{prop:strata-cylinder-coordinates}
\lean{Pyramids.strata_cylinder_1_jump, Pyramids.strata_cylinder_1_then,
Pyramids.strata_cylinder_1_integrability, Pyramids.strata_cylinder_1_gaussian,
Pyramids.strata_cylinder_1_both, Pyramids.strata_cylinder_2, Pyramids.strata_cylinder_3}
\leanok
\uses{thm:knot-atom-cascades, thm:surface-measure, lem:line-exponent-data}
Let $\eta(d\theta) := d\theta/(\ell\theta)$ on $[1,q)$, the Haar probability measure of the circle
$(0,\infty)/q^{\ZZ}$. Let a $q^{\ZZ}$-covariant family in the canonical discrete gauge have the pair
$(\Lambda, \mathcal{B})$ of Theorem~\ref{thm:surface-measure}, with
$\varpi_t = \Lambda((0,t] \times \cdot\,)$ and $a(t) = \mathcal{B}((0,t])$, and the cylinder
measures $(\sigma, \beta)$ of Theorem~\ref{thm:knot-atom-cascades}. Write
$J_\lambda(y,v) := (\lambda y, \lambda v)$ for $\lambda > 0$.
\begin{enumerate}
\item (The jump part.) The following are equivalent: (a) $\varpi_{\lambda t} = T_\lambda\varpi_t$
for all $\lambda, t > 0$; (b) $(J_\lambda)_*\Lambda = \Lambda$ for every $\lambda > 0$;
(c) $\sigma = \eta \otimes \tilde\rho$ for a measure $\tilde\rho$ on $(0,\infty)$. Then $\Lambda$ is
the image of $(dy/(\ell y)) \otimes \tilde\rho$ under $(y,r) \mapsto (y, ry)$, and
$\varpi_t(dv) = k(v/t)\,dv/v$ with $k(x) := \ell^{-1}\tilde\rho([x,\infty))$, nonincreasing. The
integrability (\ref{eq:line-e3}) is equivalent to each of $\int(r^2 \wedge (1 + \log_+r))\,\tilde\rho(dr) < \infty$
and $\int(1 \wedge x^2)\,k(x)\,dx/x < \infty$.

(The Gaussian part.) The following are equivalent: (a$'$) $a(\lambda t) = \lambda^2a(t)$ for all
$\lambda, t > 0$; (b$'$) $\mathcal{B}(\lambda A) = \lambda^2\mathcal{B}(A)$ for all $\lambda > 0$
and Borel $A$; (c$'$) $\beta(d\theta) = b\,\theta\,d\theta$ with $b \ge 0$. Then
$a(t) = \tfrac b2t^2$.

(Both parts.) $\Psi_{\lambda t}(\omega) = \Psi_t(\lambda\omega)$ for all $\lambda, t > 0$ iff
(a) and (a$'$) hold.
\item The circle average of a knot-atom cascade is a Cin ray:
\begin{equation}\label{eq:line-e5}
  \int_{[1,q)}\Psi^{(\theta,r)}_t(\omega)\,\eta(d\theta) = \frac1\ell\,\Cin(\omega\,r\,t).
\end{equation}
\item For a compact $A \subset (0,\infty)$ and $j \in \ZZ$ let
$\widehat\sigma_j(A) := \int_{[1,q) \times A}e^{2\pi ij\log_q\theta}\,\sigma(d\theta,dr)$. The
\emph{modulation spectrum} of the family is the set of $j \ne 0$ with $\widehat\sigma_j \ne 0$. It
is empty iff $\sigma = \eta \otimes \tilde\rho$ for some $\tilde\rho$.
\end{enumerate}
% CHANGED (wave 3c, 2026-11-14): lem:line-exponent-data added to \uses. The direction "covariance
% gives (a) and (a')" of "both parts" reads its clause (1), the uniqueness of the data, and so reaches
% A39(ii)-(iii) through that lemma, as thm:surface-measure does; the node cites nothing of its own.
% CHANGED (2026-11-14, a strengthening): an unused hypothesis dropped from the Lean; the print keeps
% the setting of the node, which its other clauses read.
% CHANGED (wave 3c, 2026-11-14): the jump part's equivalences and clause (3) are machine-checked, so
% statement and proof are marked proved; the status line says what is checked.
% CHANGED (fidelity review, 2026-11-15, R64): the status line says that ``both parts'' is stated for
% the pair, the family-level reading composed with thm:surface-measure.
% CHANGED (fidelity review, 2026-11-15, R26/R50/R69): the status line lists the declarations' hypotheses beside the node's (LINKAGE rule 6).
\statusT\quad\emph{(Cites nothing of its own; the direction of ``both parts'' in (1) from the
covariance to (a) and (a$'$) reads the uniqueness of the data, Lemma~\ref{lem:line-exponent-data}(1).
Every clause is machine-checked: the equivalence of (a), (b) and (c) of the jump part, the sentence
beginning ``Then'', the integrability equivalences, the Gaussian part, ``both parts'', (2) and (3);
the measure $\tilde\rho$ of (c) and (3) is read as $s$-finite, which the $\sigma$-finite one built
in the proof is. \textbf{Hypotheses} (the node's, each with what its declarations require): $q > 1$:
same, absent from \texttt{strata\_cylinder\_1\_both}; a family with the pair $(\Lambda, \mathcal{B})$:
weaker, the declarations require the pair alone (\texttt{IsSurfacePair}), reading $\varpi_t$ as
$\Lambda((0,t] \times \cdot\,)$ and $a(t)$ as $\mathcal{B}((0,t])$, and
\texttt{strata\_cylinder\_1\_integrability} and \texttt{strata\_cylinder\_2} require no pair; the
measure $\tilde\rho$ of (c), in the sentence beginning ``Then'' and in the integrability:
stronger, harmless, \texttt{strata\_cylinder\_1\_then} and \texttt{\_integrability} require
$\tilde\rho$ $s$-finite and null on $(-\infty, 0]$, which the $\tilde\rho$ of (c) is, being finite
on compact subsets of $(0,\infty)$ by (3); (b$'$) Borel $A$: same; (2) $r, t > 0$: same; (3) $A$
compact in $(0,\infty)$: same. ``Both parts'' is stated in decomposed form, the reader composing it:
\texttt{strata\_cylinder\_1\_both} states it for the exponent (\ref{eq:line-e2}) of the pair, and
\texttt{surface\_measure\_exponent} identifies that exponent with the family's $\Psi_t$.)}
\end{proposition}

% CHANGED (wave 3c, 2026-11-14): the proof follows the machine-checked route. (b)<=>(c) works on the
% cylinder through the bijection of thm:knot-atom-cascades(2): a dilation J_lambda becomes the
% rotation of the circle by lambda, a rotation-invariant finite measure on the circle is a multiple
% of eta by Fourier uniqueness, and the product is identified from the compact slices, in place of
% Lambda' and the Haar measure on the line, whose map A -> c_A had to be shown a measure. (3) uses
% the same lemma on the circle. The description of Lambda in the "Then" sentence is the orbit
% mixture of eta (x) rho, in place of Lambda' on the strips [q^n, q^{n+1}). A closing sentence says
% what is machine-checked.
% CHANGED (the theory paper, 2026-11-16): the closing sentence on checking and cited facts shortened; no step changed.
% CHANGED (supplement review, 2026-11-16): the weight w(x) renamed Theta(x) (w is the mask).
\begin{proof}
\leanok
\uses{thm:knot-atom-cascades, thm:surface-measure, lem:line-exponent-data}
(1) \emph{The jump part}, (a)$\Leftrightarrow$(b). Fix $\lambda > 0$. For $t > 0$ and Borel $A$,
$\Lambda(J_\lambda^{-1}((0,t] \times A)) = \Lambda((0,t/\lambda] \times \lambda^{-1}A)
= \varpi_{t/\lambda}(\lambda^{-1}A) = (T_\lambda\varpi_{t/\lambda})(A)$. So
$\varpi_{\lambda s} = T_\lambda\varpi_s$ for all $s > 0$ iff $(J_\lambda)_*\Lambda$ and $\Lambda$
agree on the sets $(0,t] \times A$ with $A \subseteq [\varepsilon,\infty)$ Borel, $\varepsilon > 0$.
These sets form a $\pi$-system that generates the Borel sets of $(0,\infty)^2$, both measures are
finite on each of them, and countably many of them cover $(0,\infty)^2$. So this holds iff
$(J_\lambda)_*\Lambda = \Lambda$.

(b)$\Leftrightarrow$(c). By Theorem~\ref{thm:knot-atom-cascades}(2), $\Lambda = \Lambda_\sigma$,
and $\sigma$ satisfies (\ref{eq:line-e3}). For $\lambda > 0$ let $R_\lambda(\theta, r) := ([\lambda\theta], r)$
be the rotation of the cylinder by $\lambda$. \emph{The cylinder measure of
$(J_\lambda)_*\Lambda_\sigma$ is $(R_\lambda)_*\sigma$.} Indeed, of the points
$(\lambda\theta q^n, \lambda r\theta q^n)$, $n \in \ZZ$, of the dilated orbit of
$(\theta, r) \in \mathcal{C}$, exactly one has its first coordinate in $[1,q)$, the one with
$\lambda\theta q^n = [\lambda\theta]$, and its ratio of coordinates is $r$; so its image under
$(y,v) \mapsto (y, v/y)$ is $R_\lambda(\theta, r)$.

\emph{The case $\sigma = \eta \otimes \tilde\rho$.} Let $\tilde\rho$ be a $\sigma$-finite measure
on $(0,\infty)$ (the $\tilde\rho$ of (c) is, by the bound in the paragraph ``From the slices to the
product''), and $G \ge 0$ Borel on $(0,\infty)$. Since $y \mapsto q^ny$ maps $[1,q)$ onto
$[q^n, q^{n+1})$, these intervals partition $(0,\infty)$, and $dy/y$ is invariant under dilations,
\[
  \int_{[1,q)}\sum_{n \in \ZZ}G(\theta q^n)\,\eta(d\theta) = \int_{(0,\infty)}G(y)\,\frac{dy}{\ell y} .
\]
Applied at fixed $r$ to $G(y) = \mathbf 1_E(y, ry)$, with Tonelli, this shows that
$\Lambda_{\eta \otimes \tilde\rho}$ is the image of $(dy/(\ell y)) \otimes \tilde\rho$ under
$(y,r) \mapsto (y, ry)$. Then $J_\lambda$ composed with this map is the map composed with
$(y,r) \mapsto (\lambda y, r)$, and $dy/(\ell y)$ is invariant under $y \mapsto \lambda y$; so
$\Lambda_{\eta \otimes \tilde\rho}$ is invariant under every $J_\lambda$. This gives
(c)$\Rightarrow$(b), and, as $\Lambda = \Lambda_\sigma$, the description of $\Lambda$ in the
statement. It also shows that $\eta$ is invariant under every rotation: for
$\sigma_0 := \eta \otimes \delta_1$, the cylinder measure of $(J_\lambda)_*\Lambda_{\sigma_0}$ is
both $\sigma_0$ and $(R_\lambda)_*\sigma_0$; the first marginal of $\sigma_0$ is $\eta$, and that of
$(R_\lambda)_*\sigma_0$ is the image of $\eta$ under $\theta \mapsto [\lambda\theta]$.

\emph{Rotation-invariant measures on the circle.} For a measure $\nu$ on $[1,q)$ write
$\widehat\nu_j := \int e^{2\pi ij\log_q\theta}\,\nu(d\theta)$, $j \in \ZZ$. Since
$\log_q[\lambda\theta] = \log_q\lambda + \log_q\theta - n$ with $n$ an integer, the image of $\nu$
under $\theta \mapsto [\lambda\theta]$ has the coefficients
$e^{2\pi ij\log_q\lambda}\,\widehat\nu_j$. If $\nu$ is invariant under every rotation, then for
$j \ne 0$ the choice $\lambda = q^{1/(2j)}$ gives $\widehat\nu_j = -\widehat\nu_j$, so
$\widehat\nu_j = 0$; in particular $\widehat\eta_j = 0$ for $j \ne 0$. \emph{A finite measure $\nu$
on $[1,q)$ with $\widehat\nu_j = 0$ for every $j \ne 0$ is $\nu([1,q))\,\eta$.} The map
$\theta \mapsto \log_q\theta \operatorname{mod} 1$ carries $[1,q)$ bijectively onto $\RR/\ZZ$,
with the measurable inverse $u \mapsto q^u$, $u \in [0,1)$, and it carries $\widehat\nu_j$ to the $j$-th
Fourier coefficient of the image of $\nu$. The images of $\nu$ and of $\nu([1,q))\,\eta$ are finite
measures on $\RR/\ZZ$ with the same Fourier coefficients, by $\widehat\eta_j = 0$ for $j \ne 0$
and $\widehat\eta_0 = 1$. The trigonometric polynomials are dense in the continuous functions on
$\RR/\ZZ$ for the supremum norm, and the integral against a finite measure is continuous in that
norm, so the two images have the same integral against every continuous function; a finite Borel
measure on a metric space is determined by these integrals. So the images agree, and, the map being
injective on $[1,q)$ with a measurable inverse, so do $\nu$ and $\nu([1,q))\,\eta$.

\emph{From the slices to the product.} Let $\tilde\rho(A) := \sigma([1,q) \times A)$. For
$\varepsilon > 0$ the weight of (\ref{eq:line-e3}) is at least $\Theta(\varepsilon) > 0$ on
$[1,q) \times [\varepsilon,\infty)$, with $\Theta(x) := x^2 \wedge (1 + \log_+x)$ nondecreasing, so
$\tilde\rho([\varepsilon,\infty)) < \infty$; and $\tilde\rho$ is null on $(-\infty,0]$, hence
$\sigma$-finite. For a compact $A \subset (0,\infty)$, say $A \subseteq [\varepsilon,\infty)$,
the slice $\sigma_A := \sigma(\cdot \times A)$ is thus a finite measure on $[1,q)$. Suppose that
$\widehat{(\sigma_A)}_j = 0$ for every such $A$ and every $j \ne 0$. Then
$\sigma_A = \tilde\rho(A)\,\eta$, that is $\sigma(E \times A) = \eta(E)\,\tilde\rho(A)$ for Borel
$E$. For fixed $E$ both sides are measures in $A$, null on $(-\infty,0]$, finite on every
$[\varepsilon,\infty)$, and equal on every $[a,b]$ with $a > 0$. Restricted to $[\varepsilon,\infty)$
they are finite measures that agree on the closed intervals, hence equal; letting
$\varepsilon \downarrow 0$, they agree on every Borel $A$. So $\sigma$ and $\eta \otimes \tilde\rho$
agree on the rectangles, and since $\eta$ and $\tilde\rho$ are $\sigma$-finite,
$\sigma = \eta \otimes \tilde\rho$.

(b)$\Rightarrow$(c). Assume (b). Then
$(R_\lambda)_*\sigma = \sigma$ for every $\lambda > 0$, as the cylinder measure of
$(J_\lambda)_*\Lambda = \Lambda$. Since $R_\lambda$ does not move $r$, every slice $\sigma_A$ is
invariant under every rotation, so $\widehat{(\sigma_A)}_j = 0$ for $j \ne 0$, and
$\sigma = \eta \otimes \tilde\rho$.

\emph{The Gaussian part.} (a$'$)$\Leftrightarrow$(b$'$): $a(\lambda t) = \lambda^2a(t)$ is
$\mathcal{B}(\lambda A) = \lambda^2\mathcal{B}(A)$ on the intervals $A = (0,t]$, and two measures
finite on these intervals that agree on them agree on all Borel sets. (a$'$)$\Rightarrow$(c$'$):
$a(t) = a(1)t^2$, so $\mathcal{B}(dy) = 2a(1)\,y\,dy$ and $\beta(d\theta) = b\,\theta\,d\theta$
with $b = 2a(1)$. (c$'$)$\Rightarrow$(b$'$): if $\beta = b\theta\,d\theta$, then
$\mathcal{B} = b\,y\,dy$ by Theorem~\ref{thm:knot-atom-cascades}(2), which is
$\lambda^2$-homogeneous, and $a(t) = \tfrac b2t^2$.

\emph{Both parts.} By Theorem~\ref{thm:surface-measure}, $\Psi_t$ has the data $(a(t), \varpi_t)$.
By Lemma~\ref{lem:line-exponent-data}(1),(2), $\Psi_{\lambda t} = \Psi_t(\lambda\,\cdot)$ iff
$a(\lambda t) = \lambda^2a(t)$ and $\varpi_{\lambda t} = T_\lambda\varpi_t$.

% CHANGED (wave 3c, 2026-11-14): Lambda is read as the image of (dy/(ell y)) (x) rho, from the
% (b)<=>(c) paragraph, in place of Lambda'.
\emph{The profile.} With $\Lambda$ the image of $(dy/(\ell y)) \otimes \tilde\rho$ under
$(y,r) \mapsto (y, ry)$ and the substitution $v = ry$ at fixed $r$, for which $dy/y = dv/v$,
\[
  \varpi_t(A) = \frac1\ell\int\tilde\rho(dr)\int\mathbf 1_{\{y \le t\}}\mathbf 1_A(ry)\,\frac{dy}{y}
  = \int_A\frac1\ell\,\tilde\rho\bigl([v/t, \infty)\bigr)\,\frac{dv}{v}.
\]
By Tonelli, $\int(1 \wedge x^2)k(x)\,dx/x = \ell^{-1}\int g(r)\,\tilde\rho(dr)$ with
$g(r) := \int_0^r(1 \wedge x^2)\,dx/x$, which is $r^2/2$ for $r \le 1$ and $\tfrac12 + \log r$ for
$r > 1$. So $\tfrac12\Theta(r) \le g(r) \le \Theta(r)$ with $\Theta(r) := r^2 \wedge (1 + \log_+r)$. Finally, for
$\theta \in [1,q)$, $\Theta(r) \le \Theta(r\theta) \le \max(q^2, 1 + \ell)\,\Theta(r)$. So (\ref{eq:line-e3}) for
$\eta \otimes \tilde\rho$ is $\int \Theta\,d\tilde\rho < \infty$.
% CHANGED (wave 3c, 2026-11-14): the finiteness of k, which the "Then" sentence asserts and the
% machine-checked proof reads from (\ref{eq:line-e3}), is said.
In particular $k$ is finite: $\sigma$ satisfies (\ref{eq:line-e3}) by Theorem~\ref{thm:knot-atom-cascades}(2),
so $\int \Theta\,d\tilde\rho < \infty$, and $\Theta(x)\,\tilde\rho([x,\infty)) \le \int \Theta\,d\tilde\rho$ for
$x > 0$, since $\Theta$ is nondecreasing and positive on $(0,\infty)$. It is nonincreasing because
$x \mapsto \tilde\rho([x,\infty))$ is.

(2) By Tonelli and the substitution $u = \theta q^n$, which maps $[1,q)$ onto $[q^n,q^{n+1})$ with
$du/u = d\theta/\theta$,
\[
  \frac1\ell\int_{[1,q)}\sum_n\mathbf 1_{\{\theta q^n \le t\}}\bigl(1 - \cos(\omega r\theta q^n)\bigr)
  \frac{d\theta}{\theta}
  = \frac1\ell\int_0^t\bigl(1 - \cos(\omega ru)\bigr)\frac{du}{u} = \frac1\ell\,\Cin(\omega rt),
\]
the last step by $\Cin(z) = \int_0^z(1 - \cos v)\,dv/v$ and $v = \omega ru$ for $\omega r > 0$, and
by evenness otherwise.

(3) For compact $A \subset (0,\infty)$, say $A \subseteq [\varepsilon,\infty)$ with
$\varepsilon > 0$, $\sigma([1,q) \times A) \le \tilde\rho([\varepsilon,\infty)) < \infty$, as in (1).
% CHANGED (wave 3c, 2026-11-14): the two directions of the equivalence are argued through the slices
% and the lemma on the circle of (1), and the passage from the compact A to the product is said.
The coefficient $\widehat\sigma_j(A)$ is the coefficient $\widehat{(\sigma_A)}_j$ of the slice
$\sigma_A = \sigma(\cdot \times A)$, since $\sigma$ is carried by $\mathcal{C}$. If the spectrum is
empty, these vanish for every compact $A \subset (0,\infty)$ and every $j \ne 0$, and the paragraph
``From the slices to the product'' of (1) gives $\sigma = \eta \otimes \tilde\rho$. Conversely, if
$\sigma = \eta \otimes \tilde\rho$, then $\widehat\sigma_j(A) = \widehat\eta_j\,\tilde\rho(A) = 0$
for $j \ne 0$, by the rotation invariance of $\eta$.

Every step is machine-checked, and the cited fact enters only through
Lemma~\ref{lem:line-exponent-data}(1).
\end{proof}

The cone of the continuous stratum is the one of the discrete stratum seen at a coarser resolution.
% CHANGED (2026-11-23, review round 3 batch C, item 15): the cited fact is the node thm:cited-A15.
Its extreme rays are the Gaussian ray and the Cin rays \prog{V-cone v0.1}{Prop.~3.3, pp.~15--16}
(Theorem~\ref{thm:cited-A15}).
The Cin rays are the circle averages of the knot-atom cascades, by (2). The Gaussian ray is the
circle average of the knot Gaussians: the substitution of the proof of (2) gives
$\int\sum_{\theta q^n \le t}(\theta q^n)^2\,\eta(d\theta) = \ell^{-1}\int_0^tu\,du = t^2/(2\ell)$.
A Cin ray spreads its displacement atoms log-uniformly over the scales, and a knot-atom cascade
switches them on at knots. In (1) the measure $\ell^{-1}\tilde\rho = -dk$ is the measure
$\kappa^{\mathrm C}$ of (\ref{eq:generator-form}), so the integrability (\ref{eq:line-e3}) at Haar $\sigma$ is
the one of (\ref{eq:admissible-exponent}).

% draft: Proposition 2.12 (section 2.3), the sentence "The marginal projection: ...", lines
% 195-197 of draft/spatial-hemigroup-pyramids.md.
% CHANGED: split from Proposition 2.12. The cone is defined by T_q varpi >= varpi; the
% representation by q-Cin rays is stated with its mixing measure, and the limit q -> 1 with an
% explicit bound (the draft's "Riemann sum of mesh ell" is not monotone and needed one).
% CHANGED: the identification of K_q with the exponents of the symmetric semi-selfdecomposable laws
% of span q, an apposition without proof, made clause (3) with a proof that spends A43 (blind
% review 2, flag 1); the node is graded [A].
\begin{proposition}[the $q$-Cin rays]
\label{prop:strata-cylinder-coordinates-q-cin}
\lean{Pyramids.qcin_1_cone, Pyramids.qcin_1_repr, Pyramids.qcin_1_range, Pyramids.qcin_1_rays,
Pyramids.qcin_2, Pyramids.qcin_2_limit, Pyramids.qcin_3}
\leanok
\uses{lem:line-exponent-data, def:self-decomposable, def:folded-data, thm:cited-A43}
Let $\mathcal{K}_q$ be the set of $\psi \in \mathrm{ND}_s$ whose data $(a, \varpi)$ have
$T_q\varpi \ge \varpi$. For $\tau > 0$ let
\[
  E^{(q)}_\tau(\omega) := \sum_{i \ge 0}\bigl(1 - \cos(\omega\tau q^{-i})\bigr),
\]
the exponent with data $(0, \sum_{i \ge 0}\delta_{\tau q^{-i}})$, the \emph{$q$-Cin ray} of
cutoff $\tau$.
\begin{enumerate}
\item $\mathcal{K}_q$ is a convex cone containing $\omega^2$ and every $E^{(q)}_\tau$. Every
$\psi \in \mathcal{K}_q$ is uniquely
\[
  \psi(\omega) = a\,\omega^2 + \int_{(0,\infty)}E^{(q)}_\tau(\omega)\,\kappa^{(q)}(d\tau)
\]
with $a \ge 0$ and $\kappa^{(q)}$ a positive measure on $(0,\infty)$. A pair $(a, \kappa)$ with
$a \ge 0$ and $\kappa$ a positive measure on $(0,\infty)$ represents in this way an element of
$\mathcal{K}_q$ iff
% CHANGED (2026-11-23, review round 3 batch C, supp.R3): the range of the representation stated.
\[
  \int_{(0,\infty)}\bigl(\tau^2 \wedge (1 + \log_+\tau)\bigr)\,\kappa(d\tau) < \infty .
\]
The extreme rays of $\mathcal{K}_q$ are those of $\omega^2$ and of the $E^{(q)}_\tau$, $\tau > 0$.
\item For $\tau, \omega \in \RR$ with $\tau > 0$,
\[
  \bigl|\ell\,E^{(q)}_\tau(\omega) - \Cin(\tau\omega)\bigr| \le \frac{q\,\ell^2}{2(q - 1)}\,|\tau\omega|,
\]
so $\ell E^{(q)}_\tau \to \Cin(\tau\,\cdot)$ as $q \downarrow 1$, uniformly on bounded sets of
$\tau\omega$.
\item $\mathcal{K}_q$ is the set of the exponents $-\log\widehat\mu$ of the symmetric
semi-selfdecomposable laws $\mu$ of span $q$ (Definition~\ref{def:self-decomposable}).
\end{enumerate}
% CHANGED (fidelity review, 2026-11-15, R26/R50/R69): the status line lists the declarations' hypotheses beside the node's (LINKAGE rule 6).
\statusA\quad\ledger{A43}\quad\emph{(\textbf{Assignment.} \ledger{A43}
[\sscite{sato1999levy}, Prop.~15.5, p.~93] carries one step of (3): a semi-selfdecomposable law is
infinitely divisible. The uniqueness of the residual, the other half of the entry, is not used.
Everything else in (3), and clause (1), are proved here from Lemma~\ref{lem:line-exponent-data};
clause (2) is proved from the definition of $\Cin$ and rests on no cited fact. The criterion
$T_q\nu \ge \nu$ of Sato's Thm.~15.8 is not cited: (3) proves its symmetric case on the line.
\textbf{Hypotheses} (the node's, each with what its declarations require): $q > 1$: same, and in
\texttt{qcin\_2\_limit} the limit $q \downarrow 1$ through $q > 1$; $\psi \in \mathcal{K}_q$: same
(\texttt{IsInKq}); the range, \texttt{qcin\_1\_range}: $a \ge 0$ and $\kappa$ a positive measure
on $(0,\infty)$: same, $\kappa$ read as a measure on $\RR$ with $\kappa((-\infty,0]) = 0$, the
weight being \texttt{cylWeight}; $\tau > 0$: same in (1), weaker in \texttt{qcin\_2} and \texttt{qcin\_2\_limit},
which hold for every real $\tau$; (2) bounded sets of $\tau\omega$: same, $|\tau\omega| \le R$ for
each $R$; (3) $\mu$ a symmetric semi-selfdecomposable law of span $q$: same
(\texttt{IsProbabilityMeasure}, \texttt{IsSymmetric},
\texttt{IsSemiSelfDecomposableWithSpan}).)}
% Lean (wave 3a, 2026-11-14, a strengthening): qcin_2 and qcin_2_limit hold for every real tau, the
% series depending on tau*omega only. The print keeps tau > 0, the range on which E^{(q)}_tau is
% defined.
\end{proposition}

% CHANGED (wave 3a, 2026-11-14): the proof follows the machine-checked route. In (1) the measure
% kappa^{(q)} is built as varpi - T_{q^{-1}} varpi; the uniqueness first shows that the spread of a
% representing measure is a Levy measure, then reads Lemma line-exponent-data(1); the extreme rays
% are proved here rather than referred to Theorem knot-atom-cascades(3). In (2) the comparison is made
% in the variable v, with the same constant. In (3) the sentence on A43 says where the entry is read.
% CHANGED (supplement review, 2026-11-16): the averaging bound and the point-mass argument are
% cited from lem:averaging-bound and lem:null-conull-point-mass; the scalar w renamed c' (w is the
% mask).
\begin{proof}
\leanok
\uses{lem:line-exponent-data, thm:cited-A43, lem:averaging-bound, lem:null-conull-point-mass}
For a measure $\kappa$ on $(0,\infty)$ write $S\kappa := \sum_{i \ge 0}T_{q^{-i}}\kappa$. Then
$\int f\,d(S\kappa) = \int\sum_{i \ge 0}f(\tau q^{-i})\,\kappa(d\tau)$ for measurable $f \ge 0$,
$S\delta_\tau = \sum_{i \ge 0}\delta_{\tau q^{-i}}$, and by Tonelli
\[
  \int E^{(q)}_\tau(\omega)\,\kappa(d\tau) = \int(1 - \cos\omega v)\,(S\kappa)(dv). \qquad (*)
\]
Splitting off the term $i = 0$ gives $S\kappa = \kappa + T_{q^{-1}}S\kappa$. Reindexing gives
$T_qS\kappa = T_q\kappa + S\kappa$, so $S\kappa \le T_qS\kappa$. Two measures that vanish on
$(-\infty,0]$ and agree on every Borel set contained in some $(\varepsilon,\infty)$,
$\varepsilon > 0$, are equal, by continuity from below.

(1) \emph{The cone.} If $\psi, \psi'$ have the data $(a,\varpi)$ and $(a',\varpi')$, then
$\psi + \psi'$ has the data $(a + a', \varpi + \varpi')$, and $c\psi$ with $c \ge 0$ has the data
$(ca, c\varpi)$. The condition $T_q\varpi \ge \varpi$ is preserved by both, since $T_q$ is linear.
The function $\omega^2$ has the data $(1, 0)$. By $(*)$, $E^{(q)}_\tau$ has the data
$(0, S\delta_\tau)$: the \Levy{} integral is $\sum_m(1 \wedge \tau^2q^{-2m}) < \infty$, and
$S\delta_\tau \le T_qS\delta_\tau$.

\emph{Existence.} Let $\psi \in \mathcal{K}_q$ have the data $(a, \varpi)$. The dilation $T_q$ is an
order-preserving bijection with inverse $T_{q^{-1}}$, so $T_{q^{-1}}\varpi \le \varpi$. By
Lemma~\ref{lem:line-exponent-data}(2) and the construction in the proof of its clause (3), there is
a measure $\kappa^{(q)}$ with $\varpi = \kappa^{(q)} + T_{q^{-1}}\varpi$. It vanishes on
% CHANGED (2026-11-23, review round 3, supp.R10(1)): wording.
$(-\infty,0]$. Applying this identity to $q^iA$ for $0 \le i < n$ gives, for every Borel set $A$,
\[
% CHANGED (2026-11-23, review round 3, supp.R10(1)): wording.
  \varpi(A) = \sum_{0 \le i < n}\kappa^{(q)}(q^iA) + \varpi(q^nA).
\]
If $A \subseteq (\varepsilon,\infty)$, then $q^nA \subseteq (\varepsilon q^n,\infty)$, and
$\varpi((\varepsilon q^n,\infty)) \to 0$ as $n \to \infty$, since $\varpi$ is finite on
$(\varepsilon,\infty)$. So $\varpi(A) = \sum_{i \ge 0}\kappa^{(q)}(q^iA) = (S\kappa^{(q)})(A)$, and
$\varpi = S\kappa^{(q)}$. By $(*)$,
$\psi(\omega) = a\omega^2 + \int(1 - \cos\omega v)\,\varpi(dv)
= a\omega^2 + \int E^{(q)}_\tau(\omega)\,\kappa^{(q)}(d\tau)$.

\emph{Uniqueness.} Let $(a', \kappa')$ also represent $\psi$. The measure $\kappa'$ is
$\sigma$-finite: $E^{(q)}_\tau(1) > 0$ for $\tau > 0$, since the terms with $\tau q^{-i} < 2\pi$ are
positive, and $\int E^{(q)}_\tau(1)\,\kappa'(d\tau) \le \psi(1) < \infty$. Next, $S\kappa'$ is a
\Levy{} measure. It is a countable sum of finite measures, as a sum of dilates of the
$\sigma$-finite $\kappa'$, and by $(*)$, $\int(1 - \cos\omega v)\,(S\kappa')(dv) \le \psi(\omega)$
for every $\omega$. So $\int(1 \wedge v^2)\,(S\kappa')(dv) < \infty$ by
Lemma~\ref{lem:averaging-bound}(3). Now $\psi$ has the data $(a', S\kappa')$,
and Lemma~\ref{lem:line-exponent-data}(1) gives $a' = a$ and $S\kappa' = \varpi$. Finally
$\kappa' + T_{q^{-1}}\varpi = \varpi = \kappa^{(q)} + T_{q^{-1}}\varpi$, and $T_{q^{-1}}\varpi$ is
finite on every $(\varepsilon,\infty)$. So $\kappa'$ and $\kappa^{(q)}$ agree on the Borel sets
contained in some $(\varepsilon,\infty)$, and $\kappa' = \kappa^{(q)}$.

% CHANGED (2026-11-23, review round 3 batch C, supp.R3): the range of the representation, proved
% through the bracket of prop:ladder-knot-kernel, as the Lean does (qcin_1_range).
\emph{The range.} Let $a \ge 0$ and let $\kappa$ be a positive measure on $(0,\infty)$. If
$(a, \kappa)$ represents an element $\psi$ of $\mathcal{K}_q$, then $S\kappa$ is a \Levy{} measure, by
the argument of the uniqueness with $\kappa$ in place of $\kappa'$. Conversely, if $S\kappa$ is a
\Levy{} measure, $(a, \kappa)$ represents the element of $\mathcal{K}_q$ with the data
$(a, S\kappa)$, since $S\kappa \le T_qS\kappa$ and by $(*)$. So $(a, \kappa)$ represents an element
of $\mathcal{K}_q$ iff $S\kappa$ is a \Levy{} measure. Now
$\int(1 \wedge v^2)\,(S\kappa)(dv) = \int s(\tau)\,\kappa(d\tau)$ with
$s(\tau) := \sum_{i \ge 0}(1 \wedge \tau^2q^{-2i}) = (1 \wedge \tau^2) + b(\tau)$, where
$b(\tau) := \sum_{m \ge 1}(1 \wedge \tau^2q^{-2m})$. Write $w(\tau) := \tau^2 \wedge (1 + \log_+\tau)$,
which is $\tau^2$ for $\tau \le 1$ and $1 + \log\tau$ for $\tau > 1$, since
$1 + \log\tau \le \tau \le \tau^2$ there. For $\tau > 0$, let $N$ be the integer part of
$\log_+\tau/\ell$. The terms of $b(\tau)$ with $m \le N$ are $1$, and the others are at most
$\tau^2q^{-2m} = (\tau q^{-N-1})^2q^{-2(m - N - 1)}$ with $(\tau q^{-N-1})^2 \le 1 \wedge \tau^2$, a
geometric series of ratio $q^{-2}$. So
\[
  N \le b(\tau) \le \frac{\log_+\tau}{\ell} + \frac{q^2}{q^2 - 1}\,(1 \wedge \tau^2) .
\]
For $\tau \le 1$ this gives $w(\tau) = \tau^2 \le s(\tau) \le (1 + q^2/(q^2 - 1))\,w(\tau)$. For
$\tau > 1$, $\log\tau < \ell(N + 1) \le \ell s(\tau)$ and $1 \le s(\tau)$, so
$w(\tau) \le (1 + \ell)\,s(\tau)$, and $s(\tau) \le 1 + q^2/(q^2 - 1) + \log\tau/\ell
\le (1 + q^2/(q^2 - 1) + \ell^{-1})\,w(\tau)$. Hence
\[
  \frac{w(\tau)}{1 + \ell} \le s(\tau) \le \Bigl(1 + \frac{q^2}{q^2 - 1} + \frac1\ell\Bigr)\,w(\tau)
  \qquad (\tau > 0),
\]
and $S\kappa$ is a \Levy{} measure iff $\int w\,d\kappa < \infty$.

\emph{The extreme rays.} By the uniqueness, the representation of a sum is the sum of the
representations, and that of $c\psi$, $c \ge 0$, is $(ca, c\kappa)$. By the range, every pair
$(a,\kappa)$ with $a \ge 0$ for which $S\kappa$ is a \Levy{} measure represents an element of
$\mathcal{K}_q$.

The rays of $\omega^2$ and of $E^{(q)}_\tau$ are extreme. Let $c > 0$ and $y + z = c\omega^2$, or
$y + z = cE^{(q)}_\tau$, with $y, z \in \mathcal{K}_q$. The representations of $y$ and $z$ add up to
$(c, 0)$, or to $(0, c\delta_\tau)$. So the representing measure of $y$ is $0$, or is at most
$c\delta_\tau$ and hence equal to $c'\delta_\tau$ with $c' \le c$. In both cases $y$ is a nonnegative
multiple of the ray.

Conversely, let $\psi \in \mathcal{K}_q$ span an extreme ray, with the representation $(a, \kappa)$.
For a Borel set $B$, the pairs $(0, \kappa|_B)$ and $(a, \kappa|_{B^c})$ represent elements $y, z$
of $\mathcal{K}_q$ with $y + z = \psi$. So $y = c\psi$ for some $c \ge 0$, and $\kappa|_B = c\kappa$.
Evaluated on $B^c$ and on $B$, this gives $\kappa(B) = 0$ or $\kappa(B^c) = 0$. So
$\kappa = \kappa(\{\tau\})\,\delta_\tau$ for some $\tau$, by
Lemma~\ref{lem:null-conull-point-mass}. Splitting off $a\omega^2$ in the same way shows that
$a = 0$ or $\kappa = 0$. If $\kappa = 0$, then $\psi = a\omega^2$ with $a > 0$. Otherwise $a = 0$ and
$\kappa = c'\delta_\tau$ with $c' := \kappa(\{\tau\}) > 0$. Here $\tau > 0$, since $\kappa$ vanishes on
$(-\infty,0]$, and $c' < \infty$, since $c'E^{(q)}_\tau(1) = \psi(1)$ and $E^{(q)}_\tau(1) > 0$. So
$\psi = c'E^{(q)}_\tau$.

(2) Let $x := \tau\omega > 0$, $b_i := xq^{-i}$ and $\ell = \log q$. For $b > 0$,
$\ell(1 - \cos b) = \int_{b/q}^b(1 - \cos b)\,dv/v$ and
$\Cin(b) - \Cin(b/q) = \int_{b/q}^b(1 - \cos v)\,dv/v$. Since $|\cos v - \cos b| \le b - v$ for
$v \le b$,
\[
  \bigl|\ell(1 - \cos b) - (\Cin(b) - \Cin(b/q))\bigr| \le \int_{b/q}^b\frac{b - v}{v}\,dv
  = b\Bigl(\ell - 1 + \frac1q\Bigr).
\]
The differences $\Cin(b_i) - \Cin(b_{i+1})$ are nonnegative, and their partial sums
$\Cin(x) - \Cin(xq^{-n})$ tend to $\Cin(x)$, since $\Cin$ is continuous with $\Cin(0) = 0$. Summing
over $i \ge 0$, with $\sum_ib_i = xq/(q - 1)$,
\[
  \bigl|\ell E^{(q)}_\tau(\omega) - \Cin(\tau\omega)\bigr|
  \le \Bigl(\ell - 1 + \frac1q\Bigr)\frac{xq}{q - 1} \le \frac{q\,\ell^2}{2(q - 1)}\,x .
\]
The last step is $e^{-\ell} \le 1 - \ell + \tfrac12\ell^2$, which follows from
$e^\ell \ge 1 + \ell + \tfrac12\ell^2 + \tfrac16\ell^3$, since
$(1 - \ell + \tfrac12\ell^2)(1 + \ell + \tfrac12\ell^2 + \tfrac16\ell^3)
= 1 + \tfrac16\ell^3 + \tfrac1{12}\ell^4 + \tfrac1{12}\ell^5 \ge 1$. If $\tau\omega < 0$, replace
$\omega$ by $-\omega$: both sides are even in $\omega$. If $\tau\omega = 0$, both sides vanish. Since
$\ell \le q - 1$, the bound is at most $\tfrac12q\ell\,|\tau\omega|$, which tends to $0$ as
$q \downarrow 1$, uniformly for bounded $\tau\omega$.

(3) Let $\psi \in \mathrm{ND}_s$ have the data $(a, \varpi)$. The function
$\psi(q^{-1}\,\cdot)$ has the data $(q^{-2}a, T_{q^{-1}}\varpi)$, by
Lemma~\ref{lem:line-exponent-data}(2), and $q^{-2}a \le a$. The dilation $T_q$ is an order-preserving
bijection of the measures on $(0,\infty)$ with inverse $T_{q^{-1}}$, so $T_{q^{-1}}\varpi \le \varpi$
iff $\varpi \le T_q\varpi$. By Lemma~\ref{lem:line-exponent-data}(3), therefore,
\[
  \psi \in \mathcal{K}_q \iff \psi - \psi(q^{-1}\,\cdot) \in \mathrm{ND}_s .
\]
Let $\psi \in \mathcal{K}_q$. By Lemma~\ref{lem:line-exponent-data}(1), $\psi = -\log\widehat\mu$
for a symmetric infinitely divisible law $\mu$, and $\psi - \psi(q^{-1}\,\cdot) = -\log\widehat\varrho$
for a symmetric infinitely divisible law $\varrho$. Then
$\widehat\mu = \widehat\mu(q^{-1}\,\cdot)\,\widehat\varrho$, so $\mu$ is semi-selfdecomposable of
span $q$ with the residual $\varrho$. Conversely, let $\mu$ be a symmetric semi-selfdecomposable law
of span $q$ with residual $\varrho$. By Theorem~\ref{thm:cited-A43} (read on $\RR^1$, to which $x \mapsto xe_1$ carries $\mu$ and its
residual), $\mu$ is infinitely divisible. By
Lemma~\ref{lem:line-exponent-data}(1), $\widehat\mu = e^{-\psi}$ with $\psi \in \mathrm{ND}_s$, and
in particular $\widehat\mu > 0$. So
$\widehat\varrho = \widehat\mu/\widehat\mu(q^{-1}\,\cdot) = e^{-(\psi - \psi(q^{-1}\,\cdot))}$ is real
and even. The transform determines the law, so $\varrho$ equals its reflection and is symmetric. It
is infinitely divisible by definition, so $\psi - \psi(q^{-1}\,\cdot) = -\log\widehat\varrho$ lies
in $\mathrm{ND}_s$ by Lemma~\ref{lem:line-exponent-data}(1). By the display, $\psi \in \mathcal{K}_q$.
\end{proof}

The \Levy{} measure of the $q$-Cin kernel of cutoff $\tau$, the law with exponent
$\ell E^{(q)}_\tau$, is carried by $(0,\tau]$, as that of the Cin kernel of step $\tau$ is.

% draft: Lemma 2.13 (section 2.3), lines 201-203 of draft/spatial-hemigroup-pyramids.md.
% CHANGED: clause (1), the series for Psi_t at every phase, added; it is what the proof uses and
% what the paragraph after the draft's lemma states for t = 1.
\begin{lemma}[the $q$-step field determines the family]
\label{lem:step-field-determines-family}
\lean{Pyramids.step_field_1, Pyramids.step_field_2, Pyramids.step_field_2_knots}
\leanok
\uses{thm:annulus-parametrization}
Let a $q^{\ZZ}$-covariant family in the canonical discrete gauge have the exponents $\Psi_t$, and
let $g_{t,qt} := \Psi_{qt} - \Psi_t$ be the exponent of the step $\mu_{t,qt}$.
\begin{enumerate}
\item For every $t > 0$, $\Psi_t = \sum_{m \ge 1}g_{t,qt}(q^{-m}\,\cdot)$.
\item Two such families with the same $g_{t,qt}$ for all $t \in [1,q)$ coincide. Two such families
with the same octave step $\mu_{1,q}$ have the same knot kernels.
\end{enumerate}
% CHANGED (fidelity review, 2026-11-15, R26/R50/R69): the status line lists the declarations' hypotheses beside the node's (LINKAGE rule 6).
\statusT\quad\emph{(\textbf{Hypotheses} (the node's, each with what its declarations require): $q > 1$:
same; a $q^{\ZZ}$-covariant family in the canonical discrete gauge: same (\texttt{IsQFamily}), for
each of the two families of (2); (1) $t > 0$: same; (2) the same $g_{t,qt}$ for $t \in [1,q)$:
same; the same octave step $\mu_{1,q}$: same.)}
\end{lemma}

% CHANGED (wave 2f, 2026-11-14): the proof follows the machine-checked route. It says why the
% series converges (its terms are nonnegative), how the kernels on [1, q) determine the family
% (Lemma discrete-covariance-levy read in R^1), and how the octave step gives g_{1,q}.
\begin{proof}
\leanok
\uses{lem:discrete-covariance-levy, lem:family-structure, lem:family-dilations}
(1) Covariance at the pair $(0, t)$ gives $\Psi_{qt}(\omega) = \Psi_t(q\omega)$, so
$g_{t,qt}(q^{-m}\omega) = \Psi_t(q^{1-m}\omega) - \Psi_t(q^{-m}\omega)$ and
$\sum_{m=1}^{M}g_{t,qt}(q^{-m}\omega) = \Psi_t(\omega) - \Psi_t(q^{-M}\omega)$. Each term is
nonnegative: $g_{t,qt}$ is the exponent of the step $\mu_{t,qt}$, and the transforms of the kernels
are positive and nonincreasing in $t$ (Lemma~\ref{lem:family-dilations}, through
Lemma~\ref{lem:family-structure}). The last term tends to $\Psi_t(0) = 0$, since
$\Psi_t = -\log\widehat\mu_t$ is continuous. So the series of nonnegative terms converges, to
$\Psi_t(\omega)$.
(2) By (1) the field determines $\Psi_t$ for $t \in [1,q)$, hence the kernels $\mu_t$ there, a
symmetric law being determined by its transform $e^{-\Psi_t}$. A $q$-covariant family in the
canonical gauge is determined by its kernels on $[1, q)$: the images of the two families in $\RR^1$
under $x \mapsto xe_1$ are such families, and Lemma~\ref{lem:discrete-covariance-levy} at $d = 1$
says that they coincide; the image map is injective on laws. For the second sentence, the
octave step $\mu_{1,q}$ has the transform $e^{-g_{1,q}}$, since $\widehat\mu_q = \widehat\mu_{1,q}\,
\widehat\mu_1$. So (1) at $t = 1$ gives $\Psi_1$, hence $\mu_1$, and the knot kernels are its dilates,
$\mu_{q^n} = T_{q^n}\mu_1$, by covariance iterated.
\end{proof}

So two pyramids with the same ladder can differ only between their knots. The ladder fixes the
knot kernels (Proposition~\ref{prop:ladder-knot-kernel}(1)), and through
Theorem~\ref{thm:surface-measure} it fixes the mass of $\Lambda$ on the slab
$(1,q] \times (0,\infty)$ against each set of displacements. It does not fix how that mass is spread
over the scales of the octave. This is the interpolation problem of a pyramid. A construction that
stores only the levels, as the classical ones do, does not determine it.

% draft: the paragraph after Lemma 2.13 (section 2.3), line 205 of
% draft/spatial-hemigroup-pyramids.md, which defines the ladder, its step measure and the
% modulation orbit in prose.
% CHANGED: the three facts the definition asserted (the ladder determines Psi_1; the two
% descriptions of the orbit agree; Lambda((1,q] x .) is the step measure) move to
% lem:knot-exponent-ladder, with proofs that need no family (blind review 2, flags 8 and 21). The
% step measure is written mu, the name the later nodes use (flag 20).
\begin{definition}[knot exponent, step measure, modulation orbit]
\label{def:modulation-orbit}
\lean{Pyramids.knotStep, Pyramids.IsKnotExponent, Pyramids.IsKnotLadder, Pyramids.IsStepMeasure,
Pyramids.HasNoGaussianStep, Pyramids.modulationOrbit}\leanok
\uses{def:knots-and-ladder}
A \emph{knot exponent} is a function $\Psi_1 \in \mathrm{ND}_s$ whose \emph{step}
$\phi := \Psi_1(q\,\cdot) - \Psi_1$ lies in $\mathrm{ND}_s$ and is not identically zero. Its
\emph{ladder} is the sequence of the laws with the exponents $\phi(q^n\,\cdot)$, $n \in \ZZ$, as in
Definition~\ref{def:knots-and-ladder}. The \emph{step measure} $\mu$ of the ladder is the folded
\Levy{} measure of $\phi$. (Kernels always carry indices, as $\mu_1$ and $\mu_{s,t}$, and the letter
$\mu$ alone is the step measure.) The ladder has \emph{no Gaussian part} if the Gaussian
coefficient of $\phi$ is $0$. The \emph{modulation orbit} of the ladder of $\Psi_1$ is the set of
$q^{\ZZ}$-covariant families in the canonical discrete gauge whose knot kernel $\mu_1$ has the
exponent $\Psi_1$.
% CHANGED (fidelity review, 2026-11-15, R8): the status line says which range hypotheses of the
% definition are not fields of its declarations, and that every consumer carries them.
\statusT\quad\emph{(A definition. The range $q > 1$ is not a field of \texttt{modulationOrbit} or
of \texttt{IsQFamily}, the $q^{\ZZ}$-covariant families in the canonical gauge; every statement that
uses them carries $1 < q$.)}
\end{definition}

% draft: the same paragraph (line 205), its three assertions. New node at the transcription.
\begin{lemma}[the ladder of a knot exponent]
\label{lem:knot-exponent-ladder}
\lean{Pyramids.knot_exponent_ladder_1, Pyramids.knot_exponent_ladder_1_determines,
Pyramids.knot_exponent_ladder_2, Pyramids.knot_exponent_ladder_3}
\leanok
\uses{def:modulation-orbit, lem:step-field-determines-family, thm:surface-measure}
Let $\Psi_1$ be a knot exponent with step $\phi$ and step measure $\mu$.
\begin{enumerate}
\item $\Psi_1 = \sum_{m \ge 1}\phi(q^{-m}\,\cdot)$. In particular the ladder determines $\Psi_1$,
and the modulation orbit depends on the ladder only.
\item A $q^{\ZZ}$-covariant family in the canonical discrete gauge lies in the modulation orbit iff
its octave step $\mu_{1,q}$ has the exponent $\phi$.
\item Every member of the modulation orbit has $\Lambda((1,q] \times \cdot\,) = \mu$.
\end{enumerate}
% CHANGED (2026-11-14, a strengthening): an unused hypothesis dropped from the Lean of (3); the
% print keeps the setting of the node, which (1) and (2) read.
% CHANGED (fidelity review, 2026-11-15, R26/R50/R69): the status line lists the declarations' hypotheses beside the node's (LINKAGE rule 6).
\statusT\quad\emph{(Cites nothing of its own. \textbf{Hypotheses} (the node's, each with what its
declarations require): $q > 1$: same; $\Psi_1$ a knot exponent: same (\texttt{IsKnotExponent}) for
\texttt{knot\_exponent\_ladder\_1}, \texttt{\_1\_determines} and \texttt{\_2}, weaker in
\texttt{knot\_exponent\_ladder\_3}, which does not require it: membership of the orbit makes $\Psi_1$ the
exponent of the knot kernel, and the step measure is given; $\mu$ the step measure: same
(\texttt{IsStepMeasure}); (1) ``the ladder determines $\Psi_1$'': two knot exponents with one ladder,
same; (2) a $q^{\ZZ}$-covariant family in the canonical discrete gauge: same (\texttt{IsQFamily});
(3) a member of the orbit, with its pair $(\Lambda, \mathcal{B})$: same, the pair taken as a
hypothesis (\texttt{IsFamilyPair}), which exists by Theorem~\ref{thm:surface-measure}.)}
\end{lemma}

% CHANGED (wave 3a, 2026-11-14): the proof follows the machine-checked route. In (3) the pair of the
% octave step is written down and added to the data at 1, in place of lem:line-exponent-data(3).
% CHANGED (the theory paper, 2026-11-16): the closing sentence on checking and cited facts shortened; no step changed.
\begin{proof}
\leanok
\uses{lem:step-field-determines-family, thm:surface-measure, lem:line-exponent-data,
lem:family-structure, lem:discrete-covariance-levy}
(1) $\sum_{m=1}^M\phi(q^{-m}\omega) = \sum_{m=1}^M\bigl(\Psi_1(q^{1-m}\omega) - \Psi_1(q^{-m}\omega)\bigr)
= \Psi_1(\omega) - \Psi_1(q^{-M}\omega)$. The last term tends to $\Psi_1(0) = 0$, because a function
in $\mathrm{ND}_s$ is continuous, by dominated convergence. The ladder's law at $n = 0$ has the
exponent $\phi$, which the law determines, and $\phi$ determines $\Psi_1$ by the series.

(2) For a family, the exponent of $\mu_{1,q}$ is $\Psi_q - \Psi_1 = \Psi_1(q\,\cdot) - \Psi_1$, by
Lemmas~\ref{lem:family-structure} and~\ref{lem:discrete-covariance-levy}. If the family's
$\Psi_1$ is the knot exponent, this is $\phi$. Conversely, if $\Psi_q - \Psi_1 = \phi$, then the
family's $\Psi_1$ is $\sum_{m \ge 1}\phi(q^{-m}\,\cdot)$ by
Lemma~\ref{lem:step-field-determines-family}(1) at $t = 1$, which is the knot exponent by (1).

(3) By (\ref{eq:line-e1}) the kernels $\mu_1$ and $\mu_q$ have the folded data $(a(1), \varpi_1)$ and
$(a(q), \varpi_q)$, with $\varpi_t = \Lambda((0,t] \times \cdot\,)$ and $a(t) = \mathcal{B}((0,t])$.
So $a(1) \le a(q)$, and $\varpi_q = \varpi_1 + \Lambda((1,q] \times \cdot\,)$, the strip
$(0,q]$ being the disjoint union of $(0,1]$ and $(1,q]$. The measure $\Lambda((1,q] \times \cdot\,)$
is carried by $(0,\infty)$ and is at most $\varpi_q$, so it is a folded \Levy{} measure, and
$(a(q) - a(1), \Lambda((1,q] \times \cdot\,))$ is a pair whose sum with $(a(1), \varpi_1)$ is
$(a(q), \varpi_q)$. Its exponent is therefore $\Psi_q - \Psi_1$, which is $\phi$ by (2), the
family's $\Psi_1$ being the knot exponent. By the uniqueness of the data
(Lemma~\ref{lem:line-exponent-data}(1)) its measure is the step measure $\mu$.

Every step is machine-checked, and the cited facts enter only through
Lemmas~\ref{lem:step-field-determines-family}, \ref{lem:family-structure} and
\ref{lem:line-exponent-data}.
\end{proof}

The Gaussian coefficients of a member are tied to the ladder in the same way.

% draft: none. New at the supplement review (2026-11-16): the argument "a(q) - a(1) is the
% Gaussian coefficient of the step, and a = 0 by monotonicity", proved inside the proofs of
% thm:phase-entropy-selection(3)-(4) and thm:phase-entropy-modulated-selection, stated once and
% cited there.
\begin{lemma}[the Gaussian part of a member]
\label{lem:member-gaussian-part}
\lean{Pyramids.exists_step_gauss, Pyramids.gauss_eq_zero_of_noGaussianStep}
\leanok
\uses{def:modulation-orbit, thm:surface-measure}
Let $\Psi_1$ be a knot exponent with step $\phi$, and let a member of the modulation orbit of its
ladder have the pair $(\Lambda, \mathcal{B})$ and $a(t) := \mathcal{B}((0,t])$.
\begin{enumerate}
\item $a(q) - a(1)$ is the Gaussian coefficient of $\phi$.
\item If the ladder has no Gaussian part, then $a(t) = 0$ for every $t > 0$.
\end{enumerate}
\statusT\quad\emph{(Cites nothing of its own; the uniqueness of the data is
Lemma~\ref{lem:line-exponent-data}(1). \textbf{Hypotheses} (the node's, each with what its
declarations require): $q > 1$: same; $\Psi_1$ a knot exponent: absent, the step being read as the
exponent of the member's octave step; a member with its pair: same, membership of the orbit and
\texttt{IsFamilyPair}; (2) no Gaussian part: same (\texttt{HasNoGaussianStep}). Clause (1) is
stated in decomposed form, the reader composing it: \texttt{exists\_step\_gauss} gives data of
$\phi$ whose Gaussian coefficient added to $a(1)$ is $a(q)$, and these are the data of $\phi$ by
\texttt{line\_exponent\_data\_1\_unique}.)}
\end{lemma}

\begin{proof}
\leanok
\uses{thm:surface-measure, lem:line-exponent-data, lem:family-structure,
lem:discrete-covariance-line}
(1) The exponent $\Psi_q - \Psi_1$ of the octave step $\mu_{1,q}$ lies in $\mathrm{ND}_s$
(Lemma~\ref{lem:family-structure}), and it is $\Psi_1(q\,\cdot) - \Psi_1 = \phi$
(Lemma~\ref{lem:discrete-covariance-line}). By (\ref{eq:line-e1}) at $t = 1$ and $t = q$ and
Lemma~\ref{lem:line-exponent-data}(3), it has the data $(a(q) - a(1), \Lambda((1,q] \times \cdot\,))$.
By the uniqueness of the data (Lemma~\ref{lem:line-exponent-data}(1)), $a(q) - a(1)$ is the
Gaussian coefficient of $\phi$.

(2) By (1), $a(q) = a(1)$. With $a(q) = q^2a(1)$, which is Theorem~\ref{thm:surface-measure}(ii)
on the interval $(0,1]$, this gives $a(1) = 0$, then $a(q^n) = q^{2n}a(1) = 0$ for $n \ge 0$, and
$a(t) = 0$ for every $t > 0$ because $a$ is nondecreasing and nonnegative.

Every step is machine-checked, and the cited facts enter only through the statements the proof
cites.
\end{proof}

\section{Phase entropy: what it selects, and where it selects nothing}
\label{sec:line-phase-entropy}

Among the families with one ladder, which interpolation between the knots is canonical? A natural
principle requires that no phase within an octave be preferred: each new displacement should appear
spread evenly over the octave's scales, in the logarithmic coordinate that the dilation $J$ makes
uniform. Relative entropy against the Haar measure of the circle of phases measures the departure
from that. On a ladder with no Gaussian part, the principle selects the fully covariant member when
the ladder has one, which is close to a restatement of its hypothesis. It selects a modulated
covariant member on a strictly larger class of ladders. It selects nothing on a ladder whose step
measure has a singular part, which includes every grid pyramid. A second functional, the entropy
of the scale at which each displacement appears, selects one member on every ladder with no
Gaussian part, and never the similarity member. For probability measures
$\pi, \eta$ write $\mathrm{KL}(\pi \,\|\, \eta) := \int\log\frac{d\pi}{d\eta}\,d\pi \in [0,\infty]$
if $\pi$ is absolutely continuous with respect to $\eta$, and $+\infty$ otherwise. It vanishes iff
$\pi = \eta$. Both functionals integrate it over a measurable family of probability measures.

% draft: none. New at the supplement review (2026-11-16): the measurability of r -> KL(pi_r || eta),
% proved twice nearly word for word (in the proofs of thm:phase-entropy-selection and
% prop:phase-entropy-ladder-bound(2)), stated once and cited there.
% CHANGED (blueprint round, 2026-11-16, the dispositions' UNSOURCED flag): the Mathlib parenthetical
% is replaced by Kallenberg's Thm. 9.27, p. 200 (kernel densities, T and U Borel; read by the
% librarian from the held copy), with the two-line passage from a conditional density to a density
% with respect to eta. The Lean reads Mathlib's Kernel.rnDeriv; the route is the same fact.
\begin{lemma}[the relative entropy of a measurable family]
\label{lem:kl-measurable}
\lean{Pyramids.measurable_klDiv_kernel}
\leanok
Let $\eta$ be a probability measure on $\RR$, and let $\pi_r$, $r \in \RR$, be probability measures
on $\RR$, measurable in $r$. Then $r \mapsto \mathrm{KL}(\pi_r \,\|\, \eta)$ is measurable, into
$[0,\infty]$.
\statusT\quad\emph{(Cites nothing. \textbf{Hypotheses} (the node's, each with what its declaration
requires): $\eta$ a probability measure: weaker, a finite measure; $\pi_r$ probability measures,
measurable in $r$: same.)}
\end{lemma}

\begin{proof}
\leanok
Where $\pi_r$ is absolutely continuous with respect to $\eta$,
$\mathrm{KL}(\pi_r \,\|\, \eta) = \int f(d\pi_r/d\eta)\,d\eta$ with $f(x) = x\log x + 1 - x \ge 0$.
By [\sscite{kallenberg2021foundations}, Thm.~9.27, p.~200], applied to the probability kernel
$r \mapsto \tfrac12(\pi_r \otimes \delta_1 + \eta \otimes \delta_0)$ from $\RR$ to
$\RR \times \{0, 1\}$, the density $g(r, \theta)$ of $\pi_r$ with respect to $\pi_r + \eta$ has a
version with values in $[0, 1]$, jointly measurable in $(r, \theta)$. As
$\eta = (1 - g)(\pi_r + \eta)$ puts no mass on $\{g = 1\}$, the function $p := g/(1 - g)$ on
$\{g < 1\}$, $p := 0$ on $\{g = 1\}$, is a density of the absolutely continuous part of $\pi_r$
with respect to $\eta$, jointly measurable in $(r, \theta)$. Then $\pi_r$ is absolutely continuous with respect
to $\eta$ iff $\int p(r, \theta)\,\eta(d\theta) = 1$, the total mass of $\pi_r$, and
$r \mapsto \int p(r, \theta)\,\eta(d\theta)$ is measurable by Tonelli's theorem. So the set of these
$r$ is measurable, and off it the value is $+\infty$. On it $r \mapsto \int f(p(r, \theta))\,
\eta(d\theta)$ is measurable, by Tonelli's theorem again.

Every step is machine-checked, with no cited fact. The formal proof takes the jointly measurable
density from Mathlib's Radon--Nikodym derivative of kernels, the fact read here from Kallenberg.
\end{proof}

% CHANGED (skeleton review A1, m2): the claims of def:phase-entropy split out as a lemma, so that
% they have a typed statement and the statements quantifying over every disintegration are seen
% not to be vacuous.
% CHANGED (fidelity review, 2026-11-15, R63): the tag adds Pyramids.exists_phaseDisintegration and
% Pyramids.phaseDisintegration_ae_eq, the existence and the uniqueness for every measure on the
% cylinder with (\ref{eq:line-e3}), which the status line names; the tagged phase_disintegration_exists takes
% the pair of a family. The statement is unchanged.
% CHANGED (supplement review, 2026-11-16): the lemma moved before def:phase-entropy, which cites it
% and no longer asserts its claims; the statement defines the r-marginal itself, which it read from
% the definition, and its uses drop def:phase-entropy. No claim changed.
\begin{lemma}[the phase disintegration exists]
\label{lem:phase-disintegration-exists}
\lean{Pyramids.phase_disintegration_exists, Pyramids.exists_phaseDisintegration,
Pyramids.phaseDisintegration_ae_eq}
\leanok
\uses{thm:knot-atom-cascades, thm:surface-measure}
Let a $q^{\ZZ}$-covariant family in the canonical discrete gauge have the cylinder measure $\sigma$
(Theorem~\ref{thm:knot-atom-cascades}), and let $\tilde\rho(A) := \sigma([1,q) \times A)$ be its
\emph{$r$-marginal}. Then $\tilde\rho([r_0,\infty)) < \infty$ for every $r_0 > 0$.
There are probability measures $\pi_r$ on $[1,q)$, measurable in $r$, with
$\sigma(d\theta, dr) = \pi_r(d\theta)\,\tilde\rho(dr)$, and any two such families agree for
$\tilde\rho$-almost every $r$. The last two sentences hold for every measure $\sigma$ on $\RR^2$
carried by $\mathcal{C}$ that satisfies (\ref{eq:line-e3}), with $\tilde\rho$ defined in the same way.
% CHANGED (2026-11-14, a strengthening): an unused hypothesis dropped from the Lean; the print
% keeps the setting of the node.
% CHANGED (theory paper register pass, 2026-11-16, a strengthening): the last sentence prints the
% general form that exists_phaseDisintegration and phaseDisintegration_ae_eq already state, so that
% prop:appearance-entropy cites a statement and not the inside of this proof. The Lean is unchanged.
% CHANGED (fidelity review, 2026-11-15, R26/R50/R69): the status line lists the declarations' hypotheses beside the node's (LINKAGE rule 6).
\statusT\quad\emph{(Cites nothing of its own. \textbf{Hypotheses} (the node's, each with what its
declarations require): $q > 1$: same, absent from \texttt{phaseDisintegration\_ae\_eq}; a
$q^{\ZZ}$-covariant family with the cylinder measure $\sigma$: weaker,
\texttt{phase\_disintegration\_exists} requires the pair $(\Lambda, \mathcal{B})$ with (i)--(iv) and
that the kernels $\mu_{0,t}$ have its data (\texttt{IsFamilyPair}), not that the laws form a family,
and its proof reads only condition (i); \texttt{exists\_phaseDisintegration} and
\texttt{phaseDisintegration\_ae\_eq} require only a measure $\sigma$ carried by the cylinder with
(\ref{eq:line-e3}).)}
\end{lemma}

% CHANGED (wave 3c, 2026-11-14): the proof follows the machine-checked route. The weight is w(r)
% itself, in place of a piecewise constant g; the exceptional fibres are replaced by the point mass
% at 1, in place of eta; the uniqueness is the uniqueness of a disintegration of a finite measure,
% applied after weighting by w, in place of the pi-system of rational intervals.
% CHANGED (the theory paper, 2026-11-16): the closing sentence on checking and cited facts shortened; no step changed.
\begin{proof}
\leanok
\uses{thm:knot-atom-cascades, thm:surface-measure}
Write $\Theta(x) := x^2 \wedge (1 + \log_+x)$, which is nondecreasing on $[0,\infty)$ and positive on
$(0,\infty)$. The family's $\Lambda$ satisfies Theorem~\ref{thm:surface-measure}(i), so $\sigma$,
which is carried by $\mathcal{C}$, satisfies (\ref{eq:line-e3}) by Theorem~\ref{thm:knot-atom-cascades}(2).
The rest of the proof reads of $\sigma$ only these two properties, so it also proves the last
sentence. For $(\theta, r) \in \mathcal{C}$ one has $r\theta \ge r > 0$, so $\Theta(r) \le \Theta(r\theta)$, and
\[
  \int \Theta\,d\tilde\rho = \int_{\mathcal{C}}\Theta(r)\,\sigma(d\theta,dr)
  \le \int_{\mathcal{C}}\Theta(r\theta)\,\sigma(d\theta,dr) < \infty .
\]
For $r_0 > 0$ this gives $\Theta(r_0)\,\tilde\rho([r_0,\infty)) \le \int_{[r_0,\infty)}\Theta\,d\tilde\rho < \infty$,
and $\Theta(r_0) > 0$. Since $\sigma$ is carried by $\mathcal{C}$, $\tilde\rho$ is null on $(-\infty,0]$,
so $\Theta > 0$ holds $\tilde\rho$-almost everywhere. The measure $\Theta\tilde\rho$ is finite, and
$\tilde\rho = \Theta^{-1}(\Theta\tilde\rho)$ is $\sigma$-finite.

Let $\tau$ be the image of $\sigma$ under $(\theta, r) \mapsto (r, \theta)$, whose first marginal
is $\tilde\rho$. The measure $\tau_\Theta(dr, d\theta) := \Theta(r)\,\tau(dr, d\theta)$ is finite, on the
product of two standard Borel spaces, with first marginal $\Theta\tilde\rho$. By the disintegration
theorem for finite measures, $\tau_\Theta(dr, d\theta) = \Theta(r)\tilde\rho(dr)\,\kappa_r(d\theta)$ for a
measurable family of probability measures $\kappa_r$ on $\RR$. Dividing by $\Theta(r)$, which is positive
for $\tilde\rho$-almost every $r$ and finite, gives $\tau(dr, d\theta) = \tilde\rho(dr)\,\kappa_r(d\theta)$.
Testing this on $\RR \times (\RR \setminus [1,q))$, where $\tau$ is null, gives
$\int\kappa_r(\RR \setminus [1,q))\,\tilde\rho(dr) = 0$, so $\kappa_r$ is carried by $[1,q)$ for
$\tilde\rho$-almost every $r$. Put $\pi_r := \kappa_r$ for those $r$ and $\pi_r := \delta_1$ for
the others, a measurable choice since the exceptional set is measurable. The identity is unchanged on
a $\tilde\rho$-null set, so $\sigma(d\theta, dr) = \pi_r(d\theta)\,\tilde\rho(dr)$, each $\pi_r$ a
probability measure on $[1,q)$.

For uniqueness let $\pi$ and $\pi'$ both disintegrate $\sigma$. Then
$\tilde\rho(dr)\,\pi_r(d\theta) = \tau = \tilde\rho(dr)\,\pi'_r(d\theta)$, and multiplying by $\Theta(r)$,
$\Theta(r)\tilde\rho(dr)\,\pi_r(d\theta) = \Theta(r)\tilde\rho(dr)\,\pi'_r(d\theta)$. The first marginal
$\Theta\tilde\rho$ is finite and $\RR$ is countably generated, so by the uniqueness of the
disintegration of a finite measure $\pi_r = \pi'_r$ for $\Theta\tilde\rho$-almost every $r$. As $\Theta > 0$
$\tilde\rho$-almost everywhere, $\tilde\rho$ is absolutely continuous with respect to $\Theta\tilde\rho$,
and $\pi_r = \pi'_r$ for $\tilde\rho$-almost every $r$.

Every step is machine-checked, and the proof cites no fact of its own.
\end{proof}

% draft: Definition 2.14 (section 2.4), line 209 of draft/spatial-hemigroup-pyramids.md.
% CHANGED: the relative entropy written KL (the draft's D collides with D_c and the diagonal
% dilation; review flag); the definition is given for every family, and the existence of the
% disintegration is stated; the sentence "E is intrinsic, since r, theta and Haar are defined by
% the q-action alone" moves to the prose after it, corrected: E is computed in a fixed canonical
% discrete gauge.
\begin{definition}[phase entropy]
\label{def:phase-entropy}
\lean{Pyramids.haarPhase, Pyramids.rMarginal, Pyramids.IsPhaseDisintegration, Pyramids.phaseEntropy}\leanok
\uses{thm:knot-atom-cascades, prop:strata-cylinder-coordinates, lem:phase-disintegration-exists,
lem:kl-measurable}
% CHANGED (supplement review, 2026-11-16): the definition cites lem:phase-disintegration-exists,
% now before it, for the r-marginal and the disintegration, and no longer asserts their finiteness
% and uniqueness; the integrand is measurable by lem:kl-measurable.
Let a $q^{\ZZ}$-covariant family in the canonical discrete gauge have the cylinder measure $\sigma$
and the $r$-marginal $\tilde\rho$, and let $\sigma(d\theta, dr) = \pi_r(d\theta)\,\tilde\rho(dr)$ be
a disintegration of $\sigma$ as in Lemma~\ref{lem:phase-disintegration-exists}. The
\emph{phase entropy} of the family is
\[
  \mathfrak{E} := -\int\mathrm{KL}(\pi_r \,\|\, \eta)\,\tilde\rho(dr) \in [-\infty, 0],
\]
with $\eta$ the Haar measure of Proposition~\ref{prop:strata-cylinder-coordinates}.
\statusT
\end{definition}

The integrand is measurable by Lemma~\ref{lem:kl-measurable}, and two disintegrations agree for
$\tilde\rho$-almost every $r$ by Lemma~\ref{lem:phase-disintegration-exists}, so $\mathfrak{E}$ does
not depend on the one chosen.

% CHANGED (pass 6, 2026-11-22, the author's review): the argument that a change of the scale unit
% leaves the phase entropy unchanged, made here in prose, is the next proposition; the prose cites it.
The coordinates $r$ and $\theta$ and the measure $\eta$ are defined by the action of $J$. So a
change of the scale unit does not change the phase entropy.

% draft: none. New at pass 6 (2026-11-22, the author's review of the supplement): the claim of the
% paragraph after def:phase-entropy, stated and proved
% (Formalization/Pyramids/LineAppendix/PhaseEntropyScaleUnit.lean). The family relabelled by
% t -> ct is mu_{s/c,t/c}; the statement names its pair, which the prose left implicit.
\begin{proposition}[the scale unit]
\label{prop:phase-entropy-scale-unit}
\lean{Pyramids.phase_entropy_scale_unit}
\leanok
\uses{def:phase-entropy, thm:surface-measure, thm:knot-atom-cascades,
lem:phase-disintegration-exists}
Let a $q^{\ZZ}$-covariant family $(\mu_{s,t})$ in the canonical discrete gauge have the pair
$(\Lambda, \mathcal{B})$ and the cylinder measure $\sigma$, and let $c > 0$. The family relabelled
by $t \mapsto ct$, with the kernels $\mu'_{s,t} := \mu_{s/c,\,t/c}$, is again $q^{\ZZ}$-covariant
in the canonical discrete gauge. Its pair is the image of $\Lambda$ under $(y, v) \mapsto (cy, v)$
and the image of $\mathcal{B}$ under $y \mapsto cy$. Its cylinder measure is the image of $\sigma$
under $(\theta, r) \mapsto ([c\theta], r/c)$, where $[x] \in [1,q)$ is the class of $x > 0$ modulo
$q^{\ZZ}$. The two families have the same phase entropy.
\statusT\quad\emph{(Cites nothing. Every clause is machine-checked.
\textbf{Hypotheses} (the node's, each with what its declaration requires): $q > 1$: same; a
$q^{\ZZ}$-covariant family in the canonical discrete gauge: same (\texttt{IsQFamily}); its pair:
same (\texttt{IsFamilyPair}); $c > 0$: same. The phase entropies are compared at every
disintegration of each cylinder measure.)}
\end{proposition}

\begin{proof}
\leanok
\uses{thm:surface-measure, thm:knot-atom-cascades, lem:phase-disintegration-exists}
Write $T(y, v) := (cy, v)$, $\Lambda' := T_*\Lambda$, and $\mathcal{B}'$ for the image of
$\mathcal{B}$ under $y \mapsto cy$. The laws $\mu'_{s,t}$ form a family because the $\mu_{s,t}$ do
at $s/c \le t/c$, and they are covariant with $S_qt = qt$ because $q(s/c) = (qs)/c$. The preimage
of $(s,t] \times A$ under $T$ is $(s/c, t/c] \times A$. So
$\Lambda'((0,t] \times \cdot) = \Lambda((0,t/c] \times \cdot)$ and
$\mathcal{B}'((0,t]) = \mathcal{B}((0,t/c])$, which are the data of $\mu_{0,t/c} = \mu'_{0,t}$.
The pair $(\Lambda', \mathcal{B}')$ satisfies (i)--(iv) of Theorem~\ref{thm:surface-measure}:
$T$ commutes with $J$, $\mathcal{B}'(qA) = q^2\mathcal{B}'(A)$ because the dilations by $q$ and
by $c$ commute, and
the slabs, lines and intervals of $\Lambda'$ and $\mathcal{B}'$ are those of $\Lambda$ and
$\mathcal{B}$ at the scales divided by $c$. So $(\Lambda', \mathcal{B}')$ is the pair of $\mu'$.

By Theorem~\ref{thm:knot-atom-cascades}(2), $\Lambda$ is the mixture over $\sigma$ of the orbits
of $J$, and $T$ maps the orbit of $(\theta, r) \in \mathcal{C}$ to the points
$(c\theta q^n, r\theta q^n)$, $n \in \ZZ$. Exactly one of them has its first coordinate in
$[1,q)$, at $n = -\lfloor\log_q c\theta\rfloor$. There the first coordinate is $[c\theta]$ and the
ratio of the coordinates is $r/c$. So $\Lambda'$ is the mixture of the orbits over the image
$\sigma'$ of $\sigma$ under $(\theta, r) \mapsto ([c\theta], r/c)$, a measure carried by
$\mathcal{C}$, and $\sigma'$ is the cylinder measure of $\Lambda'$
(Theorem~\ref{thm:knot-atom-cascades}(2)).

Let $\sigma(d\theta, dr) = \pi_r(d\theta)\,\tilde\rho(dr)$, and let $R(\theta) := [c\theta]$ on
$[1,q)$. It is a bijection of $[1,q)$ with the inverse $\theta \mapsto [\theta/c]$, both Borel.
Put $\pi'_r := R_*\pi_{cr}$. For Borel $E$,
\[
  \sigma'(E) = \int\pi_r(\{\theta : ([c\theta], r/c) \in E\})\,\tilde\rho(dr)
  = \int\pi'_{r/c}(\{\theta : (\theta, r/c) \in E\})\,\tilde\rho(dr).
\]
At $E = [1,q) \times A$ this shows that the $r$-marginal of $\sigma'$ is the image $\tilde\rho'$
of $\tilde\rho$ under $r \mapsto r/c$, and then that $\pi'$ disintegrates $\sigma'$. In the
coordinate $\log_q\theta$ the map $R$ is a rotation of the circle $\RR/\ZZ$ and $\eta$ is its
Lebesgue measure, so $R_*\eta = \eta$. A Borel bijection with a Borel inverse carries the density
of $\pi_{cr}$ against $\eta$ to the density of $R_*\pi_{cr}$ against $R_*\eta$, so
$\mathrm{KL}(\pi'_r \,\|\, \eta) = \mathrm{KL}(\pi_{cr} \,\|\, \eta)$. Substituting $r \mapsto r/c$,
\[
  \int\mathrm{KL}(\pi'_r \,\|\, \eta)\,\tilde\rho'(dr)
  = \int\mathrm{KL}(\pi'_{r/c} \,\|\, \eta)\,\tilde\rho(dr)
  = \int\mathrm{KL}(\pi_r \,\|\, \eta)\,\tilde\rho(dr).
\]
Any other disintegration of $\sigma'$ agrees with $\pi'$ for $\tilde\rho'$-almost every $r$
(Lemma~\ref{lem:phase-disintegration-exists}, for the family $\mu'$), so the two phase entropies
are equal.

Every step is machine-checked, and the proof cites no fact.
\end{proof}

A relabelling in $\mathcal{G}_q$ that is not a dilation moves $y$ nonlinearly and changes $\sigma$
in general, so $\mathfrak{E}$ is computed in the canonical discrete gauge in which the family is
written.

% draft: Theorem 2.15 (section 2.4), lines 211-213 of draft/spatial-hemigroup-pyramids.md.
% CHANGED: "admits a fully covariant member" made precise as "has a member with
% Psi_t = Psi_1(t .)" (by Theorem 2.0 every fully covariant member gives one, after a relabelling in
% G_q that fixes the knot kernel; thm:annulus-parametrization-phase(3)).
% CHANGED: narrowed to ladders with no Gaussian part, and clause (4) added: the phase entropy does
% not see the Gaussian part, so with a Gaussian step the maximizer is not unique (the draft's
% uniqueness claim is false then).
% CHANGED: the reading "in Fourier terms ... empty modulation spectrum" and "the sharp-knot
% couplings sit at the opposite extreme" move to the prose after the theorem (review flag,
% interpretation inside a statement); the sharp-knot couplings are not families.
% CHANGED (proof): the uniqueness of the similarity member is the uniqueness of the knot kernel
% (Lemma 2.13), which replaces the draft's deconvolution argument here; the deconvolution lemma is
% lem:haar-smearing-injective, used by thm:phase-entropy-modulated-selection.
\begin{theorem}[selection where a covariant member exists]
\label{thm:phase-entropy-selection}
\lean{Pyramids.phase_entropy_selection_1, Pyramids.phase_entropy_selection_2,
Pyramids.phase_entropy_selection_3, Pyramids.phase_entropy_selection_4}
\leanok
% CHANGED (wave 3d, 2026-11-14): lem:line-exponent-data added to \uses. Clause (3) and the
% members of (4) read its clause (1), the uniqueness of the data, and so reach A39(ii)-(iii)
% through that lemma, as prop:strata-cylinder-coordinates does; the node cites nothing of its own.
\uses{def:phase-entropy, def:modulation-orbit, lem:line-exponent-data}
Let $\Psi_1$ be a knot exponent whose ladder has a modulation orbit containing the \emph{similarity member}, the family
with $\Psi_t(\omega) = \Psi_1(t\omega)$ for all $t > 0$.
\begin{enumerate}
\item The similarity member is the only member with $\Psi_{\lambda t} = \Psi_t(\lambda\,\cdot)$ for
all $\lambda, t > 0$.
\item Every member has $\mathfrak{E} \le 0$, with equality iff $\pi_r = \eta$ for
$\tilde\rho$-almost every $r$.
\item If the ladder has no Gaussian part, the similarity member is the only member with
$\mathfrak{E} = 0$, so it is the unique maximizer of the phase entropy on the modulation orbit.
% CHANGED (fourth blind pass, 2026-11-15): "whatever their Gaussian coefficients $a(t)$" reads
% "with no condition on their Gaussian coefficients $a(t)$ beyond membership of the orbit", which
% is what the proof shows; the coefficients of a member are constrained by its being a member.
\item If the ladder has a Gaussian part, the members with $\mathfrak{E} = 0$ are those with
$\varpi_t = T_t\varpi_1$ for all $t > 0$, with no condition on their Gaussian coefficients $a(t)$
beyond membership of the orbit, and there are infinitely many.
\end{enumerate}
% CHANGED (2026-11-14, a strengthening): unused hypotheses dropped from the Lean; the print keeps
% the setting of the node.
% CHANGED (wave 3d, 2026-11-14): every clause is machine-checked, so statement and proof are
% marked proved; the status line says what is checked.
% CHANGED (fidelity review, 2026-11-15, R26/R50/R69): the status line lists the declarations' hypotheses beside the node's (LINKAGE rule 6).
\statusT\quad\emph{(Cites nothing of its own; (3) and the members of (4) read the uniqueness of the
data, Lemma~\ref{lem:line-exponent-data}(1). Every clause is machine-checked.
\textbf{Hypotheses} (the node's, each with what its declarations require): $q > 1$: same; $\Psi_1$ a
knot exponent: absent from all four declarations, the proof never reading it; the similarity member
in the orbit: same for \texttt{phase\_entropy\_selection\_1}, \texttt{\_3} and \texttt{\_4} (a member
$\mu_0$ with $\Psi_t = \Psi_1(t\,\cdot)$ for $t > 0$), absent from \texttt{\_2}; (2) a member with
its disintegration: weaker, \texttt{phase\_entropy\_selection\_2} holds for the cylinder measure of
every measure $\Lambda$ and every disintegration of it, not only for the pair of a family; (3) no
Gaussian part: same (\texttt{HasNoGaussianStep}); (3), (4) $\mathfrak{E}$ of a member: same, read
at every pair (\texttt{IsFamilyPair}) and every disintegration of it; (4) a Gaussian part: same, the
negation of \texttt{HasNoGaussianStep}, the step lying in $\mathrm{ND}_s$ as the exponent of the
octave step of the similarity member.)}
\end{theorem}

% CHANGED (wave 3d, 2026-11-14): the proof follows the machine-checked route. An opening paragraph
% proves that the phase entropy vanishes iff sigma = eta (x) rho~, which the print used without the
% uniqueness of the disintegration, and the measurability of r -> KL(pi_r || eta) that (2) needs.
% (1) names the determination of a family by its exponents. (3) derives a(q) - a(1) = 0 from the
% uniqueness of the data. (4) builds the infinitely many members as mixtures of the similarity
% member's B with a second Gaussian measure, in place of the Stieltjes measures of the a_s.
% CHANGED (fourth blind pass, 2026-11-15): the measurability of r -> KL(pi_r || eta) names the
% measurable Radon-Nikodym theorem and derives the measurability of the set where pi_r << eta; in
% (3) the data of Psi_q - Psi_1 cite lem:line-exponent-data(3).
% CHANGED (the theory paper, 2026-11-16): the closing sentence on checking and cited facts shortened; no step changed.
% CHANGED (supplement review, 2026-11-16): the measurability of r -> KL(pi_r || eta) is cited from
% lem:kl-measurable, and a(q) - a(1) and a = 0 from lem:member-gaussian-part; the mixing
% coefficients of (4) renamed alpha_i (w_i are the pyramid weights). No step changed.
\begin{proof}
\leanok
\uses{prop:strata-cylinder-coordinates, lem:knot-exponent-ladder, thm:surface-measure,
lem:line-exponent-data, lem:phase-disintegration-exists, lem:step-field-determines-family,
thm:knot-atom-cascades, lem:kl-measurable, lem:member-gaussian-part}
Two facts are used throughout. First, for a measurable family of probability measures $\pi_r$ the
map $r \mapsto \mathrm{KL}(\pi_r \,\|\, \eta)$ is measurable (Lemma~\ref{lem:kl-measurable}).
Second, for a family with the
cylinder measure $\sigma$ and a disintegration $\pi$, $\mathfrak{E} = 0$ iff
$\sigma = \eta \otimes \tilde\rho'$ for some measure $\tilde\rho'$ on $(0,\infty)$. Indeed, by (2)
below, $\mathfrak{E} = 0$ iff $\pi_r = \eta$ for $\tilde\rho$-almost every $r$, and then the identity
$\sigma(E) = \int\pi_r(E_r)\,\tilde\rho(dr)$ is the product formula for $\eta \otimes \tilde\rho$.
Conversely, if $\sigma = \eta \otimes \tilde\rho'$, then $\tilde\rho' = \tilde\rho$, since $\eta$ is a
probability measure on $[1,q)$; the constant family $\eta$ disintegrates $\sigma$, and two
disintegrations agree for $\tilde\rho$-almost every $r$
(Lemma~\ref{lem:phase-disintegration-exists}). By
Proposition~\ref{prop:strata-cylinder-coordinates}(1), the jump part, $\mathfrak{E} = 0$ is therefore
equivalent to (a) there: $\varpi_{\lambda t} = T_\lambda\varpi_t$ for all $\lambda, t > 0$.

(1) The similarity member has $\Psi_{\lambda t}(\omega) = \Psi_1(\lambda t\omega) = \Psi_t(\lambda\omega)$.
A member with $\Psi_{\lambda t} = \Psi_t(\lambda\,\cdot)$ for all $\lambda, t > 0$ has
$\Psi_t = \Psi_{t\cdot1} = \Psi_1(t\,\cdot)$, and $\Psi_1$ is the knot exponent, the same for every
member. So it has the exponents of the similarity member, and a family is determined by its
exponents on $[1,q)$ (Lemma~\ref{lem:step-field-determines-family}).

(2) $\mathfrak{E} \le 0$ since the relative entropy is nonnegative. The integrand is measurable and
nonnegative, so $\mathfrak{E} = 0$ iff $\mathrm{KL}(\pi_r \,\|\, \eta) = 0$ for $\tilde\rho$-almost
every $r$, and for two probability measures $\mathrm{KL}(\pi_r \,\|\, \eta) = 0$ iff $\pi_r = \eta$.

(3) Let $(\Lambda, \mathcal{B})$ be the pair of a member and $a(t) = \mathcal{B}((0,t])$. The
ladder has no Gaussian part, so $a \equiv 0$ by Lemma~\ref{lem:member-gaussian-part}(2). So (a$'$) of
Proposition~\ref{prop:strata-cylinder-coordinates}(1) holds, and $\mathfrak{E} = 0$ iff (a), iff (a)
and (a$'$), iff $\Psi_{\lambda t} = \Psi_t(\lambda\,\cdot)$ for all $\lambda, t > 0$ (``both parts''
there), iff the member is the similarity member, by (1). The similarity member has a pair
(Theorem~\ref{thm:surface-measure}) and a disintegration
(Lemma~\ref{lem:phase-disintegration-exists}), and its $\mathfrak{E}$ is $0$: by ``both parts'' it
satisfies (a), and this direction does not use the Gaussian part.

(4) The members with $\mathfrak{E} = 0$ are those with (a). Taking $t = 1$ in (a) gives
$\varpi_\lambda = T_\lambda\varpi_1$ for all $\lambda > 0$; conversely this gives
$\varpi_{\lambda t} = T_{\lambda t}\varpi_1 = T_\lambda T_t\varpi_1 = T_\lambda\varpi_t$. Nothing is
required of $a(t)$. For the infinitely many members, let $(\Lambda_0, \mathcal{B}_0)$ be the pair of the
similarity member, which exists by hypothesis, and $a_0(t) = \mathcal{B}_0((0,t])$. By ``both
parts'' it satisfies (a) and (a$'$), so $a_0(t) = t^2a_0(1)$; and $a_0(1) > 0$, since otherwise
$a_0(q) - a_0(1) = 0$ would be the Gaussian coefficient of $\phi$
(Lemma~\ref{lem:member-gaussian-part}(1)). Put
$c := (q + 1)a_0(1)$ and let $\mathcal{B}_{\mathrm{alt}}$ be the measure with (ii) whose restriction to
$[1,q)$ is $c$ times Lebesgue measure (Theorem~\ref{thm:knot-atom-cascades}(2)). It has no atoms, and
$\mathcal{B}_{\mathrm{alt}}((1,q]) = c(q - 1) = (q^2 - 1)a_0(1)$, so
$\mathcal{B}_{\mathrm{alt}}((0,q]) = q^2\mathcal{B}_{\mathrm{alt}}((0,1])$ gives
$\mathcal{B}_{\mathrm{alt}}((0,1]) = a_0(1)$. For $i \ge 0$ put $\alpha_i := 1/(i + 1)$ and
$\mathcal{B}_i := \alpha_i\mathcal{B}_0 + (1 - \alpha_i)\mathcal{B}_{\mathrm{alt}}$. Each $\mathcal{B}_i$
satisfies (ii) and has no atom, and $(\Lambda_0, \mathcal{B}_i)$ satisfies (iv) because
$(\Lambda_0, \mathcal{B}_0)$ does and $\alpha_i > 0$. So $(\Lambda_0, \mathcal{B}_i)$ is the pair of a
family $\mu^{(i)}$ (Theorem~\ref{thm:surface-measure}). Since $\mathcal{B}_i((0,1]) = a_0(1)$, its
knot exponent is that of the similarity member, $\Psi_1$, so $\mu^{(i)}$ is a member; and its
$\mathfrak{E}$ is $0$, since $\Lambda_0$ satisfies (a). At $t_0 := (1 + q)/2$ one has
$\mathcal{B}_0((1,t_0]) = (t_0^2 - 1)a_0(1)$ and $\mathcal{B}_{\mathrm{alt}}((1,t_0]) = c(t_0 - 1)$,
whose difference is $-a_0(1)(q - 1)^2/4 \ne 0$. So $\mathcal{B}_i((1,t_0])$ is a strictly monotone
function of $\alpha_i$, and distinct $i$ give distinct $\mathcal{B}_i$. Two families with the same kernels
have the same pair (Theorem~\ref{thm:surface-measure}), so distinct $i$ give distinct members.

Every step is machine-checked, and the cited facts, \ledger{A39} and \ledger{A41}, enter only
through the statements the proof cites.
\end{proof}

So on a ladder with no Gaussian part the principle selects the covariant member, and it is the
hypothesis that does the work. Haar fibres are full covariance
(Proposition~\ref{prop:strata-cylinder-coordinates}(1)), and the relative entropy against Haar is
largest at Haar. In Fourier terms the selected member is the one with empty modulation spectrum
(Proposition~\ref{prop:strata-cylinder-coordinates}(3)). The couplings that switch each new
displacement on entirely at a knot have every $\pi_r$ a point mass, the opposite extreme. They are
not families, since they jump at the knots (Theorem~\ref{thm:knot-atom-cascades}(4)). A ladder has
a similarity member iff its knot exponent is admissible. For a step with a \Levy{} density,
Proposition~\ref{prop:ladder-knot-kernel}(4) says when that is: the series of dilated step profiles
must be nonincreasing. The next results ask what the phase entropy does without that hypothesis.
They read a measure on the cylinder through its image on the displacements, and the first lemma
computes that image for a product with Haar fibres.

% draft: none. New at the supplement review (2026-11-16): the density of the Haar smearing, computed
% inside the proofs of lem:haar-smearing-injective and thm:phase-entropy-modulated-selection, stated
% once and cited there.
\begin{lemma}[the density of a Haar smearing]
\label{lem:haar-smearing-density}
\lean{Pyramids.map_haarPhase_prod_mul_apply, Pyramids.measurable_measure_Ioc_div}
\leanok
\uses{prop:strata-cylinder-coordinates}
Let $\rho$ be a $\sigma$-finite measure on $(0,\infty)$, and let $\nu$ be the image of
$\eta \otimes \rho$ under $(\theta, r) \mapsto r\theta$. Then $v \mapsto \rho((v/q, v])$ is
measurable, and for every Borel $A \subseteq \RR$,
\[
  \nu(A) = \int_{A \cap (0,\infty)}\rho\bigl((v/q, v]\bigr)\,\frac{dv}{\ell v} .
\]
\statusT\quad\emph{(Cites nothing. \textbf{Hypotheses} (the node's, each with what its declarations
require): $q > 1$, the range of $\eta$: weaker, $q > 0$; $\rho$ a $\sigma$-finite measure on
$(0,\infty)$: weaker, a countable sum of finite measures on $\RR$, null on $(-\infty,0]$; $A$
Borel: same.)}
\end{lemma}

\begin{proof}
\leanok
The set $\{(v, r) : v/q < r \le v\}$ is Borel, so $v \mapsto \rho((v/q, v])$, the integral in $r$ of
its indicator, is measurable by Tonelli's theorem. For Borel $A \subseteq \RR$ and fixed $r > 0$,
the substitution $v = r\theta$ carries $d\theta/(\ell\theta)$ on $[1,q)$ to $dv/(\ell v)$ on
$[r,qr)$, and $r \le v < qr$ iff $v/q < r \le v$. By Tonelli's theorem, twice,
\[
  \nu(A) = \int\rho(dr)\int_{[1,q)}\mathbf 1_A(r\theta)\,\frac{d\theta}{\ell\theta}
  = \int\rho(dr)\int_{[r,qr)}\mathbf 1_A(v)\,\frac{dv}{\ell v}
  = \int_{A \cap (0,\infty)}\rho\bigl((v/q, v]\bigr)\,\frac{dv}{\ell v}.
\]

Every step is machine-checked, with no cited fact.
\end{proof}

% draft: none. New at the transcription: the deconvolution lemma the draft's proof of Theorem 2.15
% uses (its "the closing pattern of Lemma 2.13: log-periodic and integrable is zero"), stated and
% proved; it is used by thm:phase-entropy-modulated-selection.
\begin{lemma}[Haar smearing is injective]
\label{lem:haar-smearing-injective}
\lean{Pyramids.haar_smearing_injective}
\leanok
\uses{prop:strata-cylinder-coordinates}
Let $\rho_1, \rho_2$ be positive measures on $(0,\infty)$ with $\rho_i([x,\infty)) < \infty$ for
every $x > 0$. If the images of $\eta \otimes \rho_1$ and $\eta \otimes \rho_2$ under
$(\theta, r) \mapsto r\theta$ coincide, then $\rho_1 = \rho_2$.
% CHANGED (fidelity review, 2026-11-15, R26/R50/R69): the status line lists the declarations' hypotheses beside the node's (LINKAGE rule 6).
\statusT\quad\emph{(Cites nothing. \textbf{Hypotheses} (the node's, each with what its declarations
require): $q > 1$, the range of $\eta$: same; $\rho_1, \rho_2$ positive measures on $(0,\infty)$: same,
measures on $\RR$ null on $(-\infty, 0]$; $\rho_i([x,\infty)) < \infty$ for $x > 0$: same; the
images coincide: same.)}
\end{lemma}

% CHANGED (wave 3c, 2026-11-14): the proof follows the machine-checked route. The octave masses are
% compared through the tails rho_i((x,infinity)), which are their sums by countable additivity, in
% place of the signed difference delta and its decay at infinity; the passage from almost every point
% to every point is made on the tails. The sigma-finiteness that Tonelli's theorem requires is stated.
% CHANGED (the theory paper, 2026-11-16): the closing sentence on checking and cited facts shortened; no step changed.
% CHANGED (supplement review, 2026-11-16): the density of the smeared measure is cited from
% lem:haar-smearing-density.
\begin{proof}
\leanok
\uses{lem:haar-smearing-density}
Each $\rho_i$ is null on $(-\infty,0]$ and finite on every $(1/n,\infty)$, so it is
$\sigma$-finite. Let $\nu_i$ be the image of $\eta \otimes \rho_i$ under
$(\theta, r) \mapsto r\theta$. By Lemma~\ref{lem:haar-smearing-density}, $\nu_i$ has the density
$g_i(v) := \rho_i((v/q,v])$ against $dv/(\ell v)$ on $(0,\infty)$, and $g_i$ is measurable. The measure $dv/(\ell v)$ is $\sigma$-finite, so $\nu_1 = \nu_2$ gives
$g_1 = g_2$ almost everywhere.

Let $T_i(x) := \rho_i((x,\infty))$. For $x > 0$ the interval $(x,\infty)$ is the disjoint union
of the intervals $(q^nx, q^{n+1}x]$, $n \ge 0$, since $q^nx \to \infty$. By countable additivity,
\[
  T_i(x) = \sum_{n \ge 0} g_i(q^{n+1}x).
\]
The dilation $v \mapsto q^{n+1}v$ preserves $dv/(\ell v)$, so for almost every $x > 0$ the equality
$g_1(q^{n+1}x) = g_2(q^{n+1}x)$ holds for every $n \ge 0$, and then $T_1(x) = T_2(x)$.

The tails agree at every $x > 0$. The density of $dv/(\ell v)$ is positive on $(0,\infty)$, so
every interval $(x, x + 1/(n+1))$ has positive measure and holds a point $y_n$ with
$T_1(y_n) = T_2(y_n)$. The sets $(x + 1/(n+1), \infty)$ increase to $(x,\infty)$ and
$(x + 1/(n+1), \infty) \subseteq (y_n,\infty) \subseteq (x,\infty)$, so
$T_i(x) = \sup_n T_i(y_n)$ for $i = 1, 2$, and $T_1(x) = T_2(x)$.

For $\varepsilon > 0$ the restrictions of $\rho_1$ and $\rho_2$ to $(\varepsilon,\infty)$ are
finite, have the same total mass $T_i(\varepsilon)$, and give each $(a,\infty)$ the same mass
$T_i(\max(a,\varepsilon))$. The intervals $(a,\infty)$ form a $\pi$-system that generates the Borel
sets, so the two restrictions are equal. Both $\rho_i$ are carried by $(0,\infty)$, the increasing
union of the $(1/(n+1),\infty)$, so
$\rho_1(B) = \sup_n\rho_1(B \cap (1/(n+1),\infty)) = \rho_2(B)$ for every Borel $B$.

Every step is machine-checked, with no cited fact.
\end{proof}

Fold the displacements onto the circle of phases: for $v > 0$ let $[v] \in [1,q)$ be its class
modulo $q^{\ZZ}$, and for a measure $\nu$ on $(0,\infty)$ let $\nu^\circ$ be its image under
$v \mapsto [v]$. Multiplication by $r$ descends to the rotation $\mathrm{rot}_r[\theta] := [r\theta]$
of the circle, which preserves $\eta$.

% draft: none in the draft's statements; notes/phase-entropy-without-covariant-member.md, Theorem 1
% (hub Q-0194, P-0005), which records/pyramids/PLAN-review-response.md schedules for this section.
% CHANGED (from the note): 0 < |mu| in (1), which the note leaves implicit (a ladder whose step is
% Gaussian has mu = 0); the identity in (1) is stated where E > -infinity, the bound everywhere;
% clause (3), the grid pyramids, made a statement (the note's prose).
\begin{proposition}[the ladder bound; the obstruction]
\label{prop:phase-entropy-ladder-bound}
\lean{Pyramids.ladder_bound_1, Pyramids.ladder_bound_2, Pyramids.ladder_bound_3}
\leanok
\uses{def:phase-entropy, def:modulation-orbit, def:grid-pyramid-line}
Let $\Psi_1$ be a knot exponent with step measure $\mu$, and let a member of the modulation orbit of
its ladder have the disintegration $\sigma(d\theta,dr) = \pi_r(d\theta)\,\tilde\rho(dr)$.
\begin{enumerate}
\item Suppose $0 < |\mu| := \mu((0,\infty)) < \infty$, put $\bar\mu^\circ := \mu^\circ/|\mu|$ and
$P_r := (\mathrm{rot}_r)_*\pi_r$. Then $|\tilde\rho| = |\mu|$, $\mu^\circ = \int P_r\,\tilde\rho(dr)$,
and
\begin{equation}\label{eq:line-e6}
  \mathfrak{E} \le -|\mu|\,\mathrm{KL}(\bar\mu^\circ \,\|\, \eta),
\end{equation}
a bound fixed by the ladder, and, in $[0,\infty]$,
\[
  -\mathfrak{E} = |\mu|\,\mathrm{KL}(\bar\mu^\circ \,\|\, \eta)
  + \int\mathrm{KL}(P_r \,\|\, \bar\mu^\circ)\,\tilde\rho(dr).
\]
When the bound (\ref{eq:line-e6}) is finite, equality holds in it iff $P_r = \bar\mu^\circ$ for
$\tilde\rho$-almost every $r$.
\item If $\mathfrak{E} > -\infty$ at some member, then $\mu$ is absolutely continuous with respect
to Lebesgue measure. So if $\mu$ has a nonzero singular part, $\mathfrak{E} = -\infty$ at every
member of the modulation orbit.
\item In particular, for a mask $w \ne \delta_0$ on $h\ZZ$ and every family carrying the ladder of
its grid pyramid (Definition~\ref{def:grid-pyramid-line}), $\mathfrak{E} = -\infty$.
\end{enumerate}
% CHANGED (wave 3f, 2026-11-15): every clause is machine-checked, so statement and proof are marked
% proved; the status line says what is checked.
% CHANGED (2026-11-15, a strengthening): unused hypotheses dropped from the Lean of two
% declarations, listed at the end of the status line; the print keeps the node's setting.
% CHANGED (2026-11-15, a strengthening, in print): the identity of (1) no longer assumes
% E > -infinity; it holds in [0, infinity], both sides infinite when E = -infinity. Was: "If
% E > -infinity, then -E = ...". The Lean ladder_bound_1 drops the same premise.
% CHANGED (fidelity review, 2026-11-15, R26/R50/R69): the status line lists the declarations' hypotheses beside the node's (LINKAGE rule 6).
\statusT\quad\emph{(Cites nothing of its own: the step measure is read from
Lemma~\ref{lem:knot-exponent-ladder}(3) and the step of a grid ladder from
Proposition~\ref{prop:grid-pyramids-infinitely-divisible}(1), and the cited facts those nodes
spend, the \Levy{}--Khintchine representation, Khintchine's theorem and the lattice lemma, enter
only through them. Every clause is machine-checked; the chain rule for the relative entropy, the
disintegration of a finite measure on $\RR^2$ and the jointly measurable densities of a kernel are
Mathlib's. Two declarations require less than the node, their unused hypotheses dropped on 2026-11-15.
\textbf{Hypotheses} (the node's, each with what its declarations require): $q > 1$: same; $\Psi_1$ a
knot exponent: absent from \texttt{ladder\_bound\_1} and \texttt{\_2}, and from
\texttt{ladder\_bound\_3}, which has no $\Psi_1$; $\mu$ the step measure: same
(\texttt{IsStepMeasure}); a member of the orbit with its disintegration: same, the member's pair
(\texttt{IsFamilyPair}) and a disintegration (\texttt{IsPhaseDisintegration}) taken as hypotheses;
(1) $0 < |\mu|$: absent (when $|\mu| = 0$ the marginal $\tilde\rho$ vanishes and every clause is
trivial); $|\mu| < \infty$: same; ``the bound (\ref{eq:line-e6}) is finite'': same, read as
$\mathrm{KL}(\bar\mu^\circ \,\|\, \eta) < \infty$, which it is for $0 < |\mu| < \infty$; (2)
$\mathfrak{E} > -\infty$ at a member: same; (3) $w \ne \delta_0$ a mask on $h\ZZ$: same
(\texttt{IsGridMask}, a symmetric probability law on $h\ZZ$, with $h > 0$); a family carrying the
ladder of the grid pyramid: same (\texttt{CarriesGridLadder}), with its pair and a
disintegration.)}
\end{proposition}

% CHANGED (wave 3f, 2026-11-15): the proof follows the machine-checked route. The step formula
% mu(A) = int pi_r({theta : r theta in A}) rho~(dr) is proved first and carries (1) and (2). The
% rotation is extended by the identity off [1,q) to a Borel bijection of R, under which the relative
% entropy is invariant. The identity of (1) is proved for every mixture, in [0,infinity], from the
% integral form of the conditional relative entropy, a swap of coordinates and Mathlib's chain rule,
% in place of convexity and the integrability of |log p|; the bound (\ref{eq:line-e6}) and the equality case are
% read off it. (3) names the knot exponent of the family and the infinite divisibility of its step.
% CHANGED (fourth blind pass, 2026-11-15): mu(A) is read on (0, infinity); the measurability of
% r -> KL(pi_r || eta) is proved here, where it was read from inside another proof; P_r := eta for
% r <= 0 and its measurability are said; the case of the mixture identity where Q has no density is
% argued through the uniqueness of the disintegration; (3) says what a family carrying the ladder is
% and that the step measure is nonzero.
% CHANGED (the theory paper, 2026-11-16): the closing sentence on checking and cited facts shortened; no step changed.
% CHANGED (supplement review, 2026-11-16): the clauses are proved in the stated order, (1), (2),
% (3), after the step formula; the measurability of r -> KL(pi_r || eta) is cited from
% lem:kl-measurable. No step changed.
\begin{proof}
\leanok
\uses{thm:surface-measure, thm:knot-atom-cascades, prop:grid-pyramids-infinitely-divisible,
lem:knot-exponent-ladder, lem:family-structure, lem:line-axioms-kernels, lem:kl-measurable}
Let $(\Lambda, \mathcal{B})$ be the pair of the member and $\sigma$ its cylinder measure. The step
measure $\mu$ lives on $(0,\infty)$; for a Borel $A \subseteq \RR$ write $\mu(A)$ for
$\mu(A \cap (0,\infty))$. The \emph{step formula} comes first: for Borel $A \subseteq \RR$,
\[
  \mu(A) = \sigma(\{(\theta, r) : r\theta \in A\}) = \int\pi_r(\{\theta : r\theta \in A\})\,\tilde\rho(dr).
  \qquad (*)
\]
By Lemma~\ref{lem:knot-exponent-ladder}(3), $\mu(A) = \Lambda((1,q] \times A)$. By
Theorem~\ref{thm:surface-measure}(iii) the lines $\{1\} \times (0,\infty)$ and
$\{q\} \times (0,\infty)$ carry no mass, and $\Lambda$ is carried by $(0,\infty)^2$, so
$\Lambda((1,q] \times A) = \Lambda([1,q) \times (A \cap (0,\infty)))$. The cylinder measure is the
image of $\Lambda$ on $[1,q) \times (0,\infty)$ under $(y, v) \mapsto (\theta, r) = (y, v/y)$, and
$v = r\theta$ there; this is the first equality, and the disintegration gives the second. By
Theorem~\ref{thm:knot-atom-cascades}(2), $\sigma$ is carried by the cylinder, so $\tilde\rho$ is
carried by $(0,\infty)$.

(1) The rotation $\mathrm{rot}_r$ is defined for $r > 0$. Put $P_r := \eta$ for $r \le 0$, a
$\tilde\rho$-null set of $r$, which changes no integral against $\tilde\rho$. The map
$r \mapsto P_r(E) = \pi_r(\{\theta : [r\theta] \in E\})$, $r > 0$, is measurable for every Borel
$E$, since $(r, \theta) \mapsto [r\theta]$ is Borel on $(0,\infty) \times [1,q)$ and $(\pi_r)$ is a
measurable family; so $(P_r)$ is a measurable family. The measure $\mu$ is carried by $(0,\infty)$,
so $(*)$ at $A = \RR$ gives $|\mu| = \int\pi_r(\RR)\,\tilde\rho(dr) = |\tilde\rho|$. At
$A = \{v : [v] \in E\}$ it gives
$\mu^\circ(E) = \int\pi_r(\{\theta : [r\theta] \in E\})\,\tilde\rho(dr) = \int P_r(E)\,\tilde\rho(dr)$.

\emph{The rotation does not change the entropy.} Let $r > 0$. For $\theta \in [1,q)$ the class
$[r\theta]$ lies in $[1,q)$; it is $q^{-n}r\theta$ for an integer $n$, and $r^{-1}q^{-n}r\theta = q^{-n}\theta$
has the class $\theta$. So $\mathrm{rot}_{1/r} \circ \mathrm{rot}_r$ is the identity on $[1,q)$,
and $\mathrm{rot}_r$, extended by the identity off $[1,q)$, is a bijection $R$ of $\RR$ with $R$ and
$R^{-1}$ Borel. Both $\pi_r$ and $\eta$ are carried by $[1,q)$, so $P_r = R_*\pi_r$, and
$R_*\eta = (\mathrm{rot}_r)_*\eta = \eta$. Under a Borel bijection with Borel inverse the density
of the images is the image of the density, so the relative entropy is unchanged:
$\mathrm{KL}(P_r \,\|\, \eta) = \mathrm{KL}(\pi_r \,\|\, \eta)$, for $\tilde\rho$-almost every $r$.

\emph{The mixture identity.} Let $\bar\rho$ be a probability measure on $\RR$, $(P_r)$ a measurable
family of probability measures, $M := \int P_r\,\bar\rho(dr)$, and $\eta$ a probability measure.
Then, in $[0,\infty]$,
\[
  \int\mathrm{KL}(P_r \,\|\, \eta)\,\bar\rho(dr)
  = \mathrm{KL}(M \,\|\, \eta) + \int\mathrm{KL}(P_r \,\|\, M)\,\bar\rho(dr). \qquad (**)
\]
Let $Q(dr, d\theta) := \bar\rho(dr)P_r(d\theta)$. First, for a probability measure $\nu$,
$\mathrm{KL}(Q \,\|\, \bar\rho \otimes \nu) = \int\mathrm{KL}(P_r \,\|\, \nu)\,\bar\rho(dr)$. If $P_r$
is absolutely continuous with respect to $\nu$ for $\bar\rho$-almost every $r$, then $Q$ has the
density $(r, \theta) \mapsto (dP_r/d\nu)(\theta)$ against $\bar\rho \otimes \nu$, with the densities
chosen jointly measurable, and Tonelli's theorem integrates $f(dQ/d(\bar\rho \otimes \nu))$,
$f(x) = x\log x + 1 - x$, first in $\theta$. Otherwise the measurable set of $r$ at which $P_r$ is
not absolutely continuous with respect to $\nu$ has positive $\bar\rho$-mass, and the right side is
$+\infty$. So is the left: if $Q$ had a density $g$ against $\bar\rho \otimes \nu$, then
$\bar\rho(dr)\,g(r, \theta)\nu(d\theta)$ would be a second disintegration of $Q$ over $\bar\rho$,
and two disintegrations agree for $\bar\rho$-almost every $r$, so $P_r = g(r, \cdot)\,\nu$ would be
absolutely continuous with respect to $\nu$ for $\bar\rho$-almost every $r$. Second, the relative entropy is
unchanged when both coordinates are swapped. Swapped, $\bar\rho \otimes \nu$ becomes
$\nu \otimes \bar\rho$, and $Q$ becomes a finite measure on $\RR^2$ whose first marginal is $M$;
it disintegrates as $M(d\theta)Q'_\theta(dr)$ for a measurable family of probability measures
$Q'_\theta$. The chain rule for the relative entropy gives
$\mathrm{KL}(M \otimes Q' \,\|\, \eta \otimes \bar\rho)
= \mathrm{KL}(M \,\|\, \eta) + \mathrm{KL}(M \otimes Q' \,\|\, M \otimes \bar\rho)$. Swapping back and
applying the first fact with $\nu = \eta$ and with $\nu = M$ gives $(**)$.

Let $0 < |\mu| < \infty$ and $\bar\rho := \tilde\rho/|\mu|$, a probability measure; then
$\bar\mu^\circ = \int P_r\,\bar\rho(dr)$. By the rotation step and $(**)$, multiplied by $|\mu|$,
\[
  -\mathfrak{E} = \int\mathrm{KL}(\pi_r \,\|\, \eta)\,\tilde\rho(dr)
  = \int\mathrm{KL}(P_r \,\|\, \eta)\,\tilde\rho(dr)
  = |\mu|\,\mathrm{KL}(\bar\mu^\circ \,\|\, \eta) + \int\mathrm{KL}(P_r \,\|\, \bar\mu^\circ)\,\tilde\rho(dr),
\]
in $[0,\infty]$, whether or not $\mathfrak{E} > -\infty$. The last integral is nonnegative, which is
(\ref{eq:line-e6}). When the bound is finite, equality holds in (\ref{eq:line-e6}) iff the last integral vanishes. Its
integrand is measurable (Lemma~\ref{lem:kl-measurable}) and nonnegative, so this holds iff
$\mathrm{KL}(P_r \,\|\, \bar\mu^\circ) = 0$ for $\tilde\rho$-almost every $r$, that is, iff
$P_r = \bar\mu^\circ$ for $\tilde\rho$-almost every $r$, since both are probability measures.

(2) The map $r \mapsto \mathrm{KL}(\pi_r \,\|\, \eta)$ is measurable
(Lemma~\ref{lem:kl-measurable}). So if $\mathfrak{E} > -\infty$, then
$\mathrm{KL}(\pi_r \,\|\, \eta)$ is finite, and $\pi_r$ is absolutely continuous with respect to $\eta$, for $\tilde\rho$-almost every $r > 0$; and $\eta$ is absolutely continuous with respect to
Lebesgue measure. Let $N$ be Lebesgue-null, and enlarge it to a Borel null set $T$. For $r > 0$
the set $\{\theta : r\theta \in T\} = r^{-1}T$ is Lebesgue-null, hence $\pi_r$-null for
$\tilde\rho$-almost every $r$, and $\mu(N) \le \mu(T) = 0$ by $(*)$. A measure with a nonzero
singular part is not absolutely continuous, so then $\mathfrak{E} = -\infty$. This argument does not
use that $\Psi_1$ is a knot exponent.

(3) A family carrying the ladder of the grid pyramid of $w$ is, as in
Definition~\ref{def:grid-pyramid-line}, a $q^{\ZZ}$-covariant family in the canonical discrete
gauge whose kernels have the steps $\mu_{q^n,q^{n+1}} = T_{q^n}w$. Its kernels are symmetric,
$q$-covariant with $S_qt = qt$ (Lemma~\ref{lem:line-axioms-kernels}), with the step
$\mu_{1,q} = w$. It is
therefore a member of the modulation orbit of $\Psi_1 := -\log\widehat\mu_{0,1}$, whose step
$\Psi_1(q\,\cdot) - \Psi_1 = \Psi_q - \Psi_1$ is the exponent of $\mu_{1,q} = w$. The step of a
family is infinitely divisible (Lemma~\ref{lem:family-structure}), so by
Proposition~\ref{prop:grid-pyramids-infinitely-divisible}(1) the mask $w$ has a finite folded
\Levy{} measure $\varpi_w \ne 0$, carried by $h\ZZ_{>0}$, with
$\widehat w(\omega) = \exp(-\int(1 - \cos\omega v)\,\varpi_w(dv))$. So $(0, \varpi_w)$ are the data
of the step, and the step measure is $\varpi_w$, nonzero as $w \ne \delta_0$. A countable set is
Lebesgue-null, so the nonzero $\varpi_w$, carried by one, is not absolutely continuous with respect
to Lebesgue measure, and $\mathfrak{E} = -\infty$ by (2).

Every step is machine-checked, and the proof cites no fact of its own.
\end{proof}

The bound (\ref{eq:line-e6}) depends only on the ladder. When the step is finite and the folded step
$\bar\mu^\circ$ has infinite entropy against Haar, every member has $\mathfrak{E} = -\infty$. The same holds by (2) for the ladder of Proposition~\ref{prop:discrete-gaussian-pyramid},
whose step is carried by one orbit $\{hq^i\}$, and for a ladder carried by one knot-atom orbit.
Where the folded step has finite entropy, the bound is attained exactly by the members whose
rotated fibres all equal the folded step. The next theorem finds them when, in addition, the ladder
has no Gaussian part, the density of the folded step is positive, and the step is that density times
a covariant step.

% draft: none in the draft's statements; notes/phase-entropy-without-covariant-member.md, Theorem 3.
% CHANGED (from the note): the hypothesis KL(folded step || Haar) < infinity added. Without it the
% bound (\ref{eq:line-e6}) is -infinity, every member is a maximizer, and the note's uniqueness claim is false.
% The hypotheses "no Gaussian part" (as in thm:phase-entropy-selection(4)) and the integrability of
% k are stated; the note leaves them implicit.
% CHANGED: the density of the folded step renamed p (as in the proof of
% prop:phase-entropy-ladder-bound), and its log-periodic lift written p([v]); the letter G is the step
% exponent of prop:ladder-non-creation-line (blind review 2, flag 20).
\begin{theorem}[selection with a log-periodic factor]
\label{thm:phase-entropy-modulated-selection}
\lean{Pyramids.phase_entropy_modulated_selection,
Pyramids.phase_entropy_modulated_selection_converse, Pyramids.phase_entropy_modulated_selection_iff}
\leanok
\uses{prop:phase-entropy-ladder-bound, def:modulation-orbit, thm:phase-entropy-selection}
Let $\Psi_1$ be a knot exponent whose ladder has no Gaussian part, with step measure $\mu$,
$0 < |\mu| < \infty$. Suppose
\begin{enumerate}
\item $\bar\mu^\circ$ is absolutely continuous with respect to $\eta$, with a density $p$ that is
positive $\eta$-almost everywhere, and
$\mathrm{KL}(\bar\mu^\circ \,\|\, \eta) = \int p\log p\,d\eta < \infty$;
\item for a nonincreasing $k \ge 0$ with $\int(1 \wedge x^2)\,k(x)\,dx/x < \infty$,
\[
  \mu(dv) = p([v])\,\bigl(k(v/q) - k(v)\bigr)\frac{dv}{v}.
\]
\end{enumerate}
Then the phase entropy has exactly one maximizer on the modulation orbit: the family with
$a \equiv 0$ and
\[
  \varpi^*_t(dv) = p([v])\,k(v/t)\,\frac{dv}{v} \qquad (t > 0),
\]
and the maximum is
$-|\mu|\,\mathrm{KL}(\bar\mu^\circ \,\|\, \eta) = -\int\log p([v])\,\mu(dv)$. If $p$ is not
almost everywhere constant, the ladder has no similarity member.
% CHANGED (2026-11-23, review round 3 batch C, supp.S3): the converse, so that the theorem is an
% equivalence between (2) and the attainment of the bound.
Conversely, if (1) holds and some member of the orbit attains the bound (\ref{eq:line-e6}), then
(2) holds for some $k$. So, under the other hypotheses and (1), the bound is attained iff (2)
holds.
% CHANGED (supplement review, 2026-11-16): the closing comparison with
% thm:phase-entropy-selection(3) moves to the remark after the proof.
% CHANGED (fidelity review, 2026-11-15, R64): the status line says that ``exactly one maximizer'' is
% stated in decomposed form.
% CHANGED (fidelity review, 2026-11-15, R26/R50/R69): the status line lists the declarations' hypotheses beside the node's (LINKAGE rule 6).
\statusT\quad\emph{(Cites nothing of its own: the \Levy{}--Khintchine representation and
Khintchine's theorem enter only through Theorem~\ref{thm:surface-measure},
Lemma~\ref{lem:knot-exponent-ladder}(3), Lemma~\ref{lem:phase-disintegration-exists} and
Proposition~\ref{prop:phase-entropy-ladder-bound}(1). Every clause is machine-checked, and the
declaration reads every hypothesis of the node; the Stieltjes measure of a monotone function, the
density of $p\,\eta$ against $\eta$ and the relative entropy as an integral are Mathlib's.
\textbf{Hypotheses} (the node's, each with what its
declarations require): $q > 1$: same; $\Psi_1$ a knot exponent: same; no Gaussian part: same; $\mu$ the
step measure: same; $0 < |\mu| < \infty$: same; (1) $\bar\mu^\circ = p\,\eta$: same, with $p$
measurable, stronger in form and harmless, a density being chosen measurable; $p > 0$
$\eta$-almost everywhere: same; $\mathrm{KL}(\bar\mu^\circ \,\|\, \eta) < \infty$: same; (2) $k \ge 0$
nonincreasing: same, on $(0,\infty)$; $\int(1 \wedge x^2)\,k(x)\,dx/x < \infty$: same; the form of
$\mu$: same. ``Exactly one maximizer'' is stated in decomposed form, the reader composing it: the
declaration characterizes the maximizers as the members with $\mathcal{B} = 0$ and
$\varpi_t = \varpi^*_t$ for $t > 0$, and that these data fix one family is
Theorem~\ref{thm:surface-measure} (\texttt{surface\_measure\_unique},
\texttt{surface\_measure\_injective}). The converse,
\texttt{phase\_entropy\_modulated\_selection\_converse}: weaker, it reads neither $\Psi_1$ a knot
exponent, nor the absence of a Gaussian part, nor $|\mu| > 0$, nor $p > 0$; the member attaining
the bound: same, with its pair $(\Lambda, \mathcal{B})$ and a disintegration $\pi$.
\texttt{phase\_entropy\_modulated\_selection\_iff} reads every hypothesis of the node but (2), the
two sides of its equivalence being the attainment and (2).)}
\end{theorem}

% CHANGED (wave 3f, 2026-11-15): the proof follows the machine-checked route. It says why
% k(infinity) = 0 and builds rho~_k as the image under exp of a Stieltjes measure in log scale. The
% knot measure is identified through the uniqueness argument of prop:ladder-knot-kernel(1), read
% through the symmetric lift, after the data of Psi_1 are computed; (iv) is proved on one orbit
% interval from the null sets of sigma*, in place of the strict decrease of k. Uniqueness inverts the
% rotation explicitly, and the converse of the equality case (a member with B = 0 and the
% candidate's slabs attains the bound) is proved. The last clause reads the bound at a similarity
% member, where the phase entropy is 0, in place of the telescoping of a similarity step.
% CHANGED (fourth blind pass, 2026-11-15): the facts this proof read from inside other proofs are
% written out: the image form and the slabs of Lambda_k (from the proof of
% prop:strata-cylinder-coordinates(1)), the knot measure as a series (prop:ladder-knot-kernel(1)'s
% argument), the step formula (prop:phase-entropy-ladder-bound's (*)), the density of the smeared
% measure (lem:haar-smearing-injective's proof), B = 0 (thm:phase-entropy-selection(3)'s proof), the
% determination of Lambda by its slabs, and the phase entropy of a similarity member; and log p is
% shown integrable against the folded step.
% CHANGED (the theory paper, 2026-11-16): the closing sentence on checking and cited facts shortened; no step changed.
% CHANGED (supplement review, 2026-11-16): the density of the smeared measure is cited from
% lem:haar-smearing-density and B = 0 from lem:member-gaussian-part(2). No step changed.
\begin{proof}
\leanok
\uses{prop:phase-entropy-ladder-bound, lem:haar-smearing-injective, thm:surface-measure,
thm:knot-atom-cascades, prop:strata-cylinder-coordinates, lem:knot-exponent-ladder,
lem:line-exponent-data, lem:phase-disintegration-exists, prop:ladder-knot-kernel,
thm:phase-entropy-selection, lem:haar-smearing-density, lem:member-gaussian-part}
Write $\bar p(v) := p([v])$ for $v > 0$. It satisfies $\bar p(q^nv) = \bar p(v)$ for $n \in \ZZ$,
and $\bar p > 0$ for $dv/v$-almost every $v$. Indeed, $\int G\,dv/(\ell v) =
\int\sum_nG(\theta q^n)\,\eta(d\theta)$ for Borel $G \ge 0$, and for the indicator $G$ of
$\{\bar p \le 0\}$ the integrand is $\sum_n\mathbf 1_{\{p \le 0\}}(\theta)$, which vanishes for
$\eta$-almost every $\theta$.

\emph{The measure $\tilde\rho_k$.} The limit $k(\infty) \ge 0$ exists because $k$ is
nonincreasing and nonnegative, and it is $0$: if $k \ge c > 0$ on $(0,\infty)$, then for $r > 1$
\[
  \int_0^\infty(1 \wedge x^2)\,k(x)\,\frac{dx}{x} \ge c\int_0^r(1 \wedge x^2)\,\frac{dx}{x}
  = c\Bigl(\frac12 + \log r\Bigr),
\]
which is unbounded in $r$. The function $G(u) := -\ell k(e^u)$ is nondecreasing on $\RR$ and tends
to $0$ as $u \to \infty$; $k(0{+})$ may be infinite, which is why the construction is made in the
variable $u = \log x$. Let $F(u) := G(u{+})$ be its right-continuous version, $dF$ the measure with
$dF((a,b]) = F(b) - F(a)$, and $\tilde\rho_k$ the image of $dF$ under $u \mapsto e^u$. Then
$\tilde\rho_k$ is carried by $(0,\infty)$, it is $\sigma$-finite, and
$\tilde\rho_k([x,\infty)) = -F(\log x\,{-})$ is finite for $x > 0$. Off the countable set of $x$
at which $F$ or $G$ jumps at $\log x$, $F(\log x\,{-}) = F(\log x) = G(\log x)$, so
$\tilde\rho_k([x,\infty)) = \ell k(x)$.

\emph{The candidate and its slabs.} Let $\sigma^*(d\theta, dr) := \bar p(r\theta)\,\eta(d\theta)\,
\tilde\rho_k(dr)$, a measure on the cylinder, and $\Lambda^* := \Lambda_{\sigma^*}$. By
Theorem~\ref{thm:knot-atom-cascades}(2), $\Lambda^*$ is $J$-invariant, carried by $(0,\infty)^2$,
and its cylinder measure is $\sigma^*$. The orbit of $(\theta, r)$ consists of the points
$(\theta q^n, r\theta q^n)$, at which $\bar p$ of the second coordinate is $\bar p(r\theta)$; so
$\Lambda^*(dy, dv) = \bar p(v)\,\Lambda_k(dy, dv)$ with $\Lambda_k := \Lambda_{\eta \otimes
\tilde\rho_k}$. For Borel $g \ge 0$ on $(0,\infty)^2$ and $r > 0$,
$\int\sum_ng(\theta q^n, r\theta q^n)\,\eta(d\theta) = \int_{(0,\infty)}g(y, ry)\,dy/(\ell y)$, the
intervals $q^n[1,q)$ tiling $(0,\infty)$ and $dy/y$ being invariant under $y \mapsto q^ny$. So
$\Lambda_k$ is the image of $(dy/(\ell y)) \otimes \tilde\rho_k$ under $(y, r) \mapsto (y, ry)$.
For Borel $A \subseteq (0,\infty)$, the substitution $v = ry$, with $dy/y = dv/v$, and Tonelli's
theorem give
\[
  \Lambda_k((0,t] \times A) = \int\int_{(0,rt]}\mathbf 1_A(v)\,\frac{dv}{\ell v}\,\tilde\rho_k(dr)
  = \int_A\ell^{-1}\tilde\rho_k([v/t,\infty))\,\frac{dv}{v} .
\]
The exceptional $v$ form a countable set, so
\[
  \varpi^*_t(dv) := \Lambda^*((0,t] \times dv) = \bar p(v)\,k(v/t)\,\frac{dv}{v} \qquad (t > 0).
\]

\emph{The knot measure.} Let $(a(1), \varpi_1)$ be the data of $\Psi_1$
(Lemma~\ref{lem:line-exponent-data}(1)). By Lemma~\ref{lem:line-exponent-data}(2), $\Psi_1(q\,\cdot)$
has the data $(q^2a(1), T_q\varpi_1)$. The step $\phi = \Psi_1(q\,\cdot) - \Psi_1$ lies in
$\mathrm{ND}_s$, so by Lemma~\ref{lem:line-exponent-data}(3) it has the data
$(q^2a(1) - a(1), \rho)$ with $\rho + \varpi_1 = T_q\varpi_1$; by the uniqueness of the data,
$\rho = \mu$ and $(q^2 - 1)a(1) = 0$, the Gaussian coefficient of the step. So $a(1) = 0$ and
$T_q\varpi_1 = \mu + \varpi_1$. Applying $T_{q^{-1}}$ $M$ times,
$\varpi_1 = \sum_{m=1}^MT_{q^{-m}}\mu + T_{q^{-M}}\varpi_1$. For a Borel
$B \subseteq (\varepsilon,\infty)$, $\varepsilon > 0$, the last term is
$\varpi_1(q^MB) \le \varpi_1((q^M\varepsilon,\infty))$, which tends to $0$, a \Levy{} measure being
finite on $(\varepsilon,\infty)$. The sets $(1/n,\infty)$ exhaust $(0,\infty)$, so
$\varpi_1 = \sum_{m \ge 1}T_{q^{-m}}\mu$. On the other hand $T_{q^{-m}}\mu$ has the density
$\bar p(v)\bigl(k(q^{m-1}v) - k(q^mv)\bigr)/v$, and the sum over $m \ge 1$ telescopes to
$\bar p(v)k(v)/v$ because $k(\infty) = 0$. So $\varpi^*_1 = \varpi_1$, a \Levy{} measure. This is
(i) of Theorem~\ref{thm:surface-measure} at $t = 1$, hence at every $t > 0$.

\emph{The candidate is a member.} Take $\mathcal{B} := 0$, so (ii) holds. For (iii), no line
$\{\theta\} \times (0,\infty)$ carries mass for $\sigma^*$, since $\eta$ has no atoms; by
Theorem~\ref{thm:knot-atom-cascades}(2) no line $\{y\} \times (0,\infty)$ carries mass for
$\Lambda^*$. For (iv), first $\tilde\rho_k \ne 0$: otherwise $\Lambda^* = 0$, so $\varpi^*_q = 0$,
while $\mu \le \varpi^*_q$ because $k \ge 0$, and $\mu \ne 0$. Next, $\eta \otimes \tilde\rho_k$ and
$\sigma^*$ have the same null sets: for $r > 0$, $\mathrm{rot}_r$ preserves $\eta$, so
$\bar p(r\theta) = p(\mathrm{rot}_r\theta) > 0$ for $\eta$-almost every $\theta$. Let $s < t$, put
$u := (s + t)/2$, take $n$ with $q^n \le u < q^{n+1}$, and let
$I := (q^{-n}u, \min(q, q^{-n}t))$, a nonempty interval in $[1,q)$. For $\theta \in I$ and $r > 0$,
$\theta q^n \in (s,t]$ and $r\theta q^n > 0$, so
\[
  \Lambda^*((s,t] \times (0,\infty)) \ge \sigma^*(I \times (0,\infty)) > 0,
\]
because $\eta(I)\,\tilde\rho_k((0,\infty)) > 0$. By Theorem~\ref{thm:surface-measure} the pair
$(\Lambda^*, 0)$ is the pair of a family. Its knot kernel has the exponent
$\int(1 - \cos\omega v)\,\varpi^*_1(dv)$, which is $\Psi_1$ since $\Psi_1$ has the data
$(0, \varpi_1)$ and $\varpi^*_1 = \varpi_1$. So the family is a member of the orbit, with
$a \equiv 0$ and the slabs $\varpi^*_t$.

\emph{Its phase entropy.} For $r > 0$ let $\pi^*_r := (p \circ \mathrm{rot}_r)\,\eta$, and
$\pi^*_r := \eta$ for $r \le 0$. The image of $(g \circ T)\,\nu$ under $T$ is $g\,T_*\nu$, so
$(\mathrm{rot}_r)_*\pi^*_r = p\,(\mathrm{rot}_r)_*\eta = p\,\eta = \bar\mu^\circ$, a probability
measure; hence $\pi^*_r$ is one. By Tonelli,
$\sigma^*(E) = \int\pi^*_r(\{\theta : (\theta, r) \in E\})\,\tilde\rho_k(dr)$ for Borel $E$, so the
$r$-marginal of $\sigma^*$ is $\tilde\rho_k$ and $\pi^*$ is a disintegration of $\sigma^*$, with
$P^*_r = \bar\mu^\circ$ for every $r > 0$. By the equality case of
Proposition~\ref{prop:phase-entropy-ladder-bound}(1), which applies because the bound is finite by
hypothesis, $\mathfrak{E} = -|\mu|\,\mathrm{KL}(\bar\mu^\circ \,\|\, \eta)$ at the candidate.

\emph{The bound, and equality exactly at the candidate.} The bound for every member is
Proposition~\ref{prop:phase-entropy-ladder-bound}(1). Let a member with pair $(\Lambda, \mathcal{B})$
and disintegration $\pi$ attain it. By the equality case, $P_r = \bar\mu^\circ$ for
$\tilde\rho$-almost every $r > 0$. On $[1,q)$, $\mathrm{rot}_{1/r} \circ \mathrm{rot}_r$ is the
identity, and $\pi_r$ is carried by $[1,q)$, so
$\pi_r = (\mathrm{rot}_{1/r})_*(p\,\eta)$. Since $p = (p \circ \mathrm{rot}_r) \circ
\mathrm{rot}_{1/r}$ $\eta$-almost everywhere and $\mathrm{rot}_{1/r}$ preserves $\eta$, this is
$(p \circ \mathrm{rot}_r)\,\eta$. So $\sigma(d\theta, dr) = \bar p(r\theta)\,\eta(d\theta)\,
\tilde\rho(dr)$, and its image under $(\theta, r) \mapsto r\theta$ is $\bar p\,\nu$, where $\nu$ is
the image of $\eta \otimes \tilde\rho$. That image is $\mu$. Indeed, by
Lemma~\ref{lem:knot-exponent-ladder}(3), $\mu(A) = \Lambda((1,q] \times A)$ for Borel
$A \subseteq (0,\infty)$; by Theorem~\ref{thm:surface-measure}(iii) the lines $\{1\} \times
(0,\infty)$ and $\{q\} \times (0,\infty)$ carry no mass, so this is $\Lambda([1,q) \times A)$; and
$\sigma$ is the image of $\Lambda$ on $[1,q) \times (0,\infty)$ under $(y, v) \mapsto (y, v/y)$
(Theorem~\ref{thm:knot-atom-cascades}(2)), with $v = r\theta$. The measure $\tilde\rho$ is carried
by $(0,\infty)$ and finite on every $[x,\infty)$, $x > 0$
(Lemma~\ref{lem:phase-disintegration-exists}), hence $\sigma$-finite. By
Lemma~\ref{lem:haar-smearing-density}, $\nu$ has the density $\tilde\rho((v/q, v])$ against
$dv/(\ell v)$ on $(0,\infty)$. So it does not charge the $dv/v$-null set $\{\bar p = 0\}$, and
$\nu = \bar p^{-1}\mu$. The candidate
attains the bound too, so likewise $\nu_k = \bar p^{-1}\mu$ for the image $\nu_k$ of
$\eta \otimes \tilde\rho_k$. Both $\tilde\rho$ and $\tilde\rho_k$ are carried by $(0,\infty)$ and
finite on every $[x,\infty)$, $x > 0$ (Lemma~\ref{lem:phase-disintegration-exists}), so
$\tilde\rho = \tilde\rho_k$ by Lemma~\ref{lem:haar-smearing-injective}. Hence $\sigma = \sigma^*$,
and $\Lambda = \Lambda_\sigma = \Lambda^*$ by Theorem~\ref{thm:knot-atom-cascades}(2). Next
$\mathcal{B} = 0$, because the ladder has no Gaussian part: $\mathcal{B}((0,t]) = 0$ for every
$t > 0$ by Lemma~\ref{lem:member-gaussian-part}(2), and $\mathcal{B}$ is carried by $(0,\infty)$. So
$\varpi_t = \varpi^*_t$ for every $t > 0$.

Conversely, let a member have $\mathcal{B} = 0$ and $\varpi_t = \varpi^*_t$ for every $t > 0$. Then
$\Lambda$ and $\Lambda^*$ agree on the sets $(s,t] \times (a,b]$, $0 \le s < t$, $0 < a < b$, both
giving $\varpi_t((a,b]) - \varpi_s((a,b])$, a finite number. These sets form a $\pi$-system that
generates the Borel sets of $(0,\infty)^2$, and $(0,\infty)^2$ is the increasing union of the sets
$(0,n] \times (1/n, n]$ of finite mass, so $\Lambda = \Lambda^*$. Two disintegrations of $\sigma^*$ agree for $\tilde\rho_k$-almost every $r$
(Lemma~\ref{lem:phase-disintegration-exists}), so $P_r = P^*_r = \bar\mu^\circ$ almost everywhere,
and equality holds.

% CHANGED (2026-11-23, review round 3 batch C, supp.S3): the converse, on the checked route
% (phase_entropy_modulated_selection_converse); it reads only (1) and the member at equality.
\emph{The converse.} Let (1) hold, and let a member with pair $(\Lambda, \mathcal{B})$, cylinder
measure $\sigma$, $r$-marginal $\tilde\rho$ and disintegration $\pi$ attain the bound. The first
part of the uniqueness argument above uses only (1) and the equality. It gives $\mu = \bar p\,\nu$,
where $\nu$ is the image of $\eta \otimes \tilde\rho$ under $(\theta, r) \mapsto r\theta$, with
the density $\tilde\rho((v/q, v])$ against $dv/(\ell v)$ on $(0,\infty)$
(Lemma~\ref{lem:haar-smearing-density}). Put $k(x) := \ell^{-1}\tilde\rho((x,\infty))$ for $x > 0$,
finite by Lemma~\ref{lem:phase-disintegration-exists}. Then $k \ge 0$ is nonincreasing, and
$k(v/q) - k(v) = \ell^{-1}\tilde\rho((v/q, v])$ for $v > 0$, so
$\mu(dv) = \bar p(v)\,(k(v/q) - k(v))\,dv/v$. For the integrability, $\Lambda = \Lambda_\sigma$
satisfies (i) of Theorem~\ref{thm:surface-measure}, so $\sigma$ satisfies (\ref{eq:line-e3}) by
Theorem~\ref{thm:knot-atom-cascades}(2). With $w(r) := r^2 \wedge (1 + \log_+r)$, nondecreasing,
and $\theta \ge 1$ on $\mathcal{C}$, this gives $\int w\,d\tilde\rho \le \int w(r\theta)\,
\sigma(d\theta, dr) < \infty$. By Tonelli's theorem,
$\int_0^\infty(1 \wedge x^2)\,\tilde\rho([x,\infty))\,dx/x = \int g\,d\tilde\rho$ with
$g(r) := \int_0^r(1 \wedge x^2)\,dx/x$, which is $r^2/2$ for $r \le 1$ and $\tfrac12 + \log r$ for
$r > 1$, so $g \le w$. Since $k(x) \le \ell^{-1}\tilde\rho([x,\infty))$,
$\int(1 \wedge x^2)\,k(x)\,dx/x \le \ell^{-1}\int w\,d\tilde\rho < \infty$. So (2) holds.

\emph{The value.} The measure $\bar\mu^\circ = p\,\eta$ is absolutely continuous with respect to
$\eta$, with density $p$, positive $\bar\mu^\circ$-almost everywhere, and both are probability
measures. The function $\log p$ is $\bar\mu^\circ$-integrable: $\int|\log p|\,d\bar\mu^\circ =
\int p\,|\log p|\,d\eta$, the negative part of $p\log p$ is at most $1/e$ and $\eta$ is a
probability, and the positive part has a finite integral since $\int p\log p\,d\eta < \infty$. So
$\mathrm{KL}(\bar\mu^\circ \,\|\, \eta) = \int p\log p\,d\eta = \int\log p\,d\bar\mu^\circ$. Since
$\bar\mu^\circ$ is $|\mu|^{-1}$ times the image of $\mu$ under $v \mapsto [v]$,
\[
  |\mu|\,\mathrm{KL}(\bar\mu^\circ \,\|\, \eta) = \int\log p([v])\,\mu(dv).
\]

\emph{No similarity member.} A similarity member has $\Psi_{\lambda t} = \Psi_t(\lambda\,\cdot)$
for all $\lambda, t > 0$ (Theorem~\ref{thm:phase-entropy-selection}(1)), so by ``both parts'' of
Proposition~\ref{prop:strata-cylinder-coordinates}(1) it satisfies (a) there,
$\varpi_{\lambda t} = T_\lambda\varpi_t$, and so (c): $\sigma = \eta \otimes \tilde\rho$. The
constant family $\eta$ is then a disintegration of $\sigma$, two disintegrations agree for
$\tilde\rho$-almost every $r$ (Lemma~\ref{lem:phase-disintegration-exists}), so $\pi_r = \eta$ for
$\tilde\rho$-almost every $r$, and its phase entropy is $0$ by
Theorem~\ref{thm:phase-entropy-selection}(2). The bound (\ref{eq:line-e6}) then gives
$|\mu|\,\mathrm{KL}(\bar\mu^\circ \,\|\, \eta) \le 0$. Since $|\mu| > 0$, the relative entropy
vanishes, so $\bar\mu^\circ = \eta$, that is $p\,\eta = \eta$, and $p = 1$ $\eta$-almost
everywhere.

Every step is machine-checked, and the cited facts, \ledger{A39} and \ledger{A41}, enter only
through the statements the proof cites.
\end{proof}

% CHANGED (pass 6, 2026-11-22, the author's review): the example that the class is strictly larger,
% argued here in prose, is the next proposition; the prose cites it.
With $p \equiv 1$ the theorem is Theorem~\ref{thm:phase-entropy-selection}(3) for a finite step.
Among ladders with a finite step and no Gaussian part its class is strictly larger. The ladders of
the next proposition meet its hypotheses and have no similarity member.

% draft: none. New at pass 6 (2026-11-22, the author's review of the supplement): the example of
% the remark after thm:phase-entropy-modulated-selection, stated and proved
% (Formalization/Pyramids/LineAppendix/ModulatedStrict.lean). The prose's "covariant step with
% finite mass" is stated as hypotheses on k: the covariant step has the mass (log q) k(0+), so
% "finite and nonzero" is "k bounded and not identically zero".
\begin{proposition}[the class is strictly larger]
\label{prop:modulated-selection-strict}
\lean{Pyramids.modulated_selection_strict}
\leanok
\uses{thm:phase-entropy-modulated-selection, def:modulation-orbit, lem:line-exponent-data}
Let $0 < |\varepsilon| < 1$, and let $k \ge 0$ be nonincreasing and bounded on $(0,\infty)$, not
identically zero, with $\int(1 \wedge x^2)\,k(x)\,dx/x < \infty$. Put
$p(\theta) := 1 + \varepsilon\cos(2\pi\log_q\theta)$ and
\[
  \mu(dv) := p(v)\,\bigl(k(v/q) - k(v)\bigr)\frac{dv}{v} .
\]
Then $\mu$ is the step measure of a ladder with no Gaussian part and $0 < |\mu| < \infty$. With
this $p$ and this $k$ the ladder satisfies hypotheses (1) and (2) of
Theorem~\ref{thm:phase-entropy-modulated-selection}, and $p$ is not almost everywhere constant. The
ladder has no similarity member.
\statusT\quad\emph{(Cites nothing of its own: the \Levy{}--Khintchine representation and
Khintchine's theorem, \ledger{A39} and \ledger{A41}, enter only through
Theorem~\ref{thm:phase-entropy-modulated-selection}, whose last clause gives the absence of a
similarity member, and through the uniqueness of the data, Lemma~\ref{lem:line-exponent-data}(1).
Every clause is machine-checked. \textbf{Hypotheses} (the node's, each with what its declaration
requires): $q > 1$: same; $0 < |\varepsilon| < 1$: same; $k \ge 0$ nonincreasing: same, on
$(0,\infty)$; $k$ bounded: same, an upper bound of $k$ on $(0,\infty)$; not identically zero:
same, a point $x > 0$ with $k(x) \ne 0$; $\int(1 \wedge x^2)\,k(x)\,dx/x < \infty$: same. ``A
ladder'' is a knot exponent $\Psi_1$, the step measure \texttt{IsStepMeasure}, hypothesis (1) the
three conditions of the theorem's declaration on $p$, hypothesis (2) the form of $\mu$ with
$p([v])$, and the similarity member as in Theorem~\ref{thm:phase-entropy-selection}.)}
\end{proposition}

\begin{proof}
\leanok
\uses{thm:phase-entropy-modulated-selection, lem:line-exponent-data}
Write $\ell := \log q$ and $K := \sup_{x > 0}k(x)$, the limit of $k$ at $0{+}$, which is finite
and positive. The limit of $k$ at $\infty$ is $0$: if $k \ge c > 0$, then
$\int(1 \wedge x^2)\,k(x)\,dx/x$ would diverge. The factor $p$ takes its values in
$[1 - |\varepsilon|, 1 + |\varepsilon|] \subset (0, 2)$, and $p(q^n\theta) = p(\theta)$ for
$n \in \ZZ$, so $p([v]) = p(v)$ for $v > 0$.

\emph{The folded step.} For Borel $G \ge 0$,
$\int G(v)\,dv/(\ell v) = \int\sum_{n \in \ZZ}G(\theta q^n)\,\eta(d\theta)$, because the intervals
$q^n[1,q)$ tile $(0,\infty)$ and $dv/v$ is invariant under $v \mapsto q^nv$. Against $dv/(\ell v)$
the measure $\mu$ has the density $\ell p(v)(k(v/q) - k(v))$. At the points $\theta q^n$ of an
orbit, $\theta \in [1,q)$, it takes the values $\ell p(\theta)(b_{n-1} - b_n)$ with
$b_n := k(\theta q^n)$. The $b_n$ do not increase with $n$, and they tend to $0$ as $n \to \infty$
and to $K$ as $n \to -\infty$. So the sum over $n \in \ZZ$ telescopes to $\ell Kp(\theta)$, and
$\mu^\circ = \ell K\,p\,\eta$. In the coordinate $s = \log_q\theta$ the measure $\eta$ is Lebesgue
measure on $[0,1)$, so $\int\cos(2\pi\log_q\theta)\,\eta(d\theta) = \int_0^1\cos 2\pi s\,ds = 0$
and $\int p\,d\eta = 1$. Hence $|\mu| = \ell K$, positive and finite, and $\bar\mu^\circ = p\,\eta$.
With $p \equiv 1$ the same computation gives the covariant step $(k(v/q) - k(v))\,dv/v$ the mass
$\ell K$: the hypotheses on $k$ say that it has finite positive mass. Hypothesis (1) holds: $p$ is
positive, and $\log p$ is bounded, so $\mathrm{KL}(\bar\mu^\circ \,\|\, \eta) = \int p\log p\,d\eta$
is finite. Hypothesis (2) is the definition of $\mu$, as $p([v]) = p(v)$.

\emph{The ladder.} Let $\varpi(dv) := p(v)\,k(v)\,dv/v$. Since $p \le 2$,
$\int(1 \wedge v^2)\,\varpi(dv) \le 2\int(1 \wedge x^2)\,k(x)\,dx/x < \infty$, so $\varpi$ is a
\Levy{} measure, and the exponent $\Psi_1$ with the data $(0, \varpi)$ lies in $\mathrm{ND}_s$. By
Lemma~\ref{lem:line-exponent-data}(2), $\Psi_1(q\,\cdot)$ has the data $(0, T_q\varpi)$, and
$T_q\varpi(dv) = p(v/q)\,k(v/q)\,dv/v = p(v)\,k(v/q)\,dv/v = \mu + \varpi$. So the step
$\Psi_1(q\,\cdot) - \Psi_1$ is the exponent with the data $(0, \mu)$. It lies in $\mathrm{ND}_s$,
has no Gaussian part and the step measure $\mu$. It is not identically zero: the data of an
exponent are unique (Lemma~\ref{lem:line-exponent-data}(1)), and $\mu \ne 0$. So $\Psi_1$ is a knot
exponent, and its ladder has the step measure $\mu$.

\emph{No similarity member.} The function $\cos(2\pi\log_q\theta)$ is positive on $(1, q^{1/4})$
and negative on $(q^{1/4}, q^{3/4})$, and both intervals have positive $\eta$-measure. As
$\varepsilon \ne 0$, $p$ is not almost everywhere constant, and the last clause of
Theorem~\ref{thm:phase-entropy-modulated-selection} says that the ladder has no similarity member.

Every step is machine-checked, and the cited facts, \ledger{A39} and \ledger{A41}, enter only
through the statements the proof cites. The formal proof reads
$\int\cos(2\pi\log_q\theta)\,\eta(d\theta) = 0$ as the real part of the coefficient $j = 1$ of
$\eta$, which vanishes because $\eta$ is invariant under the rotations.
\end{proof}

The selected pyramid is the covariant pyramid of $\mu/p([v])$, reweighted by the log-periodic
displacement profile $p([v])$. On a ladder whose step has a singular part, every grid pyramid among
them, the phase entropy is constant, so it selects nothing.

A functional that selects on every ladder and returns the covariant member whenever there is one is
\emph{open}. The second term of the identity in
Proposition~\ref{prop:phase-entropy-ladder-bound}(1),
$-\int\mathrm{KL}(P_r \,\|\, \bar\mu^\circ)\,\tilde\rho(dr)$, agrees with $\mathfrak{E}$ up to a
constant of the ladder where $\mathfrak{E}$ is finite. But it vanishes identically on a ladder whose
step is carried by one orbit $\{\tau q^i\}$. A second candidate disintegrates the slab of an octave
over the displacements instead of the ratios. It records, for each displacement $v$, the law of
the scale $y \in [1,q)$ at which $v$ appears, and compares it with $dy/(\ell y)$. The next
proposition shows that it selects one member on every ladder with no Gaussian part, the coupling
that spreads each displacement log-uniformly over the octave, and that this member is never the
similarity member.

% draft: none. New at the supplement review (2026-11-16): the entropy of the appearance profile,
% notes/phase-entropy-without-covariant-member.md section 5(b), which the blueprint left
% "conjectured, not proved here". The argument was a fixer's (notes/HANDOFF.md, "A candidate
% blueprint node"), checked and stated here. Found in checking it: the maximizer is unique only as
% a measure Lambda, so the uniqueness of the family needs the ladder to have no Gaussian part, as
% in thm:phase-entropy-selection(3); uniqueness comes from KL(kappa || eta) = 0 iff kappa = eta and
% the J-invariance of Lambda, not from strict convexity; and (4) holds on every ladder with a
% similarity member and a nonzero step measure, of which the Matern ladder is one. The slab is
% read on [1,q) x (0,infinity), so that eta is the Haar measure of the cylinder; by
% thm:surface-measure(iii) it carries the mass of (1,q] x (0,infinity).
% CHANGED (2026-11-16, the Lean proof, Formalization/Pyramids/LineAppendix/AppearanceEntropy.lean):
% every clause is machine-checked, so statement and proof are marked proved, the status line says
% what is checked, and the proof follows the checked route: (1) reads the slab as a measure on the
% cylinder with (\ref{eq:line-e3}), so that lem:phase-disintegration-exists's argument applies to it;
% (c) => (a) compares finite measures on the half-lines (-infinity, x]; the uniqueness of (3) goes
% through the data on [1,q) and lem:step-field-determines-family, in place of the J-invariance of
% Lambda; (4) shows mu([x,infinity)) = 0 by comparing two windows, in place of the affine function
% of log x. The statement is unchanged.
\begin{proposition}[the appearance entropy]
\label{prop:appearance-entropy}
\lean{Pyramids.appearance_entropy_1, Pyramids.appearance_entropy_2, Pyramids.appearance_entropy_3,
Pyramids.appearance_entropy_3_unique, Pyramids.appearance_entropy_4}
\leanok
\uses{def:modulation-orbit, thm:surface-measure, thm:annulus-parametrization,
prop:strata-cylinder-coordinates, thm:phase-entropy-selection}
Let $\Psi_1$ be a knot exponent with step $\phi$ and step measure $\mu$, and let $\eta$ be the Haar
measure of Proposition~\ref{prop:strata-cylinder-coordinates}, on $[1,q)$. Let a member of the
modulation orbit have the pair $(\Lambda, \mathcal{B})$. An \emph{appearance disintegration} of the
member is a family of probability measures $\kappa_v$ on $[1,q)$, $v > 0$, measurable in $v$, such
that for every Borel set $E \subseteq [1,q) \times (0,\infty)$
\[
  \Lambda(E) = \int\kappa_v\bigl(\{y : (y, v) \in E\}\bigr)\,\mu(dv) .
\]
The \emph{appearance entropy} of the member is
$\mathfrak{A} := -\int\mathrm{KL}(\kappa_v \,\|\, \eta)\,\mu(dv) \in [-\infty, 0]$.
\begin{enumerate}
\item Every member has an appearance disintegration, and any two agree for $\mu$-almost every $v$.
\item The path $\Psi_t := \Psi_1 + (\log_qt)\,\phi$, $t \in [1,q]$, satisfies the conditions
(1)--(4) of Theorem~\ref{thm:annulus-parametrization}. Its family, the \emph{log-uniform member},
lies in the modulation orbit, has $\varpi_t = \varpi_1 + (\log_qt)\,\mu$ for $t \in [1,q]$, and has
$\mathfrak{A} = 0$.
\item Every member has $\mathfrak{A} \le 0$, with equality iff
$\varpi_t = \varpi_1 + (\log_qt)\,\mu$ for every $t \in [1,q]$. If the ladder has no Gaussian
part, the log-uniform member is the only member with $\mathfrak{A} = 0$, so it is the unique
maximizer of the appearance entropy on the modulation orbit.
\item If $\mu \ne 0$ and the modulation orbit contains the similarity member, then the similarity
member has $\mathfrak{A} < 0$.
\end{enumerate}
\statusT\quad\emph{(Cites nothing of its own: the \Levy{}--Khintchine representation and
Khintchine's theorem enter only through
Theorem~\ref{thm:annulus-parametrization}, Theorem~\ref{thm:surface-measure} and
Lemma~\ref{lem:line-exponent-data}. Every clause is machine-checked. \textbf{Hypotheses} (the
node's, each with what its declarations require): $q > 1$: same; $\Psi_1$ a knot exponent: same in
\texttt{appearance\_entropy\_2}, absent from the other four, which read only membership of the
orbit; $\mu$ the step measure: same (\texttt{IsStepMeasure}); a member with its pair: same
(\texttt{modulationOrbit}, \texttt{IsFamilyPair}); an appearance disintegration: same
(\texttt{IsAppearanceDisintegration}), with $\kappa_v$ given for every real $v$ and the identity
read for every Borel set of the plane on the restriction of $\Lambda$ to $[1,q) \times (0,\infty)$;
(3) no Gaussian part: same (\texttt{HasNoGaussianStep}); (4) $\mu \ne 0$ and the similarity member:
same, a member with $\Psi_t = \Psi_1(t\,\cdot)$ for $t > 0$. The family of (2) is given through its
exponents on $[1,q]$, which fix it (Theorem~\ref{thm:annulus-parametrization}). The uniqueness in
(3) is stated in decomposed form, the reader composing it: \texttt{appearance\_entropy\_3\_unique}
gives that two members with $\mathfrak{A} = 0$ have the same kernels, and
\texttt{appearance\_entropy\_2} that the log-uniform member has $\mathfrak{A} = 0$. The appearance
entropy is computed in the canonical discrete gauge in which the member is written, with the knots
as the ends of the octave.)}
\end{proposition}

\begin{proof}
\leanok
\uses{lem:knot-exponent-ladder, thm:surface-measure, thm:knot-atom-cascades,
lem:line-exponent-data, thm:annulus-parametrization, lem:phase-disintegration-exists,
lem:kl-measurable, lem:member-gaussian-part, lem:step-field-determines-family,
thm:phase-entropy-selection, prop:strata-cylinder-coordinates}
Write $F := [1,q) \times (0,\infty)$, the cylinder $\mathcal{C}$ of
Theorem~\ref{thm:knot-atom-cascades} read in the coordinates $(y, v)$, and $\Lambda_F$ for the
restriction of $\Lambda$ to $F$. The step measure is the folded \Levy{} measure of $\phi$, so it is
null on $(-\infty,0]$, finite on every $(\varepsilon,\infty)$, $\varepsilon > 0$, and
$\sigma$-finite. The lines $\{y\} \times (0,\infty)$, $y > 0$, carry no mass of $\Lambda$
(Theorem~\ref{thm:surface-measure}(iii)), and neither does the half-plane $\RR \times (-\infty,0]$.
The sets $((-\infty,x] \times A) \cap F$ and $(1, \min(x,q)] \times A$ differ only by points of the
lines $y = 1$ and $y = q$ and of that half-plane, so for real $x$ and Borel $A$
\[
  \Lambda_F((-\infty,x] \times A) = \Lambda((1, \min(x,q)] \times A). \qquad (*)
\]
At $x = q$, with Lemma~\ref{lem:knot-exponent-ladder}(3), the second marginal of $\Lambda_F$ is
$\mu$. Also $\eta((-\infty,x]) = \log_q\min(x,q)$ for $x \ge 1$, by
$\int_1^td\theta/(\ell\theta) = \log_qt$, and $\eta((-\infty,x]) = 0$ for $x < 1$.

(1) $\Lambda_F$ is carried by $F$, and in the coordinates $(\theta, r) = (y, v)$ it satisfies
(\ref{eq:line-e3}). Indeed, $\Theta(vy) \le \max(q^2, 1 + \ell)\,\Theta(v)$ for $y \in [1,q)$ and
$v > 0$, with $\Theta$ the weight of (\ref{eq:line-e3}). And $\int\Theta(v)\,\Lambda_F(dy, dv)$ is
the integral (\ref{eq:line-e3}) of the cylinder measure $\sigma$ of $\Lambda$, since $\sigma$ is the
image of $\Lambda_F$ under $(y, v) \mapsto (y, v/y)$ and $r\theta = v$ there. It is finite by
Theorem~\ref{thm:knot-atom-cascades}(2) and Theorem~\ref{thm:surface-measure}(i). The last sentence
of Lemma~\ref{lem:phase-disintegration-exists}, applied to $\Lambda_F$, whose second marginal is
$\mu$, gives an appearance disintegration, and any two agree for $\mu$-almost every $v$. By
Lemma~\ref{lem:kl-measurable} the map $v \mapsto \mathrm{KL}(\kappa_v \,\|\, \eta)$ is measurable,
so $\mathfrak{A}$ is defined, and it does not depend on the disintegration chosen.

\emph{The log-uniform slab.} For a member the following are equivalent: (a) $\Lambda_F$ is the
product $\eta \otimes \mu$, with $y$ first; (b) $\mathfrak{A} = 0$; (c)
$\Lambda((1,t] \times \cdot\,) = (\log_qt)\,\mu$ for $t \in [1,q]$; (d)
$\varpi_t = \varpi_1 + (\log_qt)\,\mu$ for $t \in [1,q]$.

(a)$\Leftrightarrow$(b). The integrand of $\mathfrak{A}$ is measurable and nonnegative, so
$\mathfrak{A} = 0$ iff $\mathrm{KL}(\kappa_v \,\|\, \eta) = 0$ for $\mu$-almost every $v$, that is,
iff $\kappa_v = \eta$ for $\mu$-almost every $v$, both being probability measures. Then
$\Lambda_F(E) = \int\eta(\{y : (y,v) \in E\})\,\mu(dv) = (\eta \otimes \mu)(E)$ by Tonelli's
theorem. Conversely, if $\Lambda_F = \eta \otimes \mu$, the constant family $\eta$ is an appearance
disintegration, and $\kappa_v = \eta$ for $\mu$-almost every $v$ by (1).

(a)$\Rightarrow$(c). By $(*)$, $\Lambda((1,t] \times A) = \Lambda_F((-\infty,t] \times A)
= \eta((-\infty,t])\,\mu(A) = (\log_qt)\,\mu(A)$.

(c)$\Rightarrow$(a). Let $A \subseteq (\varepsilon,\infty)$ be Borel, $\varepsilon > 0$, so that
$\mu(A) < \infty$. The measures $S \mapsto \Lambda_F(S \times A)$ and $\mu(A)\,\eta$ on $\RR$ are
finite, the first being at most $\Lambda_F((-\infty,q] \times A) = \mu(A)$. By $(*)$ and (c) both
give $(-\infty,x]$ the mass $(\log_q\min(x,q))\,\mu(A)$ for $x \ge 1$, and $0$ for $x < 1$. Two
finite measures on $\RR$ that agree on every $(-\infty,x]$ are equal, so
$\Lambda_F(S \times A) = \eta(S)\,\mu(A)$ for Borel $S$. For a fixed $S$, the measures
$A \mapsto \Lambda_F(S \times A)$ and $\eta(S)\,\mu$ are null on $(-\infty,0]$ and agree on the
Borel sets bounded away from $0$, hence are equal, by continuity from below. So $\Lambda_F$ and
$\eta \otimes \mu$ agree on every rectangle, and as $\eta$ and $\mu$ are $\sigma$-finite,
$\Lambda_F = \eta \otimes \mu$.

(c)$\Leftrightarrow$(d). By (\ref{eq:line-e1}), $\varpi_t = \varpi_1 + \Lambda((1,t] \times \cdot\,)$, so (c)
gives (d). Conversely, (d) gives $\Lambda((1,t] \times A) = (\log_qt)\,\mu(A)$ for $A$ bounded away
from $0$, where $\varpi_1(A)$ is finite and cancels. Both sides are null on $(-\infty,0]$, so (c)
follows by continuity from below.

(2) Let $\Psi_1$ have the data $(a(1), \varpi_1)$ and $\phi$ the data $(a_\phi, \mu)$. For $c \ge 0$
the function $\Psi_1 + c\phi$ lies in $\mathrm{ND}_s$ with the data
$(a(1) + ca_\phi, \varpi_1 + c\mu)$, by linearity of the formula. Both $\Psi_1$ and $\phi$ are
nonnegative, and $\Psi_1 \not\equiv 0$, since otherwise $\phi = \Psi_1(q\,\cdot) - \Psi_1$ would
vanish. So $\Psi_t \ge \Psi_1$ is not identically zero: this is condition (1) of
Theorem~\ref{thm:annulus-parametrization}. For $1 \le t < t' \le q$ the difference
$\Psi_{t'} - \Psi_t = (\log_qt' - \log_qt)\,\phi$ is a positive multiple of $\phi$, which lies in
$\mathrm{ND}_s$ and is not identically zero: this is (2). Condition (3) is
$\Psi_q = \Psi_1 + \phi = \Psi_1(q\,\cdot)$, and (4) holds because $\log_qt$ is continuous. By
Theorem~\ref{thm:annulus-parametrization} the path is that of a family, whose exponent at
$t \in [1,q]$ is $\Psi_t$. At $t = 1$ it is $\Psi_1$, so the family is a member. By
(\ref{eq:line-e1}) and the uniqueness of the data (Lemma~\ref{lem:line-exponent-data}(1)),
$\varpi_t = \varpi_1 + (\log_qt)\,\mu$ for $t \in [1,q]$. This is (d), so $\mathfrak{A} = 0$.

(3) $\mathfrak{A} \le 0$ since the relative entropy is nonnegative, and the equality case is
(b)$\Leftrightarrow$(d). Let the ladder have no Gaussian part, and let two members with the pairs
$(\Lambda, \mathcal{B})$ and $(\Lambda', \mathcal{B}')$ have $\mathfrak{A} = 0$. By
Lemma~\ref{lem:member-gaussian-part}(2), $\mathcal{B}((0,t]) = 0 = \mathcal{B}'((0,t])$ for every
$t > 0$. The two members have the same $\varpi_1$, the measure of the data of $\Psi_1$, which are
unique (Lemma~\ref{lem:line-exponent-data}(1)), and so by (d) the same $\varpi_t$ for
$t \in [1,q]$. By (\ref{eq:line-e1}) their kernels at every $t \in [1,q)$ have the same data, hence the same
exponents, and a family in the canonical discrete gauge is determined by its exponents on $[1,q)$
(Lemma~\ref{lem:step-field-determines-family}). So the two members coincide, and by (2) the
log-uniform member is one of them.

(4) Let the similarity member, with $\Psi_t = \Psi_1(t\,\cdot)$ for $t > 0$, have
$\mathfrak{A} = 0$. It has $\Psi_{\lambda t} = \Psi_t(\lambda\,\cdot)$ for all $\lambda, t > 0$
(Theorem~\ref{thm:phase-entropy-selection}(1)), so by ``both parts'' of
Proposition~\ref{prop:strata-cylinder-coordinates}(1), $\varpi_t = T_t\varpi_1$ for $t > 0$. By (d),
$T_t\varpi_1 = \varpi_1 + (\log_qt)\,\mu$ for $t \in [1,q]$. For $x > 0$ put
$G(x) := \varpi_1([x,\infty))$ and $M(x) := \mu([x,\infty))$, finite since $\varpi_1$ and $\mu$ are
\Levy{} measures. Evaluated at $[x,\infty)$, the identity reads
\[
  G(x/t) = G(x) + (\log_qt)\,M(x) \qquad (x > 0,\ t \in [1,q]),
\]
and at $t = q$ it is $G(x/q) = G(x) + M(x)$. Let $s \in [1,q)$. The identity at $x/s$ with $q/s$, at
$x$ with $s$, and at $x$ with $q$ give
$G(x) + (\log_qs)\,M(x) + (1 - \log_qs)\,M(x/s) = G(x/q) = G(x) + M(x)$, and $\log_qs < 1$, so
$M(x/s) = M(x)$. With $s = \sqrt q$ twice, $M(xq) = M(x)$. So
$G(xq^k) = G(xq^{k+1}) + M(x)$ for every $k \ge 0$, and $G(x) = G(xq^n) + nM(x) \ge nM(x)$ for every
$n$. As $G(x)$ is finite, $M(x) = 0$. So $\mu([x,\infty)) = 0$ for every $x > 0$, and $\mu = 0$,
against $\mu \ne 0$. So the similarity member has $\mathfrak{A} < 0$.

Every step is machine-checked, and the cited facts enter only through the statements the proof
cites.
\end{proof}

On a ladder with no Gaussian part the step measure is not zero, since the step is not identically
zero. So on every such ladder the appearance entropy selects the log-uniform member, and where the
ladder has a similarity member it selects another one. The \Matern{} ladder of
Theorem~\ref{thm:matern-pyramid-line} is such a ladder. With a Gaussian part, $\mathfrak{A}$ does
not see $\mathcal{B}$, and uniqueness of the family is not claimed. Unlike $\mathfrak{E}$, which a
change of the scale unit leaves unchanged, the appearance entropy reads the knots as the ends of the
octave. The two functionals fail in complementary places. The phase entropy returns the similarity
member and selects nothing on a ladder with a singular step. The appearance entropy selects on every
ladder with no Gaussian part and never returns the similarity member. Within the discrete stratum, stratum (2) of
Theorem~\ref{thm:lattice-line}, the selection of chapter~4 does not use the phase entropy
(Theorem~\ref{thm:selection-stratum-two}).

\section{Three example clusters}
\label{sec:line-examples}

Three kinds of example place the classical pyramids against the axioms: the finite masks, which are
steps of no family; the \Matern{} pyramid, a fully covariant family whose octave steps are single
Laplace kernels; and what the classical pyramids converge to.

% draft: Proposition 2.16 (section 2.5), lines 217-219 of draft/spatial-hemigroup-pyramids.md.
% CHANGED: the header states the scope ("no family continuous in scale"), and the "in particular"
% clause is split into the list of masks and the statement about pyramids built from them (review
% flags on Prop. 2.16's header and its two claims); the claim on the equivalent weights, which
% holds of every mask, is not made here (Proposition 3.15(1)).
% CHANGED: "and to every finite mask" narrowed to the finite masks that are probability measures and
% not point masses, the hypothesis of the first sentence (blind review 2, flag 14).
\begin{proposition}[the classical generating kernels are steps of no family]
\label{prop:classical-kernels-no-exact-cascade}
% CHANGED (fidelity review, 2026-11-16, R50): classical_kernels_no_exact_cascade_burt added, the
% clause on Burt's family (the author's decision).
\lean{Pyramids.classical_kernels_no_exact_cascade,
Pyramids.classical_kernels_no_exact_cascade_masks,
Pyramids.classical_kernels_no_exact_cascade_burt,
Pyramids.classical_kernels_no_exact_cascade_pyramid}
\leanok
\uses{lem:line-axioms-kernels}
Let $w$ be a probability measure on $\RR$ with bounded support, not a point mass, for instance a
mask of finite support on a grid $h\ZZ$. Then $w$ is not infinitely divisible, and it is neither a
kernel $\mu_{0,t}$ nor an increment $\mu_{s,t}$ of any family of operators satisfying (A1)--(A7),
covariant or not. This applies to the binomial mask $\tfrac1{16}[1,4,6,4,1]$ on $h\ZZ$, to every
five-tap mask of Burt's family that is a probability measure, to the B-spline kernels of every
order, and to every finite mask that is a probability measure and not a point mass. For each such
$w$ and each $q > 1$, no family satisfying
(A1)--(A7) has the octave steps $\mu_{q^n,q^{n+1}} = T_{q^n}w$, $n \in \ZZ$, of the pyramid built
from $w$.
% CHANGED (fidelity review, 2026-11-15, R26/R50/R69): the status line lists the declarations' hypotheses beside the node's (LINKAGE rule 6).
% CHANGED (fidelity review, 2026-11-16, R50): Burt's family reads "same", through its own declaration.
\statusT\quad\emph{(\textbf{Hypotheses} (the node's, each with what its declarations require): $w$ a
probability measure on $\RR$: same; bounded support: same, $w(\{|x| > R\}) = 0$ for some $R$; not
a point mass: same; a family of operators with (A1)--(A7): same, a \texttt{PreCascadeCore} with
\texttt{IsPositive}, read through its kernels (\texttt{IsKernelFamily}); the binomial mask on
$h\ZZ$: same, $h > 0$; B-spline kernels of every order: same, order $m \ge 1$ and $h > 0$; a
finite mask that is a probability measure and not a point mass: same, finitely supported,
nonnegative, of sum $1$; a five-tap mask of Burt's family that is a probability measure: same,
\texttt{classical\_kernels\_no\_exact\_cascade\_burt} requires $h > 0$, $a, c \ge 0$ and Burt's
normalization and equal contribution, $b \ge 0$ following from them, and not his unimodality
(weaker); (last sentence)
$q > 1$: same.)}
\end{proposition}

% CHANGED (wave 2e, 2026-11-14): the proof follows the machine-checked route. It says how the
% proposition on R^d is read on the line (through the embedding x -> x e_1), why the listed masks
% have bounded support and are no point masses, and that the last sentence needs only the step at
% n = 0; a closing paragraph says which steps are checked.
% CHANGED (the theory paper, 2026-11-16): the closing sentence on checking and cited facts shortened; no step changed.
\begin{proof}
\leanok
\uses{lem:line-axioms-kernels, prop:no-bounded-step}
By Lemma~\ref{lem:line-axioms-kernels}(1) the family acts by convolution with probability
measures $\mu_{s,t}$ satisfying (I1) and (I2). Under the embedding $x \mapsto xe_1$ of $\RR$ in
$\RR^1$, the image of $w$ is a probability measure of bounded support and not a point mass, the
images of the $\mu_{s,t}$ satisfy (I1) and (I2), and the embedding is injective on laws and
carries infinite divisibility. So Proposition~\ref{prop:no-bounded-step} at $d = 1$ says that
$w$ is not infinitely divisible and is no $\mu_{s,t}$ of such a family.

The masks listed are probability measures of bounded support: a finite mask on $h\ZZ$ is carried
by finitely many points, and the B-spline kernel of order $m$, the $m$-fold convolution of the
uniform law on $[-h/2, h/2]$, by $[-mh/2, mh/2]$. None is a point mass, since the transform of a
point mass has modulus $1$ everywhere, while the transform $\cos^4(h\omega/2)$ of the binomial
mask vanishes at $\omega = \pi/h$, and the transform $\operatorname{sinc}(h\omega/2)^m$ of the
B-spline kernel, $m \ge 1$, vanishes at $\omega = 2\pi/h$.
% CHANGED (third blind pass, Q-0767, 2026-11-15): the proof says why a mask of Burt's family is not
% a point mass; the statement names the family without that hypothesis.
% CHANGED (fidelity review, 2026-11-16, R50): b = 1/4 derived from Burt's two constraints, as the
% checked declaration classical_kernels_no_exact_cascade_burt does, in place of read from the page.
A five-tap mask $(c, b, a, b, c)$ on $h\ZZ$ of Burt's family satisfies his normalization
$a + 2b + 2c = 1$ and his equal contribution at the ratio $2$, $a + 2c = 2b$
[\sscite{burt1981fast}, pp.~23--24]. Together they give $4b = 1$, so the mask puts mass
$\tfrac14$ at each of the two points $\pm h$ and is not a point mass. For the other finite
masks it is the hypothesis. For the last sentence, the octave step at $n = 0$ is $\mu_{1,q} = T_1w = w$, which the
first sentence excludes.

% CHANGED (fix round of the third blind pass, Q-0784, 2026-11-15): the closing sentence no longer
% says every step is machine-checked and no fact is cited; the reading of Burt's family is neither.
% CHANGED (fidelity review, 2026-11-16, R50): every step is now checked; Burt's page supplies only
% the definition of the family.
% CHANGED (supplement review, 2026-11-16): the sentence on what Burt's page supplies is dropped from
% the proof of record; the citation accounting is the status line's.
Every step is machine-checked, and the cited facts enter only through
Proposition~\ref{prop:no-bounded-step}.
\end{proof}

What the classical constraints capture can be stated. Normalization is (A5), symmetry is (A3), and
Burt's equal contribution, that every fine node contributes the same total weight to the coarse
level [\sscite{burt1983laplacian}, p.~533], is the condition that the subsampled convolution preserve
unit mass on the grid. Separability is the product structure of the frame families of
\prog{affine v0.1}{Prop.~F.1, p.~93}. None of the four constrains the family across levels. What
Burt and Adelson observed, that an equivalent weighting function resembles the Gaussian
[\sscite{burt1983laplacian}, pp.~533--534], is made exact in Proposition~\ref{prop:classical-pyramids-approximate}:
the limit is the cubic B-spline.

The \Matern{} pyramid and its ladder are built from Laplace atoms $\log(1 + \omega^2/\vartheta^2)$.
The next lemma says when a sum of them is admissible.

% draft: none. New at the transcription: the Frullani-type identity of the draft's proof of Theorem
% 2.17(1) (line 227), which the proofs of Theorem 2.17(3) and Remark 2.17' also use, stated once as a
% lemma with the admissibility of the sums (blind review 2, flags 5, 6 and 18).
\begin{lemma}[sums of Laplace atoms are admissible]
\label{lem:laplace-atoms-admissible}
\lean{Pyramids.laplace_atoms_admissible_1, Pyramids.laplace_atoms_admissible_2}\leanok
Let $a \ge 0$, and let $(\vartheta_i)_{i \in I}$ be positive numbers over a countable index set
with $\sum_i\vartheta_i^{-2} < \infty$.
\begin{enumerate}
\item For $\vartheta > 0$ and $\omega \in \RR$,
\[
  \int_0^\infty\bigl(1 - \cos\omega x\bigr)\,e^{-\vartheta x}\,\frac{dx}{x}
  = \frac12\log\Bigl(1 + \frac{\omega^2}{\vartheta^2}\Bigr).
\]
\item The series $F(\omega) := a\,\omega^2 + \sum_i\log(1 + \omega^2/\vartheta_i^2)$ converges, with
$0 \le F(\omega) \le (a + \sum_i\vartheta_i^{-2})\,\omega^2$. $F$ is admissible, of the form
(\ref{eq:admissible-exponent}) with Gaussian coefficient $a$ and the completely monotone profile
% CHANGED (2026-11-23, review round 3 batch C, supp.S10): the constant 2, which the proof gives, in place of 4.
$k(x) = 2\sum_ie^{-\vartheta_ix}$, and $\int(1 \wedge x^2)\,k(x)\,dx/x \le 2\sum_i\vartheta_i^{-2}$.
$F \not\equiv 0$ unless $a = 0$ and $I$ is empty.
\end{enumerate}
% CHANGED (fidelity review, 2026-11-15, R26/R50/R69): the status line lists the declarations' hypotheses beside the node's (LINKAGE rule 6).
\statusT\quad\emph{(Cites nothing: differentiation under the integral sign, Tonelli's theorem and
term-by-term differentiation of a series. \textbf{Hypotheses} (the node's, each with what its
declarations require): $a \ge 0$, the $\vartheta_i$ and $\sum_i\vartheta_i^{-2} < \infty$: absent from
\texttt{laplace\_atoms\_admissible\_1}, which reads one $\vartheta$; (1) $\vartheta > 0$, $\omega$
real: same; (2) $a \ge 0$: same; $I$ countable: same; $\vartheta_i > 0$: same;
$\sum_i\vartheta_i^{-2} < \infty$: same; $F$ the series: same, $F$ given as the sum with no
convergence assumed, which the declaration proves.)}
\end{lemma}

% CHANGED (wave 1d, 2026-11-13): the checked route. In (1) the sine integral is computed as the imaginary part of an exponential integral, and the equality of the two sides follows from equal derivatives and equal values at 0; in (2) the bound on the derivatives is made explicit, and the integral bound uses (1 wedge x^2)/x <= x on the whole half-line, which gives 2 sum theta_i^{-2}, in place of the bounds on (0,1) and (1,infinity) taken separately.
\begin{proof}
\leanok
(1) For $\vartheta > 0$ and $x > 0$, $1 - \cos y \le |y|$ gives
$0 \le (1 - \cos\omega x)e^{-\vartheta x}/x \le |\omega|e^{-\vartheta x}$, so the left side is an
absolutely convergent integral, and both sides vanish at $\omega = 0$. The derivative of the
integrand in $\omega$ is $\sin(\omega x)e^{-\vartheta x}$, bounded by the integrable function
$e^{-\vartheta x}$, so the left side may be differentiated under the integral sign. Its derivative
is the imaginary part of
$\int_0^\infty e^{(-\vartheta + \sqrt{-1}\,\omega)x}\,dx = (\vartheta - \sqrt{-1}\,\omega)^{-1}
= (\vartheta + \sqrt{-1}\,\omega)/(\vartheta^2 + \omega^2)$, that is $\omega/(\vartheta^2 + \omega^2)$.
This is also the derivative of the right side. The two sides have the same derivative at every
$\omega$ and the same value at $0$, so they are equal.

(2) $0 \le \log(1 + y) \le y$ for $y \ge 0$ gives the convergence and the bound. For $y \ge 0$ and
$n \ge 0$, $y^{n+2}/(n+2)! \le e^y$, so for $\vartheta, x > 0$
\[
  \vartheta^ne^{-\vartheta x} \le \frac{(n+2)!}{x^{n+2}}\,\vartheta^{-2} .
\]
With $n = 0$ the series defining $k(x)$ converges for every $x > 0$. For $x$ in $(x_0, \infty)$,
$x_0 > 0$, the bound at $x_0$ dominates the $n$-th derivatives $(-\vartheta_i)^ne^{-\vartheta_ix}$
by a summable sequence, so the series $2\sum_i(-\vartheta_i)^ne^{-\vartheta_ix}$ may be
differentiated term by term. Hence $k$ is smooth on $(0, \infty)$ with
$k^{(n)}(x) = 2\sum_i(-\vartheta_i)^ne^{-\vartheta_ix}$, and $(-1)^nk^{(n)}(x) =
2\sum_i\vartheta_i^ne^{-\vartheta_ix} \ge 0$: $k$ is completely monotone, and in particular
nonincreasing. For the integral, $(1 \wedge x^2)/x \le x$ for $x > 0$: on $(0, 1]$ the two are
equal, and on $(1, \infty)$ the left side is $1/x$. So for each $\vartheta > 0$
\[
  \int_0^\infty(1 \wedge x^2)\,e^{-\vartheta x}\,\frac{dx}{x}
  \le \int_0^\infty x\,e^{-\vartheta x}\,dx = \vartheta^{-2},
\]
and the parts of the left side over $(0, 1)$ and over $(1, \infty)$ are each at most
$\vartheta^{-2}$. Summing by Tonelli's theorem,
% CHANGED (2026-11-23, review round 3 batch C, supp.S10): the weakening to 4 dropped.
$\int(1 \wedge x^2)\,k(x)\,dx/x \le 2\sum_i\vartheta_i^{-2}$. In particular
$\int_0^1x\,k(x)\,dx$ and $\int_1^\infty k(x)\,dx/x$ are finite. By (1) and Tonelli's theorem, all
terms being nonnegative,
\[
  \int_0^\infty(1 - \cos\omega x)\,k(x)\,\frac{dx}{x}
  = \sum_i2\int_0^\infty(1 - \cos\omega x)\,e^{-\vartheta_ix}\,\frac{dx}{x}
  = \sum_i\log\Bigl(1 + \frac{\omega^2}{\vartheta_i^2}\Bigr),
\]
so $F(\omega) = a\omega^2 + \int_0^\infty(1 - \cos\omega x)\,k(x)\,dx/x$, of the form
(\ref{eq:admissible-exponent}). If $a = 0$ and $I$ is empty, $F \equiv 0$. If $a > 0$, then
$F(1) \ge a > 0$, and if $I$ contains an index $i$, then $F(1) \ge \log(1 + \vartheta_i^{-2}) > 0$.
\end{proof}

% draft: Theorem 2.17 (section 2.5), lines 223-237 of draft/spatial-hemigroup-pyramids.md;
% clauses (1), (2), (4), (5) here, clause (3) in thm:matern-pyramid-line-noncreation.
% CHANGED: split from Theorem 2.17 by grade: clause (3) rests on V-selection and Karlin.
% CHANGED: clause (1) proves admissibility directly from the Frullani-type identity, with the
% profile; the naming as the symmetric Thorin member of [V-cone v0.1, Prop. 5.7] and the bridge to
% Lindeberg's time-causal limit kernel move to the prose after the theorem (review flag,
% interpretation inside a statement), with the bridge at Theorem 4.1.
% CHANGED: clause (2) states the similarity member as in thm:phase-entropy-selection.
% CHANGED: clause (4) states the symbol of the generator and its size, and the first-order
% realization through prop:recursion-helmholtz; the citation of the Thorin clause of the generator
% proposition [V-cone v0.1, Prop. 6.4(4)] and of [Lindeberg 2016] moves to the prose.
% CHANGED: clause (5) drops "the bridge image of the causal dilogarithm family [II, Prop. 5.5]"
% (Paper II is unreleased and never cited), and states the admissibility of the limit directly.
\begin{theorem}[the \Matern{} pyramid]
\label{thm:matern-pyramid-line}
\lean{Pyramids.matern_pyramid_line_display, Pyramids.matern_pyramid_line_1,
Pyramids.matern_pyramid_line_2, Pyramids.matern_pyramid_line_2_direct,
Pyramids.matern_pyramid_line_3, Pyramids.recursion_line_transforms, Pyramids.recursion_line_passes,
Pyramids.matern_pyramid_line_4_bound, Pyramids.matern_pyramid_line_4_limit,
Pyramids.matern_pyramid_line_4_laws}
\leanok
\uses{thm:cited-A4, thm:phase-entropy-selection, prop:recursion-helmholtz}
Fix $q > 1$ and $c > 0$, and let
\begin{equation}\label{eq:line-e7}
  F_q(\omega) := \sum_{i \ge 1}\log\bigl(1 + c^2q^{-2i}\omega^2\bigr), \qquad
  e^{-F_q(t\omega)} = \prod_{i \ge 1}\frac{1}{1 + c^2q^{-2i}t^2\omega^2}.
\end{equation}
\begin{enumerate}
\item $F_q$ is admissible, with $a = 0$ and the completely monotone profile
$k(x) = 2\sum_{i \ge 1}e^{-xq^i/c}$. Its similarity family $\Psi_t = F_q(t\,\cdot)$ satisfies
(A1)--(A8) and (ND), and its octave steps are single \Matern{} steps:
$F_q(qt\omega) - F_q(t\omega) = \log(1 + c^2t^2\omega^2)$, so $\mu_{t,qt}$ is the Laplace law
$L_{ct}$ of range $ct$, where $L_r$ denotes the law with the transform
$\widehat L_r(\omega) = (1 + r^2\omega^2)^{-1}$.
% CHANGED (fourth blind pass, 2026-11-15): "for all $\lambda, t$" reads "for all $\lambda, t > 0$",
% the range on which Psi_{lambda t} is defined and the one the Lean quantifies over.
\item (Selection.) On the modulation orbit of its ladder the similarity family is the only member
with $\Psi_{\lambda t} = \Psi_t(\lambda\,\cdot)$ for all $\lambda, t > 0$, and the unique maximizer of
the phase entropy. Directly: if $\tilde F$ is a function on $(0,\infty)$ with
$\tilde F(q\omega) - \tilde F(\omega) = \log(1 + c^2\omega^2)$ and $\tilde F(0{+}) = 0$, then
$\tilde F = F_q$ on $(0,\infty)$.
\item (Generator, realization.) For every $\omega$,
$\partial_te^{-F_q(t\omega)} = h_t(\omega)\,e^{-F_q(t\omega)}$ with
\[
  h_t(\omega) = \frac2t\sum_{i \ge 1}\Bigl(\widehat L_{ctq^{-i}}(\omega) - 1\Bigr).
\]
% CHANGED (fourth blind pass, 2026-11-15): the asymptotic "$-h_t(\omega) =
% \tfrac2t\log_q(ct|\omega|) + O(1/t)$ as $|\omega| \to \infty$" is stated as the uniform bound that
% the proof gives and matern_pyramid_line_3 asserts.
Each term lies in $[-2/t, 0]$, and there is a constant $C$ with
$\bigl|-h_t(\omega) - \tfrac2t\log_q(ct|\omega|)\bigr| \le C/t$ for all $t > 0$ and all $\omega$ with
$ct|\omega| \ge 1$. In particular $h_t$ is unbounded. The convolution with each octave step
$L_{ct}$ is one forward and one backward first-order pass
(Proposition~\ref{prop:recursion-helmholtz}(1)).
\item (The limit $q \downarrow 1$.) Let
$F_\infty(\omega) := \int_0^\infty\log(1 + c^2\omega^2e^{-2u})\,du$. Then
\[
  0 \le F_\infty(\omega) - \ell F_q(\omega) \le \ell\,\log(1 + c^2\omega^2),
\]
so $\ell F_q \to F_\infty$ locally uniformly as $q \downarrow 1$, and
\begin{equation}\label{eq:line-e8}
  F_\infty(\omega) = -\tfrac12\,\mathrm{Li}_2\bigl(-c^2\omega^2\bigr)
  = \int_{1/c}^\infty\log\Bigl(1 + \frac{\omega^2}{\theta^2}\Bigr)\frac{d\theta}{\theta},
\end{equation}
with the dilogarithm $\mathrm{Li}_2(z) := -\int_0^z\log(1 - s)\,ds/s$ for $z \le 0$.
$F_\infty$ is admissible, with $a = 0$ and profile $k_\infty(x) = 2\int_{1/c}^\infty e^{-\theta x}\,d\theta/\theta$.
For each $t > 0$ the laws with exponent $\ell F_q(t\,\cdot)$ converge weakly to the law with exponent
$F_\infty(t\,\cdot)$, while the kernels $\mu_t$, with transforms $e^{-F_q(t\omega)}$, have no weak
limit.
\end{enumerate}
% CHANGED (wave 3e, 2026-11-15): every clause is machine-checked, so statement and proof are marked
% proved; the status line says what the citation carries on the checked route and what the other
% clauses read.
% CHANGED (fidelity review, 2026-11-15, R26/R50/R69): the status line lists the declarations' hypotheses beside the node's (LINKAGE rule 6).
\statusA\quad\ledger{A4}\quad\emph{(\textbf{Assignment.} \ledger{A4}, the classification on the
line (Theorem~\ref{thm:cited-A4}), carries one sentence of (1): an admissible $F \not\equiv 0$ has a similarity
family satisfying (A1)--(A8) and (ND). On the checked route it is the construction of that
theorem that is read: laws with the transforms of the gauge $\chi = \mathrm{id}$ and of $F_q$
exist, and any such laws are the kernels of a family with (A1)--(A8) and (ND) whose bijections are
$S_\lambda t = \lambda t$. The admissibility of $F_q$ and $F_\infty$, the profiles, the
telescoping, (2), (3) and (4) are proved here. (2) uses
Theorem~\ref{thm:phase-entropy-selection} and (3) uses
Proposition~\ref{prop:recursion-helmholtz}, both proved here; the laws of (4) exist by
Lemma~\ref{lem:line-exponent-data}(1), and their convergence is \Levy's continuity theorem, which
Mathlib proves (\texttt{ProbabilityMeasure.tendsto\_of\_tendsto\_charFun}). Every clause is
machine-checked, Theorem~\ref{thm:cited-A4} being itself a theorem from \ledger{A39}; the display,
(2) directly, (3) and the bound and the limit of (4) use no cited fact. The naming of $F_q$ as a
symmetric Thorin exponent and its bridge to the time-causal limit kernel are not part of the
statement. \textbf{Hypotheses} (the node's, each with what its declarations require): $q > 1$: same,
absent from \texttt{recursion\_line\_transforms}, \texttt{recursion\_line\_passes} and
\texttt{matern\_pyramid\_line\_4\_limit}, which has the limit $q \downarrow 1$, and in
\texttt{\_4\_laws} required of each $q$; $c > 0$: same, the two recursion declarations reading the
range $ct > 0$ as their $t > 0$; (1) the similarity family: same, any laws with the transforms
$e^{-(F_q(t\omega) - F_q(s\omega))}$ (\texttt{IsMaternPyramid}); (2) $\tilde F$ on $(0,\infty)$ with
the functional equation and $\tilde F(0{+}) = 0$: same; (3) $t > 0$: same; the passes:
stronger, \texttt{recursion\_line\_passes} requires the signal integrable, a class the sentence of (3)
does not name and Proposition~\ref{prop:recursion-helmholtz}(1), which it cites, states; (4)
$t > 0$: same; the kernels $\mu_t$: same, of any \Matern{} pyramid for each $q > 1$.)}
\end{theorem}

% CHANGED (wave 3e, 2026-11-15): the proof follows the machine-checked route. (1) says what the
% construction of thm:cited-A4 gives (any laws with the transforms are the kernels of the family,
% with S_lambda t = lambda t) and where the transform of L_r comes from; (2) checks that a Matern
% pyramid is a member of the orbit and that F_q is a knot exponent, and argues F_q(0+) = 0; (3)
% replaces "locally uniformly" by a summable bound on (0, 2t) and makes the constant explicit;
% (4) gives the series of integrals behind the bound, the existence of the limit law, and the
% kernels' clause from the definition of weak convergence and the bound alone.
% CHANGED (fourth blind pass, 2026-11-15): (1) derives the bijections S_lambda t = lambda t from the
% transforms and lem:line-axioms-kernels(3), where it had read S_lambda = chi^{-1}(lambda chi) off
% the construction, which thm:cited-A4 does not print; (2) defines a Matern pyramid, writes T_q for
% the dilation on the line, and reads the Gaussian coefficient of the step by the uniqueness of the
% data.
\begin{proof}
\leanok
\uses{thm:cited-A4, thm:phase-entropy-selection, prop:recursion-helmholtz,
lem:laplace-atoms-admissible, lem:line-axioms-kernels, lem:line-exponent-data}
(1) With $\vartheta_i := q^i/c$, $\log(1 + c^2q^{-2i}\omega^2) = \log(1 + \omega^2/\vartheta_i^2)$,
and $\sum_{i \ge 1}\vartheta_i^{-2} = c^2/(q^2 - 1) < \infty$. So
Lemma~\ref{lem:laplace-atoms-admissible}(2) with $a = 0$ makes the series converge, with
$0 \le F_q(\omega) \le c^2\omega^2/(q^2 - 1)$, and makes $F_q$ admissible, with the completely
monotone profile $k = 2\sum_ie^{-\vartheta_ix}$, and $F_q \not\equiv 0$. The display (\ref{eq:line-e7}) is
the exponential of the series. Telescoping:
$F_q(q\omega) = \sum_{i \ge 1}\log(1 + c^2q^{-2(i-1)}\omega^2) = \log(1 + c^2\omega^2) + F_q(\omega)$.
Take the gauge $\chi = \mathrm{id}$, so that the laws $\mu_{s,t}$, $0 \le s \le t$, are to have the
transforms $e^{-(F_q(t\omega) - F_q(s\omega))}$. By the converse of Theorem~\ref{thm:cited-A4}
(\prog{V-line v0.1}{Thm.~7.3, p.~37}) such laws exist, and the operators $f \mapsto \mu_{s,t} * f$
satisfy (A1)--(A8) and (ND); any laws with these transforms are these, a law being determined by its
transform, and they are symmetric, their transforms being real. The bijections of (A8) are
$S_\lambda t = \lambda t$: for $\lambda > 0$,
$\widehat\mu_{s,t}(\lambda\omega) = e^{-(F_q(\lambda t\omega) - F_q(\lambda s\omega))} =
\widehat\mu_{\lambda s,\lambda t}(\omega)$, so $T_\lambda\mu_{s,t} = \mu_{\lambda s,\lambda t}$, and
Lemma~\ref{lem:line-axioms-kernels}(3) turns this into (A8) at $\lambda$ with
$S_\lambda t = \lambda t$. At $t\omega$ the
telescoping is the octave step: $\mu_{t,qt}$ has the transform $(1 + c^2t^2\omega^2)^{-1}$, which
is $\widehat L_{ct}$ by Proposition~\ref{prop:recursion-helmholtz}(1), and a law is determined by
its transform.

(2) Call a \emph{\Matern{} pyramid} a family of laws $\mu_{s,t}$ with the transforms
$e^{-(F_q(t\omega) - F_q(s\omega))}$ of (1), the kernels of the similarity family. Let $\mu$ be one.
By (1) and Lemma~\ref{lem:line-axioms-kernels}(1) its laws
form a symmetric family; $T_q\mu_{s,t} = \mu_{qs,qt}$, the two laws having the transform
$e^{-(F_q(qt\omega) - F_q(qs\omega))}$; and $\mu_{0,1}$ has the exponent $F_q$. So $\mu$ is a member
of the modulation orbit of the ladder of $F_q$, and it is the similarity member, its exponents
being $\Psi_t(\omega) = F_q(t\omega)$. $F_q$ is a knot exponent: $F_q$ and its step
$\log(1 + c^2\omega^2)$ lie in $\mathrm{ND}_s$, and the step is not identically zero. The ladder
has no Gaussian part, its step being the exponent of the single Laplace atom $\vartheta = 1/c$:
Lemma~\ref{lem:laplace-atoms-admissible}(2) writes it with the Gaussian coefficient $0$, and by the
uniqueness of the data (Lemma~\ref{lem:line-exponent-data}(1)) that is its Gaussian coefficient.
So clauses (1)
to (3) of Theorem~\ref{thm:phase-entropy-selection} apply. A member with
$\Psi_{\lambda t} = \Psi_t(\lambda\,\cdot)$ for all $\lambda, t > 0$ has the laws of the similarity
member, so it is a \Matern{} pyramid; a \Matern{} pyramid has that covariance. Every member has
$\mathfrak{E} \le 0$; the similarity member has a pair and a disintegration with
$\mathfrak{E} = 0$; and a member has $\mathfrak{E} = 0$ iff its laws are those of the similarity
member, that is iff it is a \Matern{} pyramid. For the direct statement, $D := \tilde F - F_q$ has
$D(q\omega) = D(\omega)$ for $\omega > 0$, by the telescoping. The bound of (1) gives
$F_q(0{+}) = 0$, so $D(0{+}) = 0$, and $D(\omega) = D(q^{-n}\omega) \to 0$ as $n \to \infty$.

(3) The term $\log(1 + c^2q^{-2i}\tau^2\omega^2)$ has the $\tau$-derivative
$2c^2q^{-2i}\tau\omega^2/(1 + c^2q^{-2i}\tau^2\omega^2)$, at most $4c^2q^{-2i}t\omega^2$ for
$\tau \in (0, 2t)$, a summable bound. So $\partial_tF_q(t\omega)$ is the sum of the derivatives,
and $-x/(1 + x) = (1 + x)^{-1} - 1$ with $x = c^2q^{-2i}t^2\omega^2$ gives
$\partial_tF_q(t\,\cdot) = -h_t$ in the stated form; the chain rule gives
$\partial_te^{-F_q(t\omega)} = h_t(\omega)e^{-F_q(t\omega)}$. Each term $\widehat L_r - 1$ lies
in $[-1, 0]$. Let $ct|\omega| \ge 1$, and put $X := ct|\omega|$, $x_i := Xq^{-i}$ and
$N := \lfloor\log_qX\rfloor$, so that $q^N \le X < q^{N+1}$. Then
$-\tfrac t2h_t(\omega) = \sum_{i \ge 1}x_i^2/(1 + x_i^2)$. For $1 \le i \le N$,
$x_i \ge q^{N-i} \ge 1$, and the term differs from $1$ by $(1 + x_i^2)^{-1} \le q^{-2(N-i)}$; these
differences sum to at most $\sum_{j \ge 0}q^{-2j} = q^2/(q^2 - 1)$. For $i > N$,
$x_i < q^{N+1-i}$, and the terms, each at most $x_i^2 < q^{-2(i-N-1)}$, sum to at most
$q^2/(q^2 - 1)$. As $0 \le \log_qX - N < 1$,
\[
  \Bigl|-h_t(\omega) - \frac2t\log_q(ct|\omega|)\Bigr| \le \frac2t\Bigl(1 + \frac{q^2}{q^2 - 1}\Bigr)
\]
for every $t > 0$ and every $\omega$ with $ct|\omega| \ge 1$. At $\omega = q^K/(ct)$, $K \ge 0$ an
integer, this gives $-h_t(\omega) \ge (2K - 2 - 2q^2/(q^2 - 1))/t$, so $h_t$ is unbounded. The
realization is Proposition~\ref{prop:recursion-helmholtz}(1) at the range $ct$:
$L_{ct} = e^+_{ct} * e^-_{ct}$, and convolution with $e^\pm_{ct}$ solves $(1 \pm ct\,\partial_x)w = y$.

(4) The function $f(u) := \log(1 + c^2\omega^2e^{-2u})$ is nonnegative, decreasing and continuous,
and integrable on $[0,\infty)$, being at most $c^2\omega^2e^{-2u}$. With $\ell := \log q$,
$c^2q^{-2i}\omega^2 = c^2\omega^2e^{-2i\ell}$, so $\ell F_q(\omega) = \sum_{i \ge 1}\ell f(i\ell)$;
and $F_\infty(\omega) = \sum_{i \ge 0}\int_{i\ell}^{(i+1)\ell}f$, the partial sums being
$\int_0^{N\ell}f \to \int_0^\infty f$. As $f$ decreases,
$\int_{i\ell}^{(i+1)\ell}f \le \ell f(i\ell) \le \int_{(i-1)\ell}^{i\ell}f$. Summing over $i \ge 1$,
$\int_\ell^\infty f \le \ell F_q(\omega) \le \int_0^\infty f$, so
$0 \le F_\infty - \ell F_q \le \int_0^\ell f \le \ell f(0)$. As $\log(1 + c^2\omega^2)$ is bounded
on bounded sets of $\omega$ and $\ell \to 0$ as $q \downarrow 1$, $\ell F_q \to F_\infty$ locally
uniformly. For $\omega \ne 0$ the substitution $v = c^2\omega^2e^{-2u}$, which maps $(0,\infty)$
onto $(0, c^2\omega^2)$, gives
\[
  F_\infty(\omega) = \frac12\int_0^{c^2\omega^2}\log(1 + v)\,\frac{dv}{v}
  = -\frac12\,\mathrm{Li}_2(-c^2\omega^2),
\]
with $\mathrm{Li}_2(z) = -\int_0^z\log(1 - s)\,ds/s$ and $s = -v$; at $\omega = 0$ the three
expressions vanish. The substitution $\theta = e^u/c$, which maps $(0,\infty)$ onto
$(1/c,\infty)$, gives the last form of (\ref{eq:line-e8}). The profile $k_\infty$ is nonnegative and
nonincreasing, $e^{-\theta x}/\theta$ decreasing in $x$. By Tonelli and
$\int_0^\infty xe^{-\theta x}\,dx = \theta^{-2}$,
$\int_0^\infty x\,k_\infty(x)\,dx = \int_{1/c}^\infty2\theta^{-3}\,d\theta = c^2$. As
$k_\infty(x)/x \le x\,k_\infty(x)$ for $x \ge 1$, this bounds both $\int_0^1x\,k_\infty(x)\,dx$
and $\int_1^\infty k_\infty(x)\,dx/x$. By Tonelli and Lemma~\ref{lem:laplace-atoms-admissible}(1),
$\int_0^\infty(1 - \cos\omega x)\,k_\infty(x)\,dx/x
= \int_{1/c}^\infty\log(1 + \omega^2/\theta^2)\,d\theta/\theta = F_\infty(\omega)$. So $F_\infty$ is
admissible, with $a = 0$ and the profile $k_\infty$.

Fix $t > 0$. The exponents $\ell F_q(t\,\cdot)$ and $F_\infty(t\,\cdot)$ are admissible, a multiple
of an admissible exponent dilated by $t$, so by Lemma~\ref{lem:line-exponent-data}(1) the laws with
these exponents exist. Their transforms converge pointwise to $e^{-F_\infty(t\omega)}$ as
$q \downarrow 1$, by the locally uniform limit, and $e^{-F_\infty(t\,\cdot)}$ is the transform of
a law. Along every sequence $q_n \downarrow 1$, \Levy's continuity theorem gives the weak
convergence, and so it holds as $q \downarrow 1$. For the kernels, fix $\omega \ne 0$ and put
$y := t\omega$. The bound gives $F_q(y) \ge F_\infty(y)/\ell - \log(1 + c^2y^2)$, and the bound at
$q = e$ gives $F_\infty(y) \ge F_e(y) \ge \log(1 + c^2y^2e^{-2}) > 0$. So $F_q(y) \to \infty$ and
$\widehat\mu_t(\omega) = e^{-F_q(t\omega)} \to 0$ as $q \downarrow 1$. If the kernels converged
weakly to a law $\nu$, testing against $x \mapsto \cos(\omega x)$ would give
$\int\cos(\omega x)\,\nu(dx) = 0$ for every $\omega \ne 0$, while this function of $\omega$ is
continuous with the value $1$ at $\omega = 0$. So there is no weak limit
(Theorem~\ref{thm:matern-pyramid}(5) in every dimension).
\end{proof}

Among the families with its ladder, the \Matern{} pyramid is the only one covariant at every
scale, by (2). Each octave costs one forward and one backward first-order pass, by (3). The
similarity family interpolates the ladder between the knots.

Its exponent is the symmetric Thorin exponent whose Thorin measure is the geometric atom train
$U = \sum_{i \ge 1}2\,\delta_{q^i/c}$, in the normalization of
\prog{V-cone v0.1}{Prop.~5.7, (5.1), p.~48}. Its generator is the Thorin clause of the generator
proposition, a sum of Laplace-kernel differences \prog{V-cone v0.1}{Prop.~6.4(4), p.~58}. That
clause is stated for $a = 0$ and for $g \in C^2$ with $g$ and $g''$ bounded, and
Theorem~\ref{thm:matern-pyramid-line}(3) computes the symbol directly. So the level differences
$u(t) - u(qt) = (I - L_{ct})u(t)$ of its Laplacian pyramid are a bank of Laplace-kernel
differences, the atomic generators of the Thorin subclass. The Laplacian pyramid of the Gaussian
family is, in the same way, a bank of differences of Gaussians.

The subordination bridge turns an exponent $\Phi(s)$ of a time-causal family into the spatial
exponent $\Phi(\omega^2/2)$ (Theorem~\ref{thm:matern-pyramid}). It carries Lindeberg's time-causal
limit kernel of ratio $q^2$ and constant $2c^2$ to the \Matern{} pyramid. Lindeberg defines that
kernel as an infinite product of first-order factors
[\sscite{lindeberg2016time}, Sect.~5, (38), p.~62], the ratio and the constant being this reading of
it. He realizes it by recursive first-order filters [\sscite{lindeberg2016time}, Sect.~6.2, p.~64].
% Anchored 2026-11-12 by the librarian (records/pyramids/CITATIONS-acquired.md, item 1).
The limit $F_\infty$ of $\ell F_q$ in (\ref{eq:line-e8}) is the symmetric Thorin exponent with the log-uniform
density $d\theta/\theta$ above the cutoff $1/c$, and it has no Gaussian part.
% The draft's comparison "unlike the limit kernel's published form" (an uncited claim, review flag
% "Uncited 'published form' of the limit kernel") is cut, not cited. Its parenthesis on Phillips'
% formula and [V-cone v0.1, Remark 6.3], with the open check N3, is cut as well.

% draft: Theorem 2.17(3) (section 2.5), lines 229 and 235-236 of draft/spatial-hemigroup-pyramids.md.
% CHANGED: split from Theorem 2.17 (different grade). The counting reading ("creates no extrema
% from the signal at any level and none rung to rung") moves to prop:ladder-non-creation-line,
% whose counting clause carries it; this node states the Polya frequency facts.
% CHANGED (fifth blind pass, 2026-11-15, flag 2): lem:laplace-atoms-admissible joins the statement's
% uses, which the proof reads for the admissibility of log(1 + omega^2).
\begin{theorem}[the \Matern{} pyramid: \Polya{} frequency on the line]
\label{thm:matern-pyramid-line-noncreation}
\lean{Pyramids.matern_pyramid_line_noncreation}
\leanok
\uses{thm:matern-pyramid-line, thm:cited-A22, thm:cited-A24, lem:laplace-atoms-admissible}
In the \Matern{} pyramid of Theorem~\ref{thm:matern-pyramid-line}, every kernel $\mu_t$, $t > 0$, is
a \Polya{} frequency density, and every octave step $\mu_{t,qt} = L_{ct}$, $t > 0$, is a \Polya{}
frequency density.
% CHANGED (fifth blind pass, 2026-11-15, flag 1, the author's decision): "t > 0" added to the octave
% step; at t = 0 the step is delta_0, no Polya frequency density.
% CHANGED (the theory paper, 2026-11-16): the next sentence had been joined onto the comment line
% above, so the printed statement lacked the clause that matern_pyramid_line_noncreation states
% (its third conjunct) and the proof proves. Restored as it stood.
An increment $\mu_{s,t}$ with $0 < s < t$ and $t/s \notin q^{\ZZ}$ is not a \Polya{}
frequency density.
% CHANGED (wave 3b, 2026-11-15): A24 is read on the imaginary axis, the form admitted (ADR-0005).
% CHANGED (fidelity review, 2026-11-15, R26/R50/R69): the status line lists the declarations' hypotheses beside the node's (LINKAGE rule 6).
\statusA\quad\ledger{A22}\quad\ledger{A24}\quad\emph{(\textbf{Assignment.} \ledger{A22}
(\prog{V-selection v0.1}{Prop.~4.5(1), p.~35}, an equivalence of which the ``if'' is used) carries
the first sentence: the similarity family of $a\omega^2 + \sum_i\log(1 + \omega^2/\vartheta_i^2)$
with $\sum_i\vartheta_i^{-2} < \infty$ has \Polya{} frequency kernels, applied to $F_q$ and to
$\log(1 + \omega^2)$, whose admissibility, the entry's hypothesis, is proved here.
\ledger{A24} is Schoenberg's representation, printed in
\prog{V-selection v0.1}{\S~4.3, pp.~34--35} from Karlin. It carries the property of a \Polya{}
frequency density used in the second sentence: on the imaginary axis the reciprocal of its
two-sided Laplace transform is the restriction of an entire function. Karlin's class has
$\gamma \ge 0$, real $a_j$, $\sum_ja_j^2 < \infty$ and the normalization
$\gamma + \sum_ja_j^2 > 0$ (Remark~3.1, p.~346). The complex analysis of the increment is proved
here; the strip of Karlin's Prop.~1.4 is not read. The statement is on the line, and nothing is
claimed in the plane. \textbf{Hypotheses} (the node's, each with what its declarations require):
$q > 1$, $c > 0$: same; the \Matern{} pyramid: same, any laws with its transforms
(\texttt{IsMaternPyramid}); $t > 0$ for the kernels and for the octave steps: same; $0 < s < t$ and
$t/s \notin q^{\ZZ}$: same.)}
\end{theorem}

\begin{proof}
\leanok
% CHANGED (wave 3b, 2026-11-15)
% The proof follows the checked route. The old proof extended 1/L from a strip, as a meromorphic
% function, and cancelled zeros; the checked proof applies the identity theorem on C to entire
% functions that agree on the imaginary axis, and reads the vanishing of a convergent product.
\uses{thm:matern-pyramid-line, thm:cited-A22, thm:cited-A24, lem:laplace-atoms-admissible}
$F_q = \sum_{i \ge 1}\log(1 + \omega^2/\vartheta_i^2)$ with $\vartheta_i := q^i/c$ and
$\sum_i\vartheta_i^{-2} = c^2/(q^2 - 1) < \infty$, it is admissible by
Theorem~\ref{thm:matern-pyramid-line}(1), and $F_q \not\equiv 0$. The kernel $\mu_t$ has the
transform $e^{-F_q(t\,\cdot)}$, since $F_q(0) = 0$ and $F_q$ is even. So
Theorem~\ref{thm:cited-A22} with $a = 0$, applied
to the family $t \mapsto \mu_{0,t}$, makes every kernel $\mu_t$ a \Polya{} frequency density. The
step is $\mu_{t,qt} = L_{ct}$ (Theorem~\ref{thm:matern-pyramid-line}(1)), and
$\widehat{L_r}(\omega) = (1 + r^2\omega^2)^{-1} = e^{-F_1(r\omega)}$ with
$F_1(\omega) := \log(1 + \omega^2)$. By Lemma~\ref{lem:laplace-atoms-admissible}(2) with $a = 0$
and the single atom $\vartheta = 1$, $F_1$ is admissible, of the form of
Theorem~\ref{thm:cited-A22}, and not identically zero. That theorem, applied to the family
$r \mapsto L_r$, makes $L_{ct}$ a \Polya{} frequency density.

For the increments, let $\vartheta_i := q^i/c$, $i \ge 1$, and for $r > 0$
\[
  N_r(p) := \prod_{i \ge 1}\Bigl(1 - \frac{r^2}{\vartheta_i^2}\,p^2\Bigr), \qquad p \in \CC .
\]
On a compact set $|p| \le R$ the $i$-th factor differs from $1$ by at most $r^2R^2\vartheta_i^{-2}$,
which is summable, so the product converges uniformly on compact sets and $N_r$ is entire. At
$p = i\omega$ the factors are $1 + r^2\omega^2/\vartheta_i^2 > 0$, and the series
$\sum_i\log(1 + r^2\omega^2/\vartheta_i^2) = F_q(r\omega)$ gives $N_r(i\omega) = e^{F_q(r\omega)}$.
So, for $0 < s < t$ and every real $\omega$,
\[
  \widehat\mu_{s,t}(-\omega)\,N_t(i\omega) = e^{-(F_q(t\omega) - F_q(s\omega))}\,e^{F_q(t\omega)}
  = N_s(i\omega).
\]
Suppose $\mu_{s,t}$ had a \Polya{} frequency density $\varphi$. Its two-sided Laplace transform
on the imaginary axis is $\mathcal{L}(i\omega) = \int e^{-i\omega x}\varphi(x)\,dx =
\widehat\mu_{s,t}(-\omega)$. Theorem~\ref{thm:cited-A24} gives an entire function $\psi$ with
$\mathcal{L}(i\omega)\,\psi(i\omega) = 1$ for every real $\omega$. Multiplying the display by
$\psi(i\omega)$ gives $N_s(i\omega)\,\psi(i\omega) = N_t(i\omega)$ for every real $\omega$. The
entire functions $N_s\psi$ and $N_t$ thus agree on the imaginary axis, which accumulates at $0$,
so by the identity theorem they agree on $\CC$. At $p_0 := q/(cs)$ the first factor of $N_s$ is
$1 - (s/\vartheta_1)^2p_0^2 = 0$, so $N_s(p_0) = 0$ and hence $N_t(p_0) = 0$. The factors of $N_t$
at $p_0$ are $1 + f_i$ with $f_i = -(t/\vartheta_i)^2p_0^2$ and $\sum_i|f_i| < \infty$, and such a
product vanishes only if one of its factors does. So $(t/\vartheta_i)^2p_0^2 = 1$ for some
$i \ge 1$, that is $(t/s)^2 = q^{2(i-1)}$, and $t/s = q^{i-1} \in q^{\ZZ}$. When
$t/s \notin q^{\ZZ}$ this is a contradiction.
\end{proof}

The non-creation axiom on the line has four readings, which select three nested classes
\prog{V-selection v0.1}{Remark~4.15, pp.~50--51}. Pointwise non-enhancement at global extrema holds for every
admissible family \prog{V-selection v0.1}{Thm.~4.2(2), p.~30}. Endpoint counting, no more zero
crossings and no more changes of monotonicity at any scale than in the signal, selects the \Polya{}
frequency subclass \prog{V-selection v0.1}{Prop.~4.5, p.~35}. Pointwise non-enhancement at local
extrema \prog{V-selection v0.1}{Thm.~4.2(3), Cor.~4.4, pp.~31, 33} and scale-monotone counting, the variation
diminution of every interior increment \prog{V-selection v0.1}{Thm.~4.10, p.~43}, both select the
Gaussian ray. The discrete stratum has a further reading, weaker than scale-monotone counting:
\emph{ladder non-creation}, the variation diminution of the increments along one ratio $q$ only.
Required between every pair of scales, the same requirement is scale-monotone counting and selects the
Gaussian. Required along one ratio, it holds for a class of families that includes the fully covariant
\Matern{} pyramid.


% draft: Remark 2.17' (section 2.5), lines 241-249 of draft/spatial-hemigroup-pyramids.md (the
% statement and proof in the remark's second and third paragraphs).
% CHANGED: the increments are defined by their transforms, so the statement does not presuppose
% the family of F; the gloss "S^-(Phi * g) <= S^-(g) for every integrable or bounded g" is not
% stated (it is the variation-diminishing property that the counting clause uses through A23); the
% admissibility of G is proved here from lem:laplace-atoms-admissible, which replaces the draft's
% route through [V-selection v0.1, Lemma 4.7], [V-cone v0.1, Prop. 5.7], Lemma 2.7 and Prop. 3.3.
% CHANGED: "the twin of Paper I, Remark 8.16" is cut (Paper I is not released).
% CHANGED (fifth blind pass, 2026-11-15, flags 3-5): thm:matern-pyramid-line, named by the
% statement's last sentence, joins its uses; the Assignment says that A23 is proved from A22 at
% thm:cited-A23, in print, where it said "in the Lean"; the proof says that every multiplicity is
% finite, so that m(theta) is a difference of integers.
\begin{proposition}[ladder non-creation, on the line]
\label{prop:ladder-non-creation-line}
\lean{Pyramids.ladder_non_creation_line, Pyramids.ladder_non_creation_line_matern}
\leanok
% CHANGED (wave 3b, 2026-11-15): the uses follow the checked proof.
\uses{thm:cited-A22, thm:cited-A23, lem:laplace-atoms-admissible, lem:line-exponent-data,
thm:matern-pyramid-line}
Let $F(\omega) = a\omega^2 + \sum_i\log(1 + \omega^2/\vartheta_i^2)$, $F \not\equiv 0$, with
$a \ge 0$, $\vartheta_i > 0$ for $i$ in a countable index set, and $\sum_i\vartheta_i^{-2} < \infty$.
Write $\mathrm{mult}(\vartheta)$ for the number of $i$ with $\vartheta_i = \vartheta$. Suppose
$\mathrm{mult}(q\vartheta) \ge \mathrm{mult}(\vartheta)$ for every $\vartheta > 0$. For $s > 0$ let
$\mu_s$ be the law with transform $e^{-F(s\omega)}$, and $\mu_{s,qs}$ the law with transform
$e^{-(F(qs\omega) - F(s\omega))}$, so that $\mu_{qs} = \mu_{s,qs} * \mu_s$.
% CHANGED (supplement review, 2026-11-16): the closing sentence "Theorem matern-pyramid-line is
% the case ..." is clause (3), a claim with its own declaration (ladder_non_creation_line_matern)
% that chapter 4 reads, so it stays in the statement rather than moving to a remark; "Then" is
% dropped so that (3) does not read as a consequence of the hypotheses. No claim changed.
\begin{enumerate}
\item Every $\mu_{s,qs}$, $s > 0$, is a \Polya{} frequency density.
\item Let $f$ be the primitive of an integrable $f'$, and $u(s) := \mu_s * f$ its record at scale
$s$. For every $s > 0$, $u(qs)$ crosses each level $\alpha \in \RR$ no more often than $u(s)$, and
changes monotonicity no more often than $u(s)$.
\item The exponent $F_q$ of Theorem~\ref{thm:matern-pyramid-line} is of this form, with $a = 0$ and
$\vartheta_i = q^i/c$, $i \ge 1$, and for it $\mathrm{mult}(q\vartheta) - \mathrm{mult}(\vartheta)$
is $1$ at $\vartheta = 1/c$ and $0$ elsewhere.
\end{enumerate}
% CHANGED (fidelity review, 2026-11-15, R26/R50/R69): the status line lists the declarations' hypotheses beside the node's (LINKAGE rule 6).
\statusA\quad\ledger{A22}\quad\ledger{A23}\quad\emph{(\textbf{Assignment.} \ledger{A22}
(\prog{V-selection v0.1}{Prop.~4.5(1), p.~35}) carries (1), applied to the exponent $G$ of the
proof, whose form is proved here. \ledger{A23} (\prog{V-selection v0.1}{Prop.~4.5(2), p.~35}, derived
there from clause (1) on p.~39, and proved from A22 at Theorem~\ref{thm:cited-A23}), the counting clause for the
records of the similarity family of $G$, carries (2), applied to the signal $u(s)$. The
rearrangement of $F(q\,\cdot) - F$, the bound on the multiplicities, $G \not\equiv 0$, the
admissibility of $F$ and of $G$, the entry's hypothesis, and the existence of the laws $\mu_s$ and
$\mu_{s,qs}$ are proved here. On the line only. \textbf{Hypotheses} (the node's, each with what its
declarations require): $q > 1$: same; $a \ge 0$, $\vartheta_i > 0$, the index set countable,
$\sum_i\vartheta_i^{-2} < \infty$: same; $F$ the series: same, given as the sum; $F \not\equiv 0$:
same; $\mathrm{mult}(q\vartheta) \ge \mathrm{mult}(\vartheta)$ for $\vartheta > 0$: same, counted in
$\NN \cup \{\infty\}$; the laws $\mu_s$, $\mu_{s,qs}$, $s > 0$: same, any probability laws with these
transforms, whose existence the declaration proves; (2) $f$ the primitive of an integrable $f'$:
same; every level $\alpha$: same; the \Matern{} case: $q > 1$, $c > 0$, same.)}
\end{proposition}

\begin{proof}
\leanok
% CHANGED (wave 3b, 2026-11-15)
% The proof follows the checked route: the rearrangement is made with multiplicities in [0, infty]
% and needs no finiteness of the fibres; G not identically zero is read off the inverse squares;
% the record u(qs) is mu_{s,qs} * u(s) at every point, and its derivative is identified a.e.
\uses{thm:cited-A22, thm:cited-A23, thm:matern-pyramid-line, lem:laplace-atoms-admissible,
lem:line-exponent-data}
\emph{The rearrangement.} For a function $g \ge 0$ on $(0,\infty)$ and a family $(x_k)$, the sum
$\sum_kg(x_k)$ in $[0,\infty]$ is $\sum_\vartheta\mathrm{mult}_x(\vartheta)\,g(\vartheta)$, summing
over the fibres of the family; multiplicities are read in $\{0, 1, \dots, \infty\}$. The family
$(\vartheta_i/q)$ has $\mathrm{mult}(q\vartheta)$ terms equal to $\vartheta$. Every multiplicity of
the family $(\vartheta_i)$ is finite, as $\mathrm{mult}(\vartheta)\,\vartheta^{-2} \le
\sum_i\vartheta_i^{-2} < \infty$. Put
$m(\vartheta) := \mathrm{mult}(q\vartheta) - \mathrm{mult}(\vartheta)$, a nonnegative integer by
the hypothesis, and let
$(\vartheta'_j)_{j \in J}$ list each $\vartheta$ exactly $m(\vartheta)$ times; $J$ is countable,
since $m(\vartheta) > 0$ only for $\vartheta$ among the countably many $\vartheta_i/q$. Then
$\mathrm{mult}(q\vartheta) = \mathrm{mult}(\vartheta) + m(\vartheta)$ for every $\vartheta$ (both
sides vanish for $\vartheta \le 0$), so
\[
  \sum_ig(\vartheta_i/q) = \sum_ig(\vartheta_i) + \sum_jg(\vartheta'_j)
\]
in $[0,\infty]$, and in $\RR$, with both sums on the right finite, whenever the left side is finite.
At $g(\vartheta) = \vartheta^{-2}$ the left side is $q^2\sum_i\vartheta_i^{-2} < \infty$, so
$\sum_j\vartheta'^{-2}_j < \infty$. At $g(\vartheta) = \log(1 + \omega^2/\vartheta^2) \le
\omega^2\vartheta^{-2}$, with $\log(1 + q^2\omega^2/\vartheta_i^2) = g(\vartheta_i/q)$,
\[
  G(\omega) := F(q\omega) - F(\omega)
  = a(q^2 - 1)\,\omega^2 + \sum_j\log\Bigl(1 + \frac{\omega^2}{\vartheta'^2_j}\Bigr).
\]
By Lemma~\ref{lem:laplace-atoms-admissible}(2), $F$ and $G$ are admissible, each with a folded
pair, and $G \not\equiv 0$ unless $a(q^2 - 1) = 0$ and $J$ is empty. If $a > 0$, then
$a(q^2 - 1) > 0$. If $a = 0$, the family $(\vartheta_i)$ is not empty, since $F \not\equiv 0$, so
$T := \sum_i\vartheta_i^{-2} \in (0, \infty)$; were $J$ empty, the identity at
$g(\vartheta) = \vartheta^{-2}$ would read $q^2T = T$. So $G \not\equiv 0$. By
Lemma~\ref{lem:line-exponent-data}(1),(2) the laws $\mu_s$ and $\mu_{s,qs}$ exist for every
$s > 0$, and $\mu_{qs} = \mu_{s,qs} * \mu_s$ because their transforms multiply.

(1) $\mu_{s,qs}$ has the transform $e^{-G(s\omega)}$, so $s \mapsto \mu_{s,qs}$ is a family of
kernels of the similarity family of $G$. $G \not\equiv 0$ is admissible and of the form of
Theorem~\ref{thm:cited-A22}. That proposition makes
$\mu_{s,qs}$ a \Polya{} frequency density.

(2) The primitive $f = C + \int_0^xf'$ is continuous and bounded by $|C| + \|f'\|_1$, so
$u(qs) = (\mu_s * \mu_{s,qs}) * f = \mu_{s,qs} * u(s)$ at every point, by Fubini's theorem for a
bounded measurable function. By Fubini's theorem on $[0,x] \times \mu_s$,
$u(s)(x) = u(s)(0) + \int_0^x(\mu_s * f')$, so $u(s)$ is the primitive of the integrable
$\mu_s * f'$. The counting clause Theorem~\ref{thm:cited-A23}, applied to the family $t \mapsto \mu_{t,qt}$ of $G$ at scale $s$
and the signal $u(s)$, gives $S^-(u(qs) - \alpha) \le S^-(u(s) - \alpha)$ for every $\alpha$, and
$S^-(\partial_xu(qs)) \le S^-(\mu_s * f')$. By the Lebesgue differentiation theorem
$\partial_xu(s) = \mu_s * f'$ almost everywhere, and $S^-$, a count in the essential sense, does not
see the null set; so the right side is $S^-(\partial_xu(s))$.

(3) By (\ref{eq:line-e7}), $F_q(\omega) = \sum_{i \ge 1}\log(1 + \omega^2/\vartheta_i^2)$ with
$\vartheta_i = q^i/c$. The map $i \mapsto q^i/c$ is injective, so $\mathrm{mult}$ is $1$ on $\{q^i/c : i \ge 1\}$
and $0$ elsewhere; $q\vartheta$ lies in that set exactly when $\vartheta = 1/c$ or $\vartheta$ does,
and $1/c$ does not. So $m(\vartheta) = 1$ at $\vartheta = 1/c$ and $0$ elsewhere.
\end{proof}

The statement is on the line, as is every counting reading here \prog{V-selection v0.1}{\S~6, p.~69}.

% draft: Remark 2.18(1)-(2) (section 2.5), lines 251-254 of draft/spatial-hemigroup-pyramids.md.
% CHANGED: Remark 2.18 made a proposition, its assertive parts each proved. The draft's (1) says
% the equivalent weight "is Gaussian to relative order 1/n in the bulk, with excess kurtosis
% -6/5 . 1/(4n)-type corrections"; its check required to compute the fourth cumulant exactly. The exact
% computation (clause (2)) gives excess kurtosis -(3/10)(4^n+1)/(4^n-1), which tends to -3/10, not
% to 0, and the rescaled weights converge to the cubic B-spline (clause (3); Proposition 3.15(2)),
% not to a Gaussian. The central-limit reading is false and is not transcribed.
% CHANGED: the lattice non-creation of the binomial mask is proved here (clause (4)) in place of
% the citation of Lindeberg's Theorem 1, for which no page is held.
\begin{proposition}[what the binomial pyramid converges to]
\label{prop:classical-pyramids-approximate}
\lean{Pyramids.classical_pyramids_approximate_1, Pyramids.classical_pyramids_approximate_2,
Pyramids.classical_pyramids_approximate_3, Pyramids.classical_pyramids_approximate_4,
Pyramids.classical_pyramids_approximate_5}
\leanok
\uses{prop:classical-kernels-no-exact-cascade}
Let $h > 0$, let
$w := \tfrac1{16}(\delta_{-2h} + 4\delta_{-h} + 6\delta_0 + 4\delta_h + \delta_{2h})$ be the binomial
mask, and let $W_n := w * T_2w * \cdots * T_{2^{n-1}}w$, $n \ge 1$, be its equivalent weight at level
$n$.
\begin{enumerate}
\item $W_n$ is carried by $[-2(2^n - 1)h,\ 2(2^n - 1)h]$, and
$\widehat W_n(\omega) = \prod_{j=0}^{n-1}\cos^4(2^{j-1}h\omega)$ vanishes at $\omega = \pi/h$.
\item $W_n$ has variance $(4^n - 1)h^2/3$ and excess kurtosis
$-\tfrac{3}{10}\,(4^n + 1)/(4^n - 1)$, which is less than $-\tfrac3{10}$ for every $n$.
\item As $n \to \infty$, $T_{2^{-n}}W_n$ converges weakly to the centred cubic B-spline of knot
spacing $h$, the law of the sum of four independent uniform variables on $[-h/2, h/2]$, with
transform $(\sin(h\omega/2)/(h\omega/2))^4$ and excess kurtosis $-\tfrac3{10}$. The rescaled
equivalent weights do not converge to a Gaussian law.
% CHANGED: the sequences of (4) are indexed by k, since n is the level (blind review 2, flag 19).
\item (Lattice non-creation.) For a real sequence $x$ write $S^-(x)$ for its number of sign changes,
zeros deleted, and $\Delta x(k) := x(k+1) - x(k)$. For every real sequence $x$ on $h\ZZ$, the
smoothed sequence $y := w * x$ and its decimation $z(k) := y(2k)$ on $2h\ZZ$ satisfy
\[
  S^-(z) \le S^-(y) \le S^-(x), \qquad S^-(\Delta z) \le S^-(\Delta y) \le S^-(\Delta x).
\]
\item The B-spline kernel of order $m \ge 1$ and spacing $h$, the $m$-fold convolution of the uniform
law on $[-h/2, h/2]$, has bounded support and the transform $(\sin(h\omega/2)/(h\omega/2))^m$, which
vanishes on $(2\pi/h)\ZZ \setminus \{0\}$.
\end{enumerate}
% CHANGED (fidelity review, 2026-11-15, R26/R50/R69): the status line lists the declarations' hypotheses beside the node's (LINKAGE rule 6).
\statusT\quad\emph{(\textbf{Hypotheses} (the node's, each with what its declarations require): $h > 0$:
same, absent from \texttt{classical\_pyramids\_approximate\_4}, whose sequences are indexed by
$\ZZ$; $n \ge 1$: same in (1) and (2), and in (3) the limit $n \to \infty$; (4) every real
sequence $x$: same; (5) $m \ge 1$: same.)}
\end{proposition}

% CHANGED (wave 3a, 2026-11-14): the proof follows the machine-checked route. In (2) the additivity
% of the variance and of the fourth cumulant is derived for laws, from the binomial expansion of
% the moments of a convolution. In (3) the limit is not Gaussian because weak limits are unique and
% the transform of the B-spline vanishes at 2 pi/h, which no Gaussian transform does; the kurtosis
% comparison it replaces left out the point masses. In (4) the sign changes are counted along chains.
% CHANGED (supplement review, 2026-11-16): the proof-local laws of (2) renamed P, Q (bare mu is the
% step measure).
\begin{proof}
\leanok
\uses{prop:equivalent-weights}
(1) $T_{2^j}w$ is carried by $[-2^{j+1}h, 2^{j+1}h]$, and the sum of these bounds over $j < n$ is
$2(2^n - 1)h$. The transform of $w$ is $\tfrac1{16}(6 + 8\cos h\omega + 2\cos 2h\omega)
= \bigl(\tfrac12(1 + \cos h\omega)\bigr)^2 = \cos^4(h\omega/2)$, and
$\widehat{T_{2^j}w}(\omega) = \widehat w(2^j\omega)$. The factor $j = 0$ vanishes at
$h\omega = \pi$.

(2) Let $P, Q$ be laws of bounded support. By Fubini's theorem
$\int x^k\,(P * Q)(dx) = \iint(x + y)^k\,P(dx)\,Q(dy)$, and by the binomial theorem this is
$\sum_{m \le k}\frac{k!}{m!\,(k-m)!}\int x^m\,P(dx)\int y^{k-m}\,Q(dy)$. Let both have mean $0$,
and write $\kappa_2 := \int x^2$ and $\kappa_4 := \int x^4 - 3\kappa_2^2$. Then $P * Q$ has mean
$0$ and $\kappa_2(P * Q) = \kappa_2(P) + \kappa_2(Q)$. Its fourth moment is
$\int x^4\,dP + 6\kappa_2(P)\kappa_2(Q) + \int y^4\,dQ$, since the terms with a first moment
vanish, so $\kappa_4(P * Q) = \kappa_4(P) + \kappa_4(Q)$. The dilate $T_c$ multiplies
$\kappa_2$ by $c^2$ and $\kappa_4$ by $c^4$. The mask $w$ has mean $0$,
$\kappa_2(w) = \tfrac{h^2}{16}(4 + 4 + 4 + 4) = h^2$ and
$\int x^4\,dw = \tfrac{h^4}{16}(16 + 4 + 4 + 16) = \tfrac52h^4$, so $\kappa_4(w) = -\tfrac12h^4$.
These are the cumulants of $hX$ with $X = B - 2$, $B$ binomial with $4$ trials and success
probability $\tfrac12$. So $W_n$ has mean $0$,
$\kappa_2(W_n) = h^2\sum_{j<n}4^j = h^2(4^n - 1)/3$ and
$\kappa_4(W_n) = -\tfrac{h^4}2\sum_{j<n}16^j = -h^4(16^n - 1)/30$. Since
$16^n - 1 = (4^n - 1)(4^n + 1)$, the excess kurtosis is
$\kappa_4/\kappa_2^2 = -\tfrac{9(16^n - 1)}{30(4^n - 1)^2} = -\tfrac3{10}\cdot\tfrac{4^n + 1}{4^n - 1}$.
It is less than $-\tfrac3{10}$, since $4^n + 1 > 4^n - 1 > 0$ for $n \ge 1$.

(3) The transform of the rescaled weight is
\[
  \widehat{T_{2^{-n}}W_n}(\omega) = \prod_{j=0}^{n-1}\cos^4(2^{j-n-1}h\omega)
  = \prod_{m=1}^{n}\cos^4\bigl(2^{-m}\,h\omega/2\bigr).
\]
It tends to $(\sin(h\omega/2)/(h\omega/2))^4$ by Viète's product
$\prod_{m \ge 1}\cos(2^{-m}x) = \sin x/x$, which is Proposition~\ref{prop:equivalent-weights}(2). The
limit is the transform of the cubic B-spline, so Lévy's continuity theorem gives weak convergence.
The uniform law on $[-h/2,h/2]$ has the transform $\sin(h\omega/2)/(h\omega/2)$, mean $0$,
$\kappa_2 = h^2/12$ and $\kappa_4 = h^4/80 - 3h^4/144 = -h^4/120$. By the additivity in (2) the
fourfold convolution has $\kappa_2 = h^2/3$ and $\kappa_4 = -h^4/30$, so its excess kurtosis is
$-\tfrac{h^4}{30}\cdot\tfrac{9}{h^4} = -\tfrac3{10}$. Weak limits are unique, so a Gaussian limit,
of mean $m$ and variance $v \ge 0$, would be the cubic B-spline. But the Gaussian transform
$e^{im\omega - v\omega^2/2}$ vanishes nowhere, while that of the B-spline vanishes at
$\omega = 2\pi/h$ by (5). So the rescaled weights have no Gaussian limit, a point mass ($v = 0$)
included.

(4) A sequence $x$ has $S^-(x) \ge p$ exactly when there are $k_0 < \dots < k_p$ with
$x(k_i)x(k_{i+1}) < 0$ for $i < p$. A finite section with $p$ sign changes, zeros deleted, has such
an alternating subsequence of $p + 1$ terms: keep the first term, and from each later term keep
it when its sign differs from the last one kept, else let it replace that one. So
$S^-(x') \le S^-(x)$ as soon as each such chain of $x'$ gives one of $x$ of the same length.

First, for $y'(k) := x(k) + x(k-1)$, $S^-(y') \le S^-(x)$. Let
$y'(k_0), \dots, y'(k_p)$, $k_0 < \dots < k_p$, be nonzero with alternating signs. Since
$y'(k_i) \ne 0$, one of $x(k_i), x(k_i - 1)$ has the sign of $y'(k_i)$. Call its index $j_i$. Then
$j_i \le k_i \le k_{i+1} - 1 \le j_{i+1}$, and $j_i \ne j_{i+1}$ because the signs differ. So
$j_0 < \dots < j_p$, and $x$ has at least $p$ sign changes. The convolution with $w$ is four such
steps, a shift and a positive factor, so $S^-(y) \le S^-(x)$. Convolution commutes with
differences, $\Delta y = w * \Delta x$, so $S^-(\Delta y) \le S^-(\Delta x)$. A subsequence has no
more sign changes than its sequence, so $S^-(z) \le S^-(y)$. Finally
$\Delta z(k) = y(2k+2) - y(2k) = \Delta y(2k+1) + \Delta y(2k)$ is a subsequence of the sums of
consecutive terms of $\Delta y$, so $S^-(\Delta z) \le S^-(\Delta y)$ by the first step.

(5) The transform is the $m$-th power of that of the uniform law, which vanishes exactly at
$h\omega \in 2\pi\ZZ \setminus \{0\}$, and the support is $[-mh/2, mh/2]$.
\end{proof}

Both pyramids create no extrema from one level to the next. For the binomial pyramid this is (4), a
statement about one mask on one lattice. For the \Matern{} pyramid it is
Proposition~\ref{prop:ladder-non-creation-line}, a statement about its octave steps. Its increments
between scales whose ratio is not a power of $q$ are not \Polya{} frequency densities
(Theorem~\ref{thm:matern-pyramid-line-noncreation}). The binomial pyramid has no increments between
levels at all, since no family continuous in scale has its steps
(Proposition~\ref{prop:classical-kernels-no-exact-cascade}).

The binomial pyramid's equivalent weights have finite support and transforms with zeros at every
level. Rescaled, they converge to the cubic B-spline, whose resemblance to the Gaussian is what Burt
and Adelson observed. The \Matern{} pyramid is a family of the axioms, covariant at every scale. By
(5) the B-spline pyramids also have kernels of bounded support whose transforms vanish, so the same
proposition excludes them.

Lindeberg's discrete scale space is exact but not a pyramid
(Proposition~\ref{prop:lindeberg-stratum-one}). Its grid is a preferred length. The grid refinement
that makes it covariant, Proposition~\ref{prop:discrete-gaussian-pyramid}, lies strictly in the
discrete stratum, stratum (2) of Theorem~\ref{thm:lattice-line}, and its kernels are no longer
lattice laws. Lindeberg characterizes the discrete Gaussian, up to a constant in scale, as the unique
symmetric normalized discrete kernel with non-creation and a semigroup law
[\sscite{lindeberg1990scalespace}]. That characterization rests on
the semigroup law. Dropping it leaves a class of kernels in place of one family
\prog{V-selection v0.1}{Remark~4.15, p.~51}. No non-creation property is claimed here for the
geometric-jump pyramid. The next proposition says why the line's two tools do not apply to it.

% draft: Remark 2.18(3) (section 2.5), line 255 of draft/spatial-hemigroup-pyramids.md, its last two
% sentences.
% CHANGED: Remark 2.18(3) made a proposition (its sentences assert and argue); split from
% prop:classical-pyramids-approximate by grade.
% CHANGED (proof): the draft cites "a Polya frequency density has an absolutely continuous Levy
% measure" to [V-selection v0.1, section 4.3]; the librarian found no such clause there
% (records/pyramids/VERBATIM-selection-iso-affine.md). It is deduced here from Schoenberg's
% representation (A24) and the uniqueness of the triplet (A39). Ledger A25 is not used.
% CHANGED (fifth blind pass, 2026-11-15, flags 6, 7): the statement names its kernel,
% mu_t := mu_{0,t}, which it used undefined; the proof writes the atom train of
% prop:discrete-gaussian-pyramid with its weights w_i(eta), where it named undefined weights W_t(j).
\begin{proposition}[the geometric-jump pyramid is neither lattice nor \Polya{} frequency]
\label{prop:classical-pyramids-approximate-discrete-gaussian}
\lean{Pyramids.discrete_gaussian_neither}
\leanok
% CHANGED (wave 3b, 2026-11-15): the uses follow the checked proof.
\uses{prop:discrete-gaussian-pyramid, thm:cited-A24}
Let $(\mu_{s,t})$ be the geometric-jump pyramid of Proposition~\ref{prop:discrete-gaussian-pyramid},
with weights not all zero. For $t > 0$ the kernel $\mu_t := \mu_{0,t}$ is carried by no lattice, and
it is not a \Polya{} frequency density.
% CHANGED (2026-11-15, a strengthening)
% Proved from fewer cited facts: the status line was "[A] on A24 and A39", with A39(ii)-(iii)
% read for laws that need not be symmetric. The checked proof reads A24 on the imaginary axis and
% reaches A39 only through lem:line-exponent-data and prop:discrete-gaussian-pyramid.
% CHANGED (fidelity review, 2026-11-15, R26/R50/R69): the status line lists the declarations' hypotheses beside the node's (LINKAGE rule 6).
\statusA\quad\ledger{A24}\quad\emph{(\textbf{Assignment.} \ledger{A24},
Schoenberg's representation (\prog{V-selection v0.1}{\S~4.3, pp.~34--35}, printing Karlin), carries
the form of the reciprocal Laplace transform of a \Polya{} frequency density on the imaginary axis,
with $\gamma \ge 0$, real $a_j$, $\sum_ja_j^2 < \infty$. No other cited fact is read: the
uniqueness of the folded pair is Lemma~\ref{lem:line-exponent-data}(1), the atomic form of the
kernel's \Levy{} measure is Proposition~\ref{prop:discrete-gaussian-pyramid}(1), and the lattice
clause is Proposition~\ref{prop:discrete-gaussian-pyramid}(3), all proved, which carry A39 and A41
through the nodes they use. That the exponent of a \Polya{} frequency density has an absolutely
continuous folded \Levy{} measure is proved here. \textbf{Hypotheses} (the node's, each with what
its declarations require): the setting of Proposition~\ref{prop:discrete-gaussian-pyramid}, $q > 1$,
$h > 0$, $w_i \ge 0$ nonincreasing, $\sum_{i < 0}w_iq^{2i} < \infty$, $\sum_{i \ge 0}w_i < \infty$:
same; weights not all zero: same; the geometric-jump pyramid: same, any laws with its transforms
(\texttt{IsDGKernels}); $t > 0$: same.)}
\end{proposition}

\begin{proof}
\leanok
% CHANGED (supplement review, 2026-11-16): the octave fraction written zeta, as in
% prop:discrete-gaussian-pyramid (eta is the Haar measure, theta the phase of a knot).
% CHANGED (wave 3b, 2026-11-15)
% The proof follows the checked route, through the modulus of Karlin's function on the imaginary
% axis. The old proof wrote the transform of a Polya frequency density as a Levy-Khintchine
% exponent by a complex Frullani integral and read the triplet's uniqueness, A39(ii)-(iii).
\uses{prop:discrete-gaussian-pyramid, thm:cited-A24, lem:line-exponent-data,
lem:laplace-atoms-admissible}
By Proposition~\ref{prop:discrete-gaussian-pyramid}(3), no kernel $\mu_t$, $t > 0$, is carried by a
lattice $\xi\ZZ$, $\xi > 0$.

\emph{The kernel's \Levy{} measure.} Write $\Psi_t$ for the exponent of $\mu_t$,
$\widehat\mu_t = e^{-\Psi_t}$. By Proposition~\ref{prop:discrete-gaussian-pyramid}(1), $\Psi_t$ has
the folded data $(0, \varpi_t)$, where, for $t = q^{n+\zeta}$ with $n \in \ZZ$ and $\zeta \in [0, 1)$,
$\varpi_t = \sum_{i \in \ZZ}w_i(\zeta)\,\delta_{hq^{i+n}}$ is an atom train with the nonnegative
weights $w_i(\zeta) = (1 - \zeta)w_i + \zeta w_{i-1}$. They are not all zero: if $w_i > 0$, then
$w_i(\zeta) \ge (1 - \zeta)w_i > 0$. So $\varpi_t(\{hq^j\}) > 0$ for some $j$.

\emph{A \Polya{} frequency density.} Suppose $\mu_t$ had a \Polya{} frequency density $\varphi$,
with two-sided Laplace transform $\mathcal{L}(p) = \int e^{-px}\varphi(x)\,dx$. Schoenberg's
representation (Theorem~\ref{thm:cited-A24}; [\sscite{karlin1968total}, Ch.~7, Thm.~3.2(a),
p.~345], printed in \prog{V-selection v0.1}{\S~4.3, pp.~34--35}) gives $\gamma \ge 0$ and real
$\delta$, $a_j$ with $\sum_ja_j^2 < \infty$ such that
$\mathcal{L}(i\omega)\,\psi(i\omega) = 1$ for every real $\omega$, with
$\psi(p) = e^{-\gamma p^2 + \delta p}\prod_j(1 + a_jp)\,e^{-a_jp}$. On the imaginary axis
$\mathcal{L}(i\omega) = \widehat\mu_t(-\omega) = e^{-\Psi_t(\omega)}$, a positive number, while
$|(1 + ia\omega)e^{-ia\omega}| = (1 + a^2\omega^2)^{1/2}$ and $|e^{\gamma\omega^2 + i\delta\omega}| =
e^{\gamma\omega^2}$. Taking moduli,
\[
  \Psi_t(\omega) = \gamma\omega^2 + \tfrac12\sum_j\log\bigl(1 + a_j^2\omega^2\bigr)
  \qquad (\omega \in \RR).
\]
The terms with $a_j = 0$ vanish. With $\vartheta_j := 1/|a_j|$ for the $j$ with $a_j \ne 0$,
\[
  2\Psi_t(\omega) = 2\gamma\,\omega^2 + \sum_{a_j \ne 0}\log\Bigl(1 + \frac{\omega^2}{\vartheta_j^2}\Bigr),
  \qquad \sum_{a_j \ne 0}\vartheta_j^{-2} = \sum_ja_j^2 < \infty .
\]
By Lemma~\ref{lem:laplace-atoms-admissible}(2), $2\Psi_t$ is the exponent of the pair
$(2\gamma, k(x)\,dx/x)$ with the profile $k(x) = 2\sum_je^{-\vartheta_jx}$, a folded \Levy{} measure
that is absolutely continuous. But $2\Psi_t$ is also the exponent of the pair $(0, 2\varpi_t)$, the
data of $\Psi_t$ added to themselves. By the uniqueness of the folded pair
(Lemma~\ref{lem:line-exponent-data}(1)), $2\varpi_t = k(x)\,dx/x$. The left side charges the point
$hq^j$ and the right side charges no point, a contradiction. So $\mu_t$ is not a \Polya{}
frequency density.
\end{proof}

Multiresolution analysis [\sscite{mallat1989theory}, (2)--(8), p.~676] is the axiomatic neighbour:
% Anchored 2026-11-12 by the librarian (records/pyramids/CITATIONS-acquired.md, item 4).
nested dyadic subspaces with the scaling relation as an axiom. It is an exact theory of dyadic
structure by projection rather than smoothing, with no positivity, no non-creation and no
hemigroup, since projections do not compose across levels as convolutions do. Its scaling relation $\phi = \sum_kc_k\,\phi(2\,\cdot - k)$ is the
knot relation of Definition~\ref{def:self-decomposable} without the requirement that the residual
be infinitely divisible, and the one-step exactness that it takes as an axiom is what
Proposition~\ref{prop:classical-kernels-no-exact-cascade} denies to a positive smoothing mask of
finite support.
% The draft's check on Remark 2.18(4), "the only nonnegative orthonormal scaling function is the
% Haar function, and whether it is worth stating here", is cut: the sentence makes no claim about it.

The theory of this appendix is discrete in scale on continuous space, and the pyramid principle of
chapter~3 says that sampling a $q^{\ZZ}$-covariant family at scale $q^nt$ on the grid $q^nh\ZZ$
stores the same array at every level. What the axioms do not fix is $h/t$, the sampling density
relative to the scale, which is the anti-aliasing trade-off that Meer and coauthors optimize for one
decimation [\sscite{meer1987frequency}]. For the \Matern{} pyramid the transform
$\prod_i(1 + c^2q^{-2i}t^2\omega^2)^{-1}$ of the level kernel at the grid's Nyquist wavenumber
$\pi/(q^nh)$ measures the aliasing of one level. The sampled Laplace kernel is a two-sided geometric
sequence, realized exactly by one causal and one anticausal first-order recursion. The sampled
Laplace family is an exact cascade on the grid, and its subsample is exact in sense (ii) of
Definition~\ref{def:exact-senses}. The grid pyramid of sampled \Matern{} levels is exact in sense (i)
at its first level and at no later one, whose kernels are convolutions of several Laplace kernels
and not the convolutions of their samples (Proposition~\ref{prop:sampled-laplace-exact}).
% Remark 2.19's check is answered by prop:sampled-laplace-exact (1)-(2), which the sentence cites.

\section{The $q$-Cin rays on refining lattices}
\label{sec:line-random-walks}

The $q$-Cin rays, the extreme rays of the cone $\mathcal{K}_q$ of
Proposition~\ref{prop:strata-cylinder-coordinates-q-cin}, and the knot-atom cascades are computed
by symmetric random walks, with step sizes shrinking by the ratio, each applied once. The kernel is
the limit of truncated cascades. For an integer ratio each truncation is carried by one lattice,
and as $q \downarrow 1$ the kernels converge to the Cin kernels.

% draft: Proposition 2.24 (section 2.7), lines 317-321 of draft/spatial-hemigroup-pyramids.md.
% CHANGED: the translations T_{+-v} of (2.14) written as the point masses delta_{+-v} (T_c is the
% dilation in this paper).
% CHANGED: "computed exactly there by N + 1 sparse convolutions" narrowed to the statement about
% supports (each factor has unbounded support, so the convolutions are not finite computations).
% CHANGED: the limit q -> 1 is stated for real q (an integer q cannot decrease to 1), with the
% bound of prop:strata-cylinder-coordinates-q-cin; the convergence of the partial products with
% the bound of the draft's proof.
% CHANGED: "every pyramid of section 2 with an atomic surface measure is a finite composition of
% such steps between its knots" narrowed to finitely many atoms (with infinitely many atoms the
% composition between two scales need not be finite), and these are members of the cone of
% Theorem 2.11, not families.
% CHANGED: the title names the q-Cin rays, the extreme rays of the cone K_q (blind review, flag 2).
% CHANGED: clause (5) names the member of the cone as the pair (Lambda_sigma, 0), with no Gaussian
% part (blind review 2, flag 16).
\begin{proposition}[the $q$-Cin rays: a cascade of random walks on refining lattices]
\label{prop:extreme-rays-random-walks}
\lean{Pyramids.extreme_rays_random_walks_1, Pyramids.extreme_rays_random_walks_2,
Pyramids.extreme_rays_random_walks_3, Pyramids.extreme_rays_random_walks_4,
Pyramids.extreme_rays_random_walks_5, Pyramids.extreme_rays_random_walks_5_finite}
\leanok
\uses{prop:strata-cylinder-coordinates-q-cin, thm:knot-atom-cascades}
Let $q > 1$, $\ell = \log q$ and $\tau > 0$. For $v, c > 0$ let
\begin{equation}\label{eq:line-e9}
  \exp\bigl[c\,\delta^2_v\bigr] := e^{-c}\sum_{n \ge 0}\frac{c^n}{n!}
  \Bigl(\tfrac12(\delta_v + \delta_{-v})\Bigr)^{*n}.
\end{equation}
\begin{enumerate}
\item $\exp[c\,\delta^2_v]$ is the law of a symmetric random walk with steps $\pm v$ and a
Poisson($c$) number of steps, a law carried by $v\ZZ$, with transform $e^{-c(1 - \cos v\omega)}$.
\item The law $K^{(q)}_\tau$ with transform $\exp[-\ell E^{(q)}_\tau(\omega)]$, the $q$-Cin kernel of
cutoff $\tau$ with per-octave normalization, is the weak limit as $N \to \infty$ of
\[
  K_N := \exp\bigl[\ell\,\delta^2_\tau\bigr] * \exp\bigl[\ell\,\delta^2_{\tau q^{-1}}\bigr] * \cdots *
  \exp\bigl[\ell\,\delta^2_{\tau q^{-N}}\bigr],
\]
and the exponents differ by at most $\ell\tau^2\omega^2q^{-2N}/(2(q^2 - 1))$.
\item If $q$ is an integer, $K_N$ is carried by the lattice $\tau q^{-N}\ZZ$, and its $i$-th factor
by the sublattice $\tau q^{-i}\ZZ$.
\item As $q \downarrow 1$ through real values, $K^{(q)}_\tau$ converges weakly to the law with
exponent $\Cin(\tau\omega)$, the Cin kernel of step $\tau$.
\item The kernel at scale $t$ of the knot-atom cascade (\ref{eq:line-e4}) is the convolution of the laws
$\exp[\delta^2_{r\theta q^n}]$ over the knots $\theta q^n \le t$. For a cylinder measure
$\sigma = \sum_{j=1}^Jc_j\delta_{(\theta_j, r_j)}$ with finitely many atoms, the step from $s$ to $t$,
$0 < s < t$, of the member $(\Lambda_\sigma, 0)$ of the cone of
Theorem~\ref{thm:knot-atom-cascades}, with $\mathcal{B} = 0$, is the convolution of the finitely
many laws $\exp[c_j\delta^2_{r_j\theta_jq^n}]$ over the pairs $(j, n)$
with $s < \theta_jq^n \le t$.
\end{enumerate}
% CHANGED (2026-11-14, a strengthening): unused hypotheses dropped from the Lean; the print keeps
% the setting of the node.
% CHANGED (fidelity review, 2026-11-15, R26/R50/R69): the status line lists the declarations' hypotheses beside the node's (LINKAGE rule 6).
\statusT\quad\emph{(Cites nothing of its own. Four declarations require less than the node, their
unused hypotheses dropped on 2026-11-14; the series and the lattices depend on $v$ and $\tau$ only
through $v\omega$, $\tau\omega$ and $v\ZZ$, $\tau q^{-N}\ZZ$. \textbf{Hypotheses} (the node's, each
with what its declarations require): $q > 1$: same, absent from \texttt{extreme\_rays\_random\_walks\_1},
which has no $q$, and in \texttt{\_4} the limit $q \downarrow 1$ through $q > 1$; $\tau > 0$: same
in \texttt{\_5}, absent from \texttt{\_2}, \texttt{\_3} and \texttt{\_4}; (1) $c > 0$: same;
$v > 0$: absent; (2) the law $K^{(q)}_\tau$: same, any probability law with its transform, whose
existence the declaration proves; (3) $q$ an integer: same, $q \ge 2$; (4) the Cin kernel and the
$K^{(q)}_\tau$: same, any probability laws with their transforms; (5) $(\theta, r) \in
\mathcal{C}$, $t > 0$, the kernel at scale $t$: same, any probability law with the transform of the
cascade; finitely many atoms with $c_j > 0$ and $(\theta_j, r_j) \in \mathcal{C}$: same;
$0 < s < t$: same.)}
\end{proposition}

% CHANGED (wave 3c, 2026-11-14): the proof follows the machine-checked route. In (1) the lattice is
% reached through convolutions of measures carried by a set closed under addition. In (2) the
% continuity of the limit transform at 0 is shown from the remainder bound at N = 0, and the law is
% obtained from the existence clause of Levy's continuity theorem. In (4) the convergence is drawn
% from the bound of prop:strata-cylinder-coordinates-q-cin(2), the Cin kernel is obtained as a
% limit along q_n = 1 + 1/(n+1), and the limit q -> 1 is taken along sequences. In (5) the knots
% below t are indexed by n <= floor(log_q(t/theta)), and the finite atomic case is computed atom by
% atom.
% CHANGED (the theory paper, 2026-11-16): the closing sentence on checking and cited facts shortened; no step changed.
\begin{proof}
\leanok
\uses{prop:strata-cylinder-coordinates-q-cin, thm:knot-atom-cascades}
Two facts are used throughout. A convolution of two measures carried by a Borel set $S$ closed
under addition is carried by $S$: if $x + y \notin S$ then $x \notin S$ or $y \notin S$, and both
of these sets are null for the product measure. And Lévy's continuity theorem in two parts: if the
transforms of a sequence of laws on $\RR$ converge pointwise to a function continuous at $0$, that
function is the transform of a law (the sequence is tight, by Prokhorov's theorem a subsequence
converges weakly, and the transform of the limit is the limit of the transforms); and if they
converge pointwise to the transform of a law, the sequence converges weakly to that law.

(1) The $n$-fold convolution power of $\tfrac12(\delta_v + \delta_{-v})$ is the law of $n$ steps of
the walk and has the transform $\cos^n(v\omega)$. Mixing over a Poisson($c$) number of steps gives
$e^{-c}\sum_nc^n\cos^n(v\omega)/n! = e^{-c(1 - \cos v\omega)}$, and the weights sum to $1$. The law
$\tfrac12(\delta_v + \delta_{-v})$ and $\delta_0$ are carried by $v\ZZ$, which is closed under
addition, so every convolution power is, and so is the mixture.

(2) By (1) the transform of a convolution is the product of the transforms, so $K_N$ has the
exponent $\ell\sum_{i=0}^N(1 - \cos(\omega\tau q^{-i}))$. Since
$0 \le 1 - \cos x \le x^2/2$, the remainder satisfies
\[
  0 \le \ell\sum_{i > N}\bigl(1 - \cos(\omega\tau q^{-i})\bigr)
  \le \tfrac\ell2\tau^2\omega^2\sum_{i>N}q^{-2i} = \frac{\ell\tau^2\omega^2q^{-2N}}{2(q^2 - 1)} .
\]
So the transforms of the $K_N$ converge pointwise to $\exp[-\ell E^{(q)}_\tau]$. At $N = 0$ the
bound gives $0 \le \ell E^{(q)}_\tau(\omega) \le \ell\tau^2\omega^2\bigl(\tfrac12 + \tfrac1{2(q^2 - 1)}\bigr)$,
so $\ell E^{(q)}_\tau(\omega) \to 0$ as $\omega \to 0$ and $\exp[-\ell E^{(q)}_\tau]$ is continuous
at $0$. By Lévy's continuity theorem a law $K^{(q)}_\tau$ with that transform exists, and the $K_N$
converge weakly to it.

(3) For an integer $q$ and $i \le N$, $\tau q^{-i}\ZZ = \tau q^{-N}(q^{N-i}\ZZ) \subseteq \tau q^{-N}\ZZ$.
By (1) the $i$-th factor is carried by $\tau q^{-i}\ZZ$, hence by $\tau q^{-N}\ZZ$, which is closed
under addition and contains $0$; so the convolution $K_N$ is carried by it.

(4) By Proposition~\ref{prop:strata-cylinder-coordinates-q-cin}(2),
$|\ell E^{(q)}_\tau(\omega) - \Cin(\tau\omega)| \le \frac{q\ell^2}{2(q - 1)}|\tau\omega|$, and since
$0 \le \ell \le q - 1$ the coefficient is at most $q(q - 1)/2$, which tends to $0$ as $q \downarrow 1$.
So $\ell E^{(q)}_\tau \to \Cin(\tau\,\cdot)$ pointwise. Along $q_n := 1 + 1/(n + 1)$ the laws
$K^{(q_n)}_\tau$ of (2) have transforms converging to $e^{-\Cin(\tau\omega)}$, which is continuous,
so by Lévy's continuity theorem the Cin kernel of step $\tau$ exists. The filter $q \downarrow 1$
has a countable base, so the weak convergence of $K^{(q)}_\tau$ as $q \downarrow 1$ is its
convergence along every sequence $q_n \downarrow 1$, which is Lévy's continuity theorem again.

(5) For $\theta > 0$, $t > 0$ and $n \in \ZZ$, $\theta q^n \le t$ if and only if
$n \le L := \lfloor\log_q(t/\theta)\rfloor$. By (1) with $c = 1$ the convolution of the laws
$\exp[\delta^2_{r\theta q^n}]$ over $-N \le n \le L$ has the exponent
$\sum_{n=-N}^{L}(1 - \cos(\omega r\theta q^n))$. The series (\ref{eq:line-e4}),
$\Psi^{(\theta,r)}_t(\omega) = \sum_{\theta q^n \le t}(1 - \cos(\omega r\theta q^n))$, has terms at most
$\tfrac12(\omega r\theta)^2q^{2n}$ and none beyond $n = L$, so it converges, and its partial sums over
$-N \le n \le N$ (equal, once $N \ge L$, to those over $-N \le n \le L$) tend to it. The transforms
converge to $e^{-\Psi^{(\theta,r)}_t}$, the transform of the kernel at $t$
(Theorem~\ref{thm:knot-atom-cascades}(4)), and by Lévy's continuity theorem the finite convolutions
converge weakly to that kernel.

For the finite atomic $\sigma = \sum_jc_j\delta_{(\theta_j,r_j)}$ one has
$\Lambda_\sigma = \sum_jc_j\Lambda^{(\theta_j,r_j)}$, and with $\mathcal{B} = 0$, (\ref{eq:line-e2}) integrates
$1 - \cos(\omega v)$ over the atoms of each orbit in $(0,t] \times (0,\infty)$: its value at $t$ is
$\sum_jc_j\Psi^{(\theta_j,r_j)}_t(\omega)$. Subtracting the value at $s$, the terms with
$\theta_jq^n \le s$ cancel, and the exponent of the step from $s$ to $t$ is
\[
  \sum_jc_j\sum_{s < \theta_jq^n \le t}\bigl(1 - \cos(\omega r_j\theta_jq^n)\bigr),
\]
a finite sum over $\lfloor\log_q(s/\theta_j)\rfloor < n \le \lfloor\log_q(t/\theta_j)\rfloor$. By (1)
it is the exponent of the convolution of the laws $\exp[c_j\delta^2_{r_j\theta_jq^n}]$ over these
pairs, which is a probability law.

Every step is machine-checked, and the proof cites no fact of its own.
\end{proof}

The Cin kernels are thus limits of cascades as simple as the \Matern{} recursion and of another
kind. Each stage is a symmetric random walk with a Poisson($\ell$) number of steps of one size, and
the sizes shrink by $q$ from one stage to the next, so the cascade runs from coarse steps to fine
ones. The extreme rays of the continuum cone are the limits $q \downarrow 1$ of these cascades, and
for an integer $q$ each truncated cascade is carried by one lattice. This is the second link between
the discrete stratum and the continuous one. The first is the circle average of
Proposition~\ref{prop:strata-cylinder-coordinates}(2).
